% Fonctions : Dérivée et primitives — Mathématiques Tle Tech — Cours signé OnBuch+
% Programme officiel MINESEC. Compiler avec : tectonic N02-fonctions-derivee-primitives.tex
\newcommand{\DOCMATIERE}{Mathématiques}
\newcommand{\DOCNIVEAU}{Tle Tech}
\newcommand{\DOCMODULE}{Mathématiques – classe de Terminale technique}
\newcommand{\DOCLECON}{N02}
\newcommand{\DOCTITRE}{Fonctions : Dérivée et primitives}
% OnBuch+ — préambule « cours signés » ENRICHI (compilé avec tectonic / XeLaTeX)
% Système visuel « Pop » d'OnBuch+ : fond crème (jamais blanc), bordures encre
% épaisses, ombres « dures » décalées, titres Archivo Black en capitales.
% Version MATHÉMATIQUES : graphes (pgfplots/tikz), boîtes pédagogiques
% (définition, propriété, méthode, attention, le savais-tu, exercice résolu,
% exercice, TP, bilan), page de garde et signature OnBuch+ ancrée en bas.
%
% Macros attendues dans main.tex AVANT \input{preamble} :
%   \DOCMATIERE  \DOCNIVEAU  \DOCMODULE  \DOCLECON (numéro)  \DOCTITRE
\documentclass[11pt]{article}

\usepackage{fontspec}
\usepackage[french]{babel}
\usepackage[a4paper,margin=2cm,top=3.4cm,bottom=2.8cm,headheight=1.4cm,footskip=1.3cm]{geometry}
\usepackage{amsmath,amssymb,amsfonts}
\usepackage{mathtools} % \xmapsto, \coloneqq... souvent utilisés par les modèles
\usepackage{cancel} % simplifications barrées
% \abs{z} (module, valeur absolue) et \norm{u} (norme), utilisés par les modèles
\providecommand{\abs}{}\let\abs\relax\DeclarePairedDelimiter{\abs}{\lvert}{\rvert}
\providecommand{\norm}{}\let\norm\relax\DeclarePairedDelimiter{\norm}{\lVert}{\rVert}
\usepackage[table]{xcolor} % [table] : fournit aussi \rowcolors (alternance de lignes)
\usepackage{pagecolor}
\usepackage{tikz}
\usetikzlibrary{arrows.meta,positioning,calc,shapes.geometric,shapes.misc,shapes.arrows,decorations.pathmorphing,decorations.pathreplacing,decorations.markings,patterns,shadows,mindmap,trees,fit,backgrounds,matrix,intersections,angles,quotes}
\usepackage{pgfplots}
\usepackage{pgf-pie} % diagrammes circulaires \pie (statistiques)
\pgfplotsset{compat=1.18}
\usepgfplotslibrary{fillbetween,groupplots} % aires entre courbes (calcul intégral)
\usepackage{enumitem}
\usepackage{fancyhdr}
% [most] charge toutes les bibliothèques tcolorbox (listings, minted, raster,
% documentation...) et gonfle inutilement la mémoire TeX sur un document long
% (~300 boîtes) : on ne charge que ce qui est réellement utilisé.
\usepackage[skins,breakable]{tcolorbox}
\usepackage{graphicx}
\usepackage{colortbl}
\usepackage{booktabs}
\usepackage{tabularx}
\usepackage{diagbox} % case d'en-tête diagonale des tableaux à double entrée
% Colonnes extensibles centrée / gauche / droite (C, L, R), souvent utilisées par les modèles
\newcolumntype{C}{>{\centering\arraybackslash}X}
\newcolumntype{L}{>{\raggedright\arraybackslash}X}
\newcolumntype{R}{>{\raggedleft\arraybackslash}X}
\usepackage{array}
\usepackage{multicol}
\usepackage{siunitx}
% parse-numbers=false : accepte aussi \SI{2,5 \times 10^{-3}}{...}
\sisetup{locale=FR,output-decimal-marker={,},group-minimum-digits=4,parse-numbers=false}
\DeclareSIUnit\ohm{\ensuremath{\symup{\Omega}}} % U+2126 absent de la police
\usepackage{qrcode}
% \degree : commande fréquemment utilisée par les modèles pour un angle ou une
% température en dehors de \SI{}{} (non fournie par les paquets chargés).
\providecommand{\degree}{^{\circ}}
% \qed (fin de démonstration), utilisé par les modèles sans amsthm
\providecommand{\qed}{\hfill$\square$}
% \barre : double barre des valeurs interdites dans un tableau de variations (tkz-tab)
\providecommand{\barre}{\ensuremath{\|}}
% environnement proof (amsthm non chargé) utilisé par les modèles
\ifdefined\proof\else\newenvironment{proof}[1][Démonstration]{\par\noindent\textit{#1.}\ }{\qed\par}\fi
\usepackage{lastpage}
\usepackage{hyperref}
\hypersetup{hidelinks}

% ============================================================
% Palette « Pop » OnBuch+ — mêmes hex que la classe P (plus_ui.dart)
% ============================================================
\definecolor{poporange}{HTML}{F59321}
\definecolor{popdark}{HTML}{E07A0C}
\definecolor{popcream}{HTML}{FFF6E8}
\definecolor{popink}{HTML}{1C1714}
\definecolor{poppurple}{HTML}{7A5AE0}
\definecolor{popgreen}{HTML}{2E9E6A}
\definecolor{poppink}{HTML}{E0457B}
\definecolor{popblue}{HTML}{2F7DE1}
\definecolor{popgold}{HTML}{C98A12}
\definecolor{popmuted}{HTML}{8A7F76}
\definecolor{popline}{HTML}{EADFD2}
\definecolor{popgray}{HTML}{8A7F76}
\colorlet{popgrey}{popgray}
% Alias tolérants (le modèle nomme parfois les couleurs différemment)
\colorlet{popwhite}{white}
\colorlet{popblack}{popink}
\colorlet{popred}{poppink}
\colorlet{popyellow}{popgold}
% Teintes claires pour fonds de tableaux / zones
\colorlet{popblueL}{popblue!10!white}
\colorlet{poporangeL}{poporange!14!white}
\colorlet{popgreenL}{popgreen!12!white}
\colorlet{poppurpleL}{poppurple!12!white}
\colorlet{poppinkL}{poppink!12!white}
\colorlet{popgoldL}{popgold!14!white}

\pagecolor{popcream}

% ============================================================
% Typographie
% ============================================================
\defaultfontfeatures{Ligatures=TeX}
\setmainfont{PlusJakartaSans-Medium.ttf}[Path=../../../fonts/,BoldFont=PlusJakartaSans-Bold.ttf]
% Maths en sans-serif (Fira Math) avec chiffres et lettres droites de Plus
% Jakarta, pour que les formules se fondent dans le texte.
\usepackage[math-style=ISO,bold-style=ISO]{unicode-math}
\setmathfont{FiraMath-Regular.otf}[Path=../../../fonts/]
\setmathfont{PlusJakartaSans-Medium.ttf}[Path=../../../fonts/,range={up/num,up/{Latin,latin},\mathup}]
\usepackage{icomma} % virgule décimale sans espace en mode math
% Caractères mathématiques Unicode absents des polices de texte (Plus Jakarta,
% Archivo Black) : rendus par leur équivalent mathématique, en texte comme en
% formule — sinon ils s'affichent en carré vide (ex. « Arithmétique dans ℤ »).
\usepackage{newunicodechar}
\newunicodechar{ℝ}{\ensuremath{\mathbb{R}}}
\newunicodechar{ℤ}{\ensuremath{\mathbb{Z}}}
\newunicodechar{ℂ}{\ensuremath{\mathbb{C}}}
\newunicodechar{ℕ}{\ensuremath{\mathbb{N}}}
\newunicodechar{ℚ}{\ensuremath{\mathbb{Q}}}
\newunicodechar{𝔻}{\ensuremath{\mathbb{D}}}
\newunicodechar{α}{\ensuremath{\alpha}}
\newunicodechar{β}{\ensuremath{\beta}}
\newunicodechar{γ}{\ensuremath{\gamma}}
\newunicodechar{δ}{\ensuremath{\delta}}
\newunicodechar{ε}{\ensuremath{\varepsilon}}
\newunicodechar{θ}{\ensuremath{\theta}}
\newunicodechar{λ}{\ensuremath{\lambda}}
\newunicodechar{μ}{\ensuremath{\mu}}
\newunicodechar{σ}{\ensuremath{\sigma}}
\newunicodechar{τ}{\ensuremath{\tau}}
\newunicodechar{φ}{\ensuremath{\varphi}}
\newunicodechar{ψ}{\ensuremath{\psi}}
\newunicodechar{ω}{\ensuremath{\omega}}
\newunicodechar{Σ}{\ensuremath{\Sigma}}
% majuscules produites par \MakeUppercase dans les titres (α-aminés -> Α)
\newunicodechar{Α}{\ensuremath{\alpha}}
\newunicodechar{Β}{\ensuremath{\beta}}
\newunicodechar{Γ}{\ensuremath{\Gamma}}
\newunicodechar{Δ}{\ensuremath{\Delta}}
\newunicodechar{Ε}{\ensuremath{\varepsilon}}
\newunicodechar{Θ}{\ensuremath{\Theta}}
\newunicodechar{Λ}{\ensuremath{\Lambda}}
\newunicodechar{Μ}{\ensuremath{\mu}}
\newunicodechar{Τ}{\ensuremath{\tau}}
\newunicodechar{Φ}{\ensuremath{\Phi}}
\newunicodechar{Ψ}{\ensuremath{\Psi}}
\newunicodechar{Ω}{\ensuremath{\Omega}}
\newunicodechar{↦}{\ensuremath{\mapsto}}
\newunicodechar{∈}{\ensuremath{\in}}
\newunicodechar{∉}{\ensuremath{\notin}}
\newunicodechar{∧}{\ensuremath{\wedge}}
\newunicodechar{⇔}{\ensuremath{\Leftrightarrow}}
\newunicodechar{⟺}{\ensuremath{\Longleftrightarrow}}
\newunicodechar{⇒}{\ensuremath{\Rightarrow}}
\newunicodechar{⟹}{\ensuremath{\Longrightarrow}}
\newunicodechar{∘}{\ensuremath{\circ}}
\newunicodechar{∪}{\ensuremath{\cup}}
\newunicodechar{∩}{\ensuremath{\cap}}
\newunicodechar{≡}{\ensuremath{\equiv}}
\newunicodechar{⊂}{\ensuremath{\subset}}
\newunicodechar{⊥}{\ensuremath{\perp}}
\newunicodechar{∃}{\ensuremath{\exists}}
\newunicodechar{∀}{\ensuremath{\forall}}
\newunicodechar{∛}{\ensuremath{\sqrt[3]{\,}}}
\newunicodechar{∜}{\ensuremath{\sqrt[4]{\,}}}
\newunicodechar{⃗}{\ensuremath{{}^{\to}}}
\newunicodechar{ⁿ}{\ifmmode^{n}\else\textsuperscript{\ensuremath{n}}\fi}
\newunicodechar{ˣ}{\ifmmode^{x}\else\textsuperscript{\ensuremath{x}}\fi}
\newunicodechar{ᵇ}{\ifmmode^{b}\else\textsuperscript{\ensuremath{b}}\fi}
\newunicodechar{ʳ}{\ifmmode^{r}\else\textsuperscript{\ensuremath{r}}\fi}
\newunicodechar{ᵘ}{\ifmmode^{u}\else\textsuperscript{\ensuremath{u}}\fi}
\newunicodechar{ʸ}{\ifmmode^{y}\else\textsuperscript{\ensuremath{y}}\fi}
\newunicodechar{ᵃ}{\ifmmode^{a}\else\textsuperscript{\ensuremath{a}}\fi}
\newunicodechar{ᵉ}{\ifmmode^{e}\else\textsuperscript{\ensuremath{e}}\fi}
\newunicodechar{ᵏ}{\ifmmode^{k}\else\textsuperscript{\ensuremath{k}}\fi}
\newunicodechar{ᵖ}{\ifmmode^{p}\else\textsuperscript{\ensuremath{p}}\fi}
\newunicodechar{ᴺ}{\ifmmode^{N}\else\textsuperscript{\ensuremath{N}}\fi}
\newunicodechar{ᶜ}{\ifmmode^{c}\else\textsuperscript{\ensuremath{c}}\fi}
\newunicodechar{ʰ}{\ifmmode^{h}\else\textsuperscript{\ensuremath{h}}\fi}
\newunicodechar{ᵠ}{\ifmmode^{q}\else\textsuperscript{\ensuremath{q}}\fi}
\newunicodechar{ₙ}{\ifmmode_{n}\else\textsubscript{\ensuremath{n}}\fi}
\newunicodechar{ₐ}{\ifmmode_{a}\else\textsubscript{\ensuremath{a}}\fi}
\newunicodechar{ᵢ}{\ifmmode_{i}\else\textsubscript{\ensuremath{i}}\fi}
\newunicodechar{ᵣ}{\ifmmode_{r}\else\textsubscript{\ensuremath{r}}\fi}
\newunicodechar{ₖ}{\ifmmode_{k}\else\textsubscript{\ensuremath{k}}\fi}
\newunicodechar{ᵦ}{\ifmmode_{\beta}\else\textsubscript{\ensuremath{\beta}}\fi}
\newunicodechar{ₓ}{\ifmmode_{x}\else\textsubscript{\ensuremath{x}}\fi}
\newfontfamily\popdisplay{ArchivoBlack-Regular.ttf}[Path=../../../fonts/]
\newfontfamily\poplogo{SpaceGrotesk-Bold.ttf}[Path=../../../fonts/]
\newfontfamily\popsemi{PlusJakartaSans-SemiBold.ttf}[Path=../../../fonts/]
\newfontfamily\popxbold{PlusJakartaSans-ExtraBold.ttf}[Path=../../../fonts/]
\newfontfamily\popxboldls{PlusJakartaSans-ExtraBold.ttf}[Path=../../../fonts/,LetterSpace=3.0]

\setlist[enumerate]{leftmargin=1.7em,itemsep=5pt,topsep=4pt}
\setlist[itemize]{leftmargin=1.7em,itemsep=4pt,topsep=4pt,label={\color{poporange}\textbullet}}
\setlength{\parindent}{0pt}
\setlength{\parskip}{5pt}
\renewcommand{\arraystretch}{1.25}

% pgfplots : style Pop par défaut
\pgfplotsset{
  every axis/.append style={axis line style={popink,line width=1pt},tick style={popink},
    grid style={popline,line width=0.5pt},label style={font=\small},tick label style={font=\footnotesize}},
  popcurve/.style={poporange,line width=1.8pt,no markers,smooth},
}
% Même style disponible dans un \draw TikZ simple (hors pgfplots)
\tikzset{popcurve/.style={poporange,line width=1.8pt}}

% ============================================================
% Pill / squiggle
% ============================================================
\newcommand{\poppill}[3]{%
  \tikz[baseline=-0.55ex]\node[draw=popink,line width=1.2pt,rounded corners=2.6pt,%
    fill=#2,inner xsep=6.5pt,inner ysep=3pt]{%
    \popxboldls\fontsize{7}{7}\selectfont\color{#3}\MakeUppercase{#1}};%
}
\newcommand{\popsquiggle}[1]{%
  \tikz[baseline=-0.4ex]{
    \draw[#1,line width=2.6pt,line cap=round]
      plot [smooth] coordinates {(0,0.09) (0.34,0.01) (0.68,0.21) (1.02,0.01) (1.36,0.10)};
  }%
}
% Mot-clé surligné (marqueur orange sous le texte)
\newcommand{\cle}[1]{{\popxbold\color{popdark}#1}}

% ============================================================
% En-tête / pied
% ============================================================
\pagestyle{fancy}
\fancyhf{}
\renewcommand{\headrulewidth}{0pt}
\renewcommand{\footrulewidth}{0pt}
\lhead{\raisebox{-0.15ex}{\poplogo\fontsize{13}{13}\selectfont\color{popink}OnBuch\color{poporange}+}}
\rhead{\poppill{\DOCMATIERE\ \textbullet\ \DOCNIVEAU\ \textbullet\ Leçon \DOCLECON}{white}{popink}}
\lfoot{\popxbold\fontsize{7.6}{7.6}\selectfont\color{popmuted}\MakeUppercase{Cours signé OnBuch+}\ \textbullet\ {\popsemi\DOCTITRE}}
\rfoot{\tikz[baseline=-0.6ex]\node[rounded corners=8.5pt,fill=popink,minimum height=17pt,minimum width=17pt,inner xsep=5pt,inner sep=0]{\popxbold\fontsize{8}{8}\selectfont\color{popcream}\thepage\,/\,\pageref*{LastPage}};}

% ============================================================
% Titres
% ============================================================
\newcommand{\chaptitre}[2]{%
  \begin{tcolorbox}[enhanced,colback=poporange,colframe=popink,boxrule=2.6pt,arc=11pt,
      drop shadow=popink,left=15pt,right=15pt,top=12pt,bottom=11pt,before skip=3pt,after skip=18pt]
    \poppill{#2}{popink}{popcream}\par\vspace{9pt}
    {\raggedright\hyphenpenalty=10000\exhyphenpenalty=10000\popdisplay\fontsize{22}{24}\selectfont\color{popcream}\MakeUppercase{#1}\par}
  \end{tcolorbox}
}
% Section numérotée : pastille orange avec le numéro + titre Archivo
\newcounter{popsec}
\newcommand{\coursec}[1]{%
  \stepcounter{popsec}\setcounter{popsub}{0}%
  \par\vspace{16pt}\noindent
  \begin{minipage}{\linewidth}
  \tikz[baseline=-0.7ex]\node[draw=popink,line width=1.6pt,fill=poporange,rounded corners=4pt,minimum size=22pt,inner sep=0]{\popdisplay\fontsize{12}{12}\selectfont\color{popcream}\thepopsec};%
  \hspace{8pt}{\popdisplay\fontsize{15}{17}\selectfont\color{popink}\MakeUppercase{#1}}\par
  \vspace{4pt}\noindent{\color{poporange}\rule{36pt}{3.6pt}}
  \end{minipage}\par\nopagebreak\vspace{8pt}
}
\newcounter{popsub}
\newcommand{\courssub}[1]{%
  \stepcounter{popsub}%
  \par\vspace{9pt}\noindent{\popxbold\fontsize{11.5}{13}\selectfont\color{popink}{\color{poporange}\thepopsec.\thepopsub}\hspace{6pt}#1}\par\nopagebreak\vspace{4pt}
}

% ============================================================
% Boîtes pédagogiques — même recette : fond blanc, bordure épaisse colorée,
% ombre dure encre, badge plein. Titre optionnel [#1].
% ============================================================
\tcbset{popbase/.style={breakable,enhanced,colback=white,boxrule=2.3pt,arc=9pt,
  shadow={3pt}{-3pt}{0pt}{popink},left=14pt,right=14pt,top=11pt,bottom=11pt,before skip=11pt,after skip=15pt,
  fontupper=\popsemi\fontsize{10.6}{14.5}\selectfont\color{popink},
  fontlower=\popsemi\fontsize{10.6}{14.5}\selectfont\color{popink}}}
% Coche dessinée (le glyphe ✓ n'existe pas dans les polices du document)
\newcommand{\popcheck}[1]{\tikz[baseline=-0.5ex]\draw[#1,line width=1.6pt,line cap=round,line join=round](-0.13,0.0)--(-0.04,-0.09)--(0.13,0.1);}
\newcommand{\popboxhead}[3]{\poppill{#1}{#2}{white}\ifx\relax#3\relax\else\hspace{6pt}{\popxbold\fontsize{10.6}{12}\selectfont\color{popink}#3}\fi\par\vspace{7pt}}

\newtcolorbox{aretenir}[1][]{popbase,colframe=popgreen,before upper={\popboxhead{À retenir}{popgreen}{#1}}}
\newtcolorbox{exemplebox}[1][]{popbase,colframe=poporange,before upper={\popboxhead{Exemple}{poporange}{#1}}}
\newtcolorbox{definition}[1][]{popbase,colframe=popblue,before upper={\popboxhead{Définition}{popblue}{#1}}}
\newtcolorbox{propriete}[1][]{popbase,colframe=poppurple,before upper={\popboxhead{Propriété}{poppurple}{#1}}}
\newtcolorbox{methode}[1][]{popbase,colframe=popgold,before upper={\popboxhead{Méthode}{popgold}{#1}}}
\newtcolorbox{attention}[1][]{popbase,colframe=poppink,before upper={\popboxhead{Attention}{poppink}{#1}}}
\newtcolorbox{experience}[1][]{popbase,colframe=popblue,colback=popblueL,before upper={\popboxhead{Expérience}{popblue}{#1}}}
% « Le savais-tu ? » : carte violette pleine (contexte, culture, Cameroun)
\newtcolorbox{savaistu}[1][]{popbase,colback=poppurple,colframe=popink,
  fontupper=\popsemi\fontsize{10.4}{14.2}\selectfont\color{white},
  before upper={\poppill{Le savais-tu\,?}{white}{poppurple}\ifx\relax#1\relax\else\hspace{6pt}{\popxbold\color{white}#1}\fi\par\vspace{7pt}}}
% Exercice résolu : énoncé en haut, \tcblower puis la solution sur fond crème
\newcounter{popexo}\newcounter{popexr}
\newtcolorbox{exoresolu}[1][]{popbase,colframe=poporange,colbacklower=popcream,
  segmentation style={popink,line width=1.2pt,dashed},
  before upper={\stepcounter{popexr}\popboxhead{Exercice résolu \thepopexr}{poporange}{#1}},
  before lower={\poppill{Solution}{popink}{popcream}\par\vspace{6pt}}}
% Exercice (non résolu en place) — la correction est dans la partie Corrigés
\newtcolorbox{exercice}[1][]{popbase,colframe=popink,
  before upper={\stepcounter{popexo}\popboxhead{Exercice \thepopexo}{popink}{#1}}}
% Corrigé
\newtcolorbox{corrige}[1][]{popbase,colframe=popgreen,colback=popgreenL,
  before upper={\popboxhead{Corrigé}{popgreen}{#1}}}
% Objectifs / prérequis / bilan
\newtcolorbox{objectifs}{popbase,colframe=popink,colback=poporangeL,
  before upper={\popboxhead{Objectifs de la leçon}{popink}{}}}
\newtcolorbox{prerequis}{popbase,colframe=popink,colback=popblueL,
  before upper={\popboxhead{Prérequis}{popblue}{}}}
\newtcolorbox{bilan}{popbase,colframe=popink,colback=white,boxrule=2.8pt,
  before upper={\popboxhead{Fiche bilan}{poporange}{L'essentiel à connaître pour l'examen}}}
% Figure encadrée avec légende
\newtcolorbox{popfigure}[1][]{enhanced,breakable=false,colback=white,colframe=popline,boxrule=1.4pt,arc=8pt,
  left=10pt,right=10pt,top=10pt,bottom=8pt,before skip=10pt,after skip=12pt,halign=center,
  lower separated=false,halign lower=center,
  fontlower=\popsemi\fontsize{9}{11.5}\selectfont\color{popmuted},
  #1}
\newcommand{\legende}[1]{\par\vspace{6pt}{\popsemi\fontsize{9}{11.5}\selectfont\color{popmuted}#1}}

% ============================================================
% Signature OnBuch+
% ============================================================
\newcommand{\poplogoinline}[1]{{\poplogo\fontsize{#1}{#1}\selectfont\color{popink}OnBuch\color{poporange}+}}
\newcommand{\signatureonbuch}{%
  \begin{tcolorbox}[enhanced,colback=popink,colframe=popink,boxrule=2.2pt,arc=10pt,
      drop shadow=poporange,left=14pt,right=14pt,top=10pt,bottom=10pt,before skip=0pt,after skip=0pt]
    \tikz[baseline=-0.7ex]\node[circle,fill=popgreen,draw=popcream,line width=1.3pt,minimum size=22pt,inner sep=0]{\popcheck{white}};%
    \hspace{9pt}\begin{minipage}[c]{0.62\linewidth}
      {\popxbold\fontsize{12}{13}\selectfont\color{popcream}Signé\ \textbullet\ {\poplogo OnBuch\color{poporange}+}}\par\vspace{2pt}
      {\raggedright\popsemi\fontsize{8}{10.5}\selectfont\color{popline}\MakeUppercase{Équipe pédagogique OnBuch+}\\\MakeUppercase{Conforme au programme officiel MINESEC}\par}
    \end{minipage}\hfill
    {\poplogo\fontsize{22}{22}\selectfont\color{popcream}OnBuch\color{poporange}+}
  \end{tcolorbox}%
}

% ============================================================
% Page de garde
% #1 titre · #2 matière · #3 niveau · #4 module · #5 n° leçon · #6 durée ·
% #7 description / objectifs (texte ou itemize)
% ============================================================
\newcommand{\pagegarde}[7]{%
  \thispagestyle{empty}
  \newgeometry{a4paper,margin=1.7cm,top=1.6cm,bottom=1.4cm}%
  \begin{tikzpicture}[remember picture,overlay]
    \fill[poporange!18] ([xshift=-3.2cm,yshift=-3.6cm]current page.north east) circle (4.6cm);
    \fill[poppurple!14] ([xshift=2.4cm,yshift=7.6cm]current page.south west) circle (3.8cm);
  \end{tikzpicture}%
  {\poplogo\fontsize{30}{30}\selectfont\color{popink}OnBuch\color{poporange}+}\hfill
  \raisebox{0.6ex}{\poppill{Cours signé \textbullet\ Premium}{popink}{popcream}}\par
  \vspace{1pt}\popsquiggle{poppurple}\par
  \vspace{8pt}
  \begin{tcolorbox}[enhanced,colback=poporange,colframe=popink,boxrule=2.8pt,arc=13pt,
      drop shadow=popink,left=18pt,right=18pt,top=12pt,bottom=13pt,after skip=10pt]
    \poppill{#2 \textbullet\ #3}{popink}{popcream}\hspace{5pt}\poppill{#4}{white}{popink}\par\vspace{8pt}
    {\popdisplay\fontsize{14}{14}\selectfont\color{popink}LEÇON #5}\par\vspace{4pt}
    {\raggedright\hyphenpenalty=10000\exhyphenpenalty=10000\popdisplay\fontsize{25}{27}\selectfont\color{popcream}\MakeUppercase{#1}\par}
    \vspace{8pt}
    {\popxbold\fontsize{9.5}{9.5}\selectfont\color{popink}\MakeUppercase{Durée conseillée : #6}}
  \end{tcolorbox}
  \begin{tcolorbox}[enhanced,colback=white,colframe=popink,boxrule=2.2pt,arc=10pt,
      drop shadow=popink,left=16pt,right=16pt,top=10pt,bottom=10pt,after skip=8pt]
    \poppill{Dans cette leçon}{poppurple}{white}\par\vspace{5pt}
    \popsemi\fontsize{9.3}{12.6}\selectfont\color{popink}#7
  \end{tcolorbox}
  \noindent\begin{minipage}[t]{0.487\linewidth}
    \begin{tcolorbox}[enhanced,colback=popgreen,colframe=popink,boxrule=2.2pt,arc=10pt,
        drop shadow=popink,left=11pt,right=11pt,top=10pt,bottom=10pt,height=3.1cm,valign=center]
      \tcbox[colback=white,colframe=popink,boxrule=1.3pt,arc=5pt,boxsep=0pt,nobeforeafter,
        left=3pt,right=3pt,top=3pt,bottom=3pt,valign=center]{\qrcode[height=1.9cm]{https://play.google.com/store/apps/details?id=com.luvvix.onbuch&hl=fr}}%
      \hspace{8pt}\begin{minipage}[c]{0.5\linewidth}
        {\popdisplay\fontsize{11}{12.5}\selectfont\color{white}\MakeUppercase{Télécharge OnBuch}}\par\vspace{4pt}
        {\raggedright\popsemi\fontsize{8.4}{11}\selectfont\color{white}Gratuit \textbullet\ Google Play\\Tuteur IA, annales, concours, résultats\par}
      \end{minipage}
    \end{tcolorbox}
  \end{minipage}\hfill
  \begin{minipage}[t]{0.487\linewidth}
    \begin{tcolorbox}[enhanced,colback=poppurple,colframe=popink,boxrule=2.2pt,arc=10pt,
        drop shadow=popink,left=11pt,right=11pt,top=10pt,bottom=10pt,height=3.1cm,valign=center]
      \tikz[baseline=-0.7ex]\node[circle,fill=white,draw=popink,line width=1.3pt,minimum size=1.6cm,inner sep=0]{\popdisplay\fontsize{24}{24}\selectfont\color{poppurple}+};%
      \hspace{8pt}\begin{minipage}[c]{0.6\linewidth}
        {\popdisplay\fontsize{11}{12.5}\selectfont\color{white}\MakeUppercase{Passe à OnBuch+}}\par\vspace{4pt}
        {\raggedright\popsemi\fontsize{8.4}{11}\selectfont\color{white}Tous les cours signés, Tuteur IA illimité, fiches PDF — onglet OnBuch+ de l'app\par}
      \end{minipage}
    \end{tcolorbox}
  \end{minipage}
  \vspace{8pt}
  \popcontenu
  \vfill
  \signatureonbuch
  \restoregeometry
  \newpage
}

% Rangée « Ce cours contient » (page de garde)
\newcommand{\poptile}[3]{% #1 couleur #2 gros texte #3 légende
  \begin{tcolorbox}[enhanced,colback=#1,colframe=popink,boxrule=2pt,arc=9pt,shadow={2.5pt}{-2.5pt}{0pt}{popink},
      left=6pt,right=6pt,top=9pt,bottom=9pt,halign=center,valign=center,height=1.75cm,width=\linewidth,nobeforeafter]
    {\popdisplay\fontsize{12.5}{13}\selectfont\color{white}\MakeUppercase{#2}}\par\vspace{3pt}
    {\centering\popsemi\fontsize{7.8}{9.5}\selectfont\color{white}#3\par}
  \end{tcolorbox}}
\newcommand{\popcontenu}{%
  \par\noindent{\popxboldls\fontsize{7.5}{7.5}\selectfont\color{popmuted}CE COURS SIGNÉ CONTIENT}\par\vspace{7pt}
  \noindent\begin{minipage}[t]{0.235\linewidth}\poptile{popblue}{Cours}{complet, illustré, pas à pas}\end{minipage}\hfill
  \begin{minipage}[t]{0.235\linewidth}\poptile{poporange}{Méthodes}{et exercices résolus}\end{minipage}\hfill
  \begin{minipage}[t]{0.235\linewidth}\poptile{poppink}{Exercices}{type examen + corrigés}\end{minipage}\hfill
  \begin{minipage}[t]{0.235\linewidth}\poptile{popgreen}{Fiche bilan}{l'essentiel à retenir}\end{minipage}\par}

% Sommaire
\newtcolorbox{sommaire}{enhanced,colback=white,colframe=popink,boxrule=2.2pt,arc=10pt,
  drop shadow=popink,left=18pt,right=18pt,top=13pt,bottom=13pt,before skip=6pt,after skip=20pt,
  before upper={\poppill{Sommaire}{poppurple}{white}\par\vspace{9pt}\popsemi\fontsize{10.6}{14.5}\selectfont\color{popink}}}

% Signature de fin de cours — poussée tout en bas de la dernière page
\newcommand{\findecours}{%
  \par\vfill\nopagebreak
  \begin{center}\popsquiggle{poporange}\end{center}\vspace{4pt}
  \signatureonbuch
}
\AtBeginDocument{\def\llbracket{\mathopen{[\mkern-3.5mu[}}\def\rrbracket{\mathclose{]\mkern-3.5mu]}}} % intervalles d'entiers (glyphe ⟦ absent de la police)
% Glyphes absents de Fira Math : redéfinis à partir de glyphes disponibles
\AtBeginDocument{%
  \def\setminus{\mathbin{\backslash}}\let\smallsetminus\setminus
  \def\vdots{\mathord{\vbox{\baselineskip3.2pt\lineskiplimit0pt\kern4pt\hbox{.}\hbox{.}\hbox{.}}}}%
  \def\ddots{\mathinner{\mkern1mu\raise6.5pt\hbox{.}\mkern2mu\raise3.6pt\hbox{.}\mkern2mu\raise0.7pt\hbox{.}\mkern1mu}}%
  \def\triangle{\mathord{\scalebox{0.9}{$\symup{\Delta}$}}}%
  \def\nabla{\mathord{\raisebox{\depth}{\scalebox{1}[-1]{$\symup{\Delta}$}}}}%
  \def\checkmark{\popcheck{popdark}}%
  \def\star{\mathbin{\ast}}\def\bigstar{\mathord{\ast}}%
  \def\diamond{\mathbin{\scalebox{0.6}{$\Diamond$}}}%
  \def\bigcup{\mathop{\mathchoice{\raisebox{-0.3em}{\scalebox{1.7}{$\cup$}}}{\scalebox{1.25}{$\cup$}}{\cup}{\cup}}\displaylimits}%
  \def\bigcap{\mathop{\mathchoice{\raisebox{-0.3em}{\scalebox{1.7}{$\cap$}}}{\scalebox{1.25}{$\cap$}}{\cap}{\cap}}\displaylimits}%
  \def\lesseqqgtr{\mathrel{\lessgtr}}\def\bigoplus{\mathop{\oplus}\displaylimits}%
}
\providecommand{\ssi}{si et seulement si} % abréviation fréquente
\AtBeginDocument{\providecommand{\wideparen}{\overparen}} % arc de cercle

%% FIGFIX-BEGIN v4 — correctifs globaux des figures (docs/onbuch-generation/tools/figfix)
\usepackage{adjustbox}\usepackage{etoolbox}
% toute figure TikZ/pgfplots plus large que la ligne est réduite à la largeur disponible
\BeforeBeginEnvironment{tikzpicture}{\begin{adjustbox}{max width=\linewidth}}
\AfterEndEnvironment{tikzpicture}{\end{adjustbox}}
% légendes pgfplots lisibles par-dessus les courbes
\pgfplotsset{every axis/.append style={legend style={fill=white,fill opacity=0.92,text opacity=1,draw=black!15,inner sep=3pt}}}
% étiquettes d'arêtes (syntaxe edge["..."]) avec fond blanc
\tikzset{every edge quotes/.append style={fill=white,inner sep=1.2pt}}
%% FIGFIX-END
\begin{document}

\pagegarde{Fonctions : Dérivée et primitives}{Mathématiques}{Tle Tech}{Mathématiques – classe de Terminale technique}{N02}{2 séances (fiche de progression harmonisée)}{Cette leçon aborde les dérivées usuelles et les primitives des fonctions polynômes, en mettant l'accent sur les règles de dérivation, le calcul de primitives et leur application à des problèmes concrets de la vie courante au Cameroun.
\vspace{4pt}
\begin{multicols}{2}\small
\begin{itemize}
\item À la fin de cette leçon, je sais définir la dérivée d'une fonction en un point et l'interpréter géométriquement comme le coefficient directeur de la tangente.
\item Je sais appliquer les règles de dérivation (linéarité, produit, quotient, chaîne) pour dériver tout polynôme.
\item Je sais dresser et utiliser le tableau des dérivées usuelles pour les puissances entières et rationnelles.
\item Je sais définir une primitive d'une fonction et expliquer le rôle de la constante d'intégration.
\item Je sais déterminer la primitive d'une fonction polynôme terme à terme et vérifier le résultat par dérivation.
\item Je sais modéliser une situation concrète (coût, revenu, bénéfice) par une fonction polynôme et utiliser la dérivée pour optimiser.
\end{itemize}\end{multicols}}

\begin{sommaire}
\begin{multicols}{2}
\begin{enumerate}
\item Rappels sur la dérivée et interprétation géométrique
\item Règles de dérivation usuelles
\item Tableau des dérivées usuelles et premières primitives
\item Primitives des fonctions polynomiales
\item Résolution de problèmes concrets : modélisation et interprétation
\item Activité d'intégration
\item Méthodes et exercices résolus
\item Exercices
\item Corrigés des exercices
\item Fiche bilan
\end{enumerate}
\end{multicols}
\end{sommaire}

\begin{objectifs}
\begin{itemize}
\item À la fin de cette leçon, je sais définir la dérivée d'une fonction en un point et l'interpréter géométriquement comme le coefficient directeur de la tangente.
\item Je sais appliquer les règles de dérivation (linéarité, produit, quotient, chaîne) pour dériver tout polynôme.
\item Je sais dresser et utiliser le tableau des dérivées usuelles pour les puissances entières et rationnelles.
\item Je sais définir une primitive d'une fonction et expliquer le rôle de la constante d'intégration.
\item Je sais déterminer la primitive d'une fonction polynôme terme à terme et vérifier le résultat par dérivation.
\item Je sais modéliser une situation concrète (coût, revenu, bénéfice) par une fonction polynôme et utiliser la dérivée pour optimiser.
\item Je sais rédiger une démarche mathématique complète, justifiant chaque étape et concluant de façon précise.
\item Je sais interpréter graphiquement la dérivée et la primitive (pente de la tangente, aire sous la courbe).
\end{itemize}
\end{objectifs}

\begin{prerequis}
\begin{itemize}
\item Notion de limite d'une fonction en un point (classe de 4ème / début Terminale).
\item Calcul algébrique : développement, factorisation, manipulation des puissances.
\item Représentation graphique de fonctions affines et quadratiques.
\item Notion de taux d'accroissement et de coefficient directeur d'une droite.
\item Utilisation basique de la notation f(x) et de la variable indépendante.
\end{itemize}
\end{prerequis}

\newpage

\coursec{Rappels sur la dérivée et interprétation géométrique}

\textbf{Situation d'accroche.} Dans un marché de Yaoundé, une vendeuse de beignets souhaite savoir combien de beignets elle doit produire chaque jour pour maximiser son bénéfice, sachant que le coût de fabrication augmente avec la quantité produite. Elle note que le coût total peut être modélisé par une fonction polynôme. Comment utiliser la dérivée pour trouver le \cle{coût marginal} (le coût de production d'un beignet supplémentaire) et la primitive pour retrouver le coût total à partir du coût marginal ? Cette situation illustre l'utilité concrète des dérivées et des primitives dans la gestion d'une petite entreprise.

\textbf{Problématique.} Comment mesurer la « vitesse de variation » d'une fonction en un point ? Comment ce nombre, appelé \cle{dérivée}, se traduit-il géométriquement sur la courbe ? C'est à ces questions que répond cette première section, avant d'apprendre à calculer des dérivées et des primitives dans les sections suivantes.

\courssub{Le taux d'accroissement : de la sécante à la tangente}

Considérons une fonction $f$ définie sur un intervalle $I$ et un nombre réel $a$ de $I$. Pour un réel $h$ non nul tel que $a+h$ appartienne à $I$, on peut comparer les images de $a$ et de $a+h$ : la quantité $f(a+h)-f(a)$ mesure la variation des images, et $h$ mesure la variation de la variable. Le rapport de ces deux variations s'appelle le \cle{taux d'accroissement} de $f$ entre $a$ et $a+h$ :
\[ \tau(h) = \frac{f(a+h)-f(a)}{h}. \]

Géométriquement, ce taux est la \textbf{pente} (ou coefficient directeur) de la droite passant par les points $A\big(a\,;\,f(a)\big)$ et $M\big(a+h\,;\,f(a+h)\big)$ de la courbe : c'est la \emph{sécante} $(AM)$.

L'idée fondatrice du calcul différentiel est la suivante : lorsque $h$ devient de plus en plus petit, le point $M$ se rapproche du point $A$ en glissant le long de la courbe, et la sécante $(AM)$ pivote autour de $A$ pour se rapprocher d'une position limite : la \cle{tangente} à la courbe au point $A$. La pente de cette tangente est la limite du taux d'accroissement.

\begin{definition}[Nombre dérivé de $f$ en $a$]
Soit $f$ une fonction définie sur un intervalle $I$ et $a$ un réel de $I$. On dit que $f$ est \cle{dérivable} en $a$ si le taux d'accroissement $\dfrac{f(a+h)-f(a)}{h}$ admet une limite finie lorsque $h$ tend vers $0$. Cette limite s'appelle le \cle{nombre dérivé} de $f$ en $a$ et se note $f'(a)$ :
\[ f'(a) = \lim_{h \to 0} \frac{f(a+h)-f(a)}{h}. \]
Si cette limite n'existe pas (ou est infinie), on dit que $f$ n'est \textbf{pas dérivable} en $a$.
\end{definition}

\begin{aretenir}[Notations de la dérivée]
Le nombre dérivé de $f$ en $a$ se note de plusieurs façons équivalentes :
\begin{itemize}
\item $f'(a)$ : notation de Lagrange, la plus utilisée au lycée ;
\item $\dfrac{df}{dx}(a)$ ou $\left.\dfrac{df}{dx}\right|_{x=a}$ : notation de Leibniz, très utilisée en physique et dans les métiers techniques car elle précise la variable de dérivation ;
\item $Df(a)$ : notation d'Euler, plus rare.
\end{itemize}
On lit indifféremment « $f$ prime de $a$ » ou « dérivée de $f$ par rapport à $x$ en $a$ ».
\end{aretenir}

\begin{attention}[$f'(a)$ n'est pas $f(a)$]
Erreur très fréquente au Probatoire et au Baccalauréat technique : confondre $f(a)$ et $f'(a)$.
\begin{itemize}
\item $f(a)$ est l'\textbf{image} de $a$ : c'est l'ordonnée du point de la courbe d'abscisse $a$.
\item $f'(a)$ est la \textbf{pente de la tangente} en ce point : c'est un taux de variation, pas une image.
\end{itemize}
Par exemple, pour $f(x)=x^2$ : $f(1)=1$ (le point de contact a pour ordonnée $1$) et $f'(1)=2$ (la tangente monte de $2$ unités quand on avance de $1$ unité). Ce sont deux nombres différents qui ne jouent pas le même rôle.
\end{attention}

\courssub{Interprétation géométrique : la pente de la tangente}

\begin{propriete}[Tangente et nombre dérivé]
Si $f$ est dérivable en $a$, alors la courbe représentative de $f$ admet au point $A\big(a\,;\,f(a)\big)$ une tangente dont le coefficient directeur (la pente) est $f'(a)$. Une équation de cette tangente est :
\[ y = f'(a)\,(x-a) + f(a). \]
\end{propriete}

Cette formule est exigible au Baccalauréat technique : elle combine les deux nombres $f(a)$ (position du point de contact) et $f'(a)$ (inclinaison de la tangente). Retenons l'intuition : la tangente est la droite qui « épouse » le mieux la courbe au voisinage de $A$ ; localement, la courbe et sa tangente sont presque confondues.

\begin{popfigure}
\begin{center}
\begin{tikzpicture}
\begin{axis}[
  width=0.9\linewidth, height=7cm,
  axis lines=middle,
  xlabel={$x$}, ylabel={$y$},
  xmin=-1.2, xmax=3.2, ymin=-0.8, ymax=5.2,
  xtick={1,2}, ytick={1,2,4},
  grid=both, grid style={popline!40},
  legend style={at={(0.02,0.98)}, anchor=north west, font=\small}
]
\addplot[popcurve, domain=-1.1:2.3, samples=200] {x^2};
\addlegendentry{$\mathcal{C}_f : y=x^2$}
\addplot[poporange, very thick, domain=-0.3:2.6, samples=2] {2*x-1};
\addlegendentry{Tangente en $x=1$ : $y=2x-1$}
\addplot[only marks, mark=*, mark size=2.5pt, popdark] coordinates {(1,1)};
\node[popdark, above right, font=\small] at (axis cs:1,1) {$A(1\,;\,f(1))$};
% Triangle de pente
\draw[popgreen, very thick] (axis cs:1,1) -- (axis cs:2,1) node[midway, below, font=\small] {$1$};
\draw[popgreen, very thick] (axis cs:2,1) -- (axis cs:2,3) node[midway, right, font=\small] {$f'(1)=2$};
\draw[popgreen, dashed] (axis cs:2,3) -- (axis cs:1,1);
\end{axis}
\end{tikzpicture}
\end{center}
\legende{Courbe de $f(x)=x^2$ et sa tangente au point $A(1\,;\,1)$. Le triangle vert illustre la pente : en avançant de $1$ en abscisse, la tangente monte de $2$, donc $f'(1)=2$.}
\end{popfigure}

\begin{savaistu}[Newton et Leibniz, les pères de la dérivée]
Au XVII\textsuperscript{e} siècle, Isaac Newton en Angleterre et Gottfried Wilhelm Leibniz en Allemagne ont formalisé \emph{indépendamment} le calcul différentiel. Newton cherchait à décrire la vitesse instantanée des planètes (ses « fluxions »), tandis que Leibniz inventait les notations $\dfrac{df}{dx}$ que nous utilisons encore aujourd'hui. Leur querelle de priorité a divisé les mathématiciens européens pendant des décennies ! Aujourd'hui, la dérivée est un outil quotidien de l'ingénieur : vitesse d'un véhicule, intensité d'un courant, coût marginal en gestion, débit d'une pompe...
\end{savaistu}

\courssub{Calculer une dérivée par la définition}

\begin{methode}[Calculer $f'(a)$ par la définition]
Pour calculer le nombre dérivé de $f$ en $a$ en revenant à la définition :
\begin{enumerate}
\item \textbf{Écrire} le taux d'accroissement $\tau(h)=\dfrac{f(a+h)-f(a)}{h}$ en remplaçant par l'expression de $f$.
\item \textbf{Simplifier} le numérateur, puis factoriser par $h$ et simplifier la fraction (le facteur $h$ au dénominateur doit disparaître).
\item \textbf{Faire tendre} $h$ vers $0$ dans l'expression simplifiée : la valeur obtenue est $f'(a)$.
\item \textbf{Conclure} : si la limite est un nombre fini, $f$ est dérivable en $a$ et $f'(a)$ est ce nombre.
\end{enumerate}
\end{methode}

\begin{experience}[Démonstration guidée : dérivée de $f(x)=x^2$ en $x=1$]
Appliquons la méthode pas à pas pour $f(x)=x^2$ et $a=1$.
\begin{enumerate}
\item \textbf{Taux d'accroissement.} On calcule $f(1+h)=(1+h)^2=1+2h+h^2$ et $f(1)=1$. Donc :
\[ \tau(h)=\frac{f(1+h)-f(1)}{h}=\frac{1+2h+h^2-1}{h}=\frac{2h+h^2}{h}. \]
\item \textbf{Simplification.} On factorise le numérateur par $h$ : $2h+h^2=h(2+h)$. Comme $h\neq 0$, on peut simplifier :
\[ \tau(h)=\frac{h(2+h)}{h}=2+h. \]
\item \textbf{Passage à la limite.} Lorsque $h$ tend vers $0$, $2+h$ tend vers $2$.
\item \textbf{Conclusion.} $f$ est dérivable en $1$ et $\boxed{f'(1)=2}$. La tangente à la parabole au point $A(1\,;\,1)$ a pour pente $2$ : c'est bien la droite $y=2(x-1)+1$, c'est-à-dire $y=2x-1$, tracée sur la figure ci-dessus.
\end{enumerate}
\end{experience}

\begin{exoresolu}[Dérivée de $f(x)=3x+2$ en $x=1$]
Énoncé. Soit $f$ la fonction définie sur $\mathbb{R}$ par $f(x)=3x+2$. En utilisant la définition, calculer $f'(1)$, puis donner l'équation de la tangente à la courbe de $f$ au point d'abscisse $1$.
\tcblower
\textbf{Solution détaillée.}
\begin{enumerate}
\item \textbf{Taux d'accroissement.} On a $f(1+h)=3(1+h)+2=5+3h$ et $f(1)=5$. Donc :
\[ \tau(h)=\frac{f(1+h)-f(1)}{h}=\frac{5+3h-5}{h}=\frac{3h}{h}=3 \quad (h\neq 0). \]
\item \textbf{Limite.} Le taux est constant, égal à $3$ : sa limite quand $h\to 0$ est $3$.
\item \textbf{Conclusion.} $f$ est dérivable en $1$ et $f'(1)=3$.
\end{enumerate}
\textbf{Tangente.} $y=f'(1)(x-1)+f(1)=3(x-1)+5$, soit $y=3x+2$. On retrouve la droite elle-même : c'est normal, la courbe d'une fonction affine est une droite, qui est sa propre tangente en tout point ! Sa pente est le coefficient de $x$, ici $3$.
\end{exoresolu}

\courssub{Deux dérivées fondamentales : constante et identité}

\begin{propriete}[Dérivée d'une constante et de l'identité]
\begin{itemize}
\item Si $f$ est une fonction \textbf{constante}, $f(x)=k$ (avec $k\in\mathbb{R}$), alors $f$ est dérivable sur $\mathbb{R}$ et $f'(x)=0$ pour tout réel $x$.
\item Si $f$ est la fonction \textbf{identité}, $f(x)=x$, alors $f$ est dérivable sur $\mathbb{R}$ et $f'(x)=1$ pour tout réel $x$.
\end{itemize}
Cette démonstration est exigible : elle se fait par la définition.
\end{propriete}

\begin{experience}[Démonstrations par la définition]
\textbf{1. Fonction constante.} Soit $f(x)=k$ pour tout $x$. Pour tout réel $a$ et tout $h\neq 0$ :
\[ \tau(h)=\frac{f(a+h)-f(a)}{h}=\frac{k-k}{h}=\frac{0}{h}=0. \]
Le taux est nul pour tout $h$, donc sa limite est $0$ : $f'(a)=0$. Géométriquement, la courbe est une droite horizontale : sa tangente est horizontale, de pente nulle.

\textbf{2. Fonction identité.} Soit $f(x)=x$. Pour tout réel $a$ et tout $h\neq 0$ :
\[ \tau(h)=\frac{f(a+h)-f(a)}{h}=\frac{(a+h)-a}{h}=\frac{h}{h}=1. \]
Le taux vaut constamment $1$, donc $f'(a)=1$. Géométriquement, la courbe est la droite $y=x$, de pente $1$ : sa tangente en tout point est elle-même.
\end{experience}

\begin{attention}[Conditions d'application]
\begin{itemize}
\item Dans le calcul de $\tau(h)$, on ne peut simplifier par $h$ que parce que $h\neq 0$ : la limite se prend pour $h$ \emph{tendant vers} $0$, jamais égal à $0$. Écrire $\tau(0)$ est une erreur : le dénominateur serait nul.
\item Une fonction constante a une dérivée nulle, mais attention à la rédaction : on écrit $f'(x)=0$ et non $f'(x)=k$.
\end{itemize}
\end{attention}

\courssub{Le lien entre dérivée et sens de variation}

Revenons à notre vendeuse de beignets : si le coût total augmente quand la quantité produite augmente, cela se voit sur la courbe du coût — elle « monte » — et cela se lit sur la dérivée — elle est positive. Ce lien entre le \textbf{signe de la dérivée} et le \textbf{sens de variation} de la fonction est l'un des résultats les plus utiles de l'année.

\begin{propriete}[Signe de la dérivée et variations]
Soit $f$ une fonction dérivable sur un intervalle $I$.
\begin{itemize}
\item Si $f'(x)>0$ pour tout $x$ de $I$ (sauf éventuellement en des points isolés où $f'$ s'annule), alors $f$ est \textbf{strictement croissante} sur $I$.
\item Si $f'(x)<0$ pour tout $x$ de $I$ (sauf éventuellement en des points isolés), alors $f$ est \textbf{strictement décroissante} sur $I$.
\item Si $f'(x)=0$ pour tout $x$ de $I$, alors $f$ est \textbf{constante} sur $I$.
\end{itemize}
\end{propriete}

L'intuition géométrique est immédiate : une tangente de pente positive « monte », donc la courbe monte ; une tangente de pente négative « descend » ; une tangente horizontale (pente nulle) signale un point où la courbe cesse momentanément de monter ou de descendre — souvent un maximum ou un minimum, précisément ce que cherche notre vendeuse pour optimiser son bénéfice.

\begin{exemplebox}[Lecture graphique du signe de la dérivée]
Pour $f(x)=x^2$, on a vu que $f'(1)=2>0$ : au voisinage de $x=1$, la parabole est croissante, ce que confirme la figure. En $x=0$, la tangente est horizontale : $f'(0)=0$, et la fonction admet un minimum en $0$ (le sommet de la parabole). Pour $x<0$, les tangentes « descendent » : $f'(x)<0$ et $f$ est décroissante sur $]-\infty\,;\,0]$.
\end{exemplebox}

\begin{exercice}[À vous de jouer]
\begin{enumerate}
\item Soit $g(x)=x^2$. En utilisant la définition, calculer $g'(3)$, puis déterminer l'équation de la tangente à la courbe de $g$ au point d'abscisse $3$.
\item Soit $h(x)=-2x+5$. Calculer $h'(a)$ pour un réel $a$ quelconque. Le résultat dépend-il de $a$ ? Interpréter géométriquement.
\item La courbe d'une fonction $f$ admet au point d'abscisse $2$ une tangente d'équation $y=4x-3$. Déterminer $f(2)$ et $f'(2)$.
\end{enumerate}
\end{exercice}

\begin{corrige}[Exercice 1]
\textbf{1.} $g(3+h)=(3+h)^2=9+6h+h^2$ et $g(3)=9$, donc $\tau(h)=\dfrac{6h+h^2}{h}=6+h$ pour $h\neq 0$. Quand $h\to 0$, $\tau(h)\to 6$ : $g'(3)=6$. Tangente : $y=6(x-3)+9$, soit $y=6x-9$.

\textbf{2.} $\tau(h)=\dfrac{-2(a+h)+5-(-2a+5)}{h}=\dfrac{-2h}{h}=-2$. Donc $h'(a)=-2$ pour tout $a$ : le résultat ne dépend pas de $a$. La courbe est une droite de pente $-2$, tangente à elle-même en tout point ; la fonction est strictement décroissante sur $\mathbb{R}$ car $h'(a)<0$.

\textbf{3.} La tangente en $x=2$ a pour équation $y=f'(2)(x-2)+f(2)$. Par identification avec $y=4x-3=4(x-2)+5$ : $f'(2)=4$ et $f(2)=5$. (Vérification : le point de contact $(2\,;\,5)$ appartient bien à la tangente : $4\times 2-3=5$.)
\end{corrige}

\begin{aretenir}[L'essentiel de la section]
\begin{itemize}
\item $f'(a)=\displaystyle\lim_{h\to 0}\frac{f(a+h)-f(a)}{h}$ : c'est la limite du taux d'accroissement.
\item Géométriquement, $f'(a)$ est la \textbf{pente de la tangente} à la courbe au point $A\big(a\,;\,f(a)\big)$, d'équation $y=f'(a)(x-a)+f(a)$.
\item Fonction constante $\Rightarrow$ dérivée nulle ; fonction identité $\Rightarrow$ dérivée égale à $1$.
\item Signe de $f'$ $\Longleftrightarrow$ sens de variation de $f$ : $f'>0$ donne $f$ croissante, $f'<0$ donne $f$ décroissante.
\item Ne jamais confondre $f(a)$ (image) et $f'(a)$ (pente).
\end{itemize}
\end{aretenir}

\coursec{Règles de dérivation usuelles}

Dans la section précédente, nous avons défini le nombre dérivé $f'(a)$ comme la limite du taux d'accroissement, et nous l'avons interprété géométriquement comme le coefficient directeur de la tangente à la courbe au point d'abscisse $a$. Calculer cette limite à chaque fois serait long et fastidieux : imaginez un technicien en électricité qui devrait refaire tout le calcul de limite chaque fois qu'il étudie une tension $u(t)$ ! Heureusement, les mathématiciens ont établi des \cle{règles de dérivation} : une fois les dérivées des fonctions usuelles connues, ces règles permettent de dériver \emph{n'importe quelle} combinaison de ces fonctions (somme, produit, quotient, composée) \textbf{sans revenir à la définition}. C'est l'objet de cette section.

\courssub{Linéarité de la dérivation}

\textbf{Intuition.} Si une recette de chantier double les quantités de ciment, le coût double. Si l'on additionne deux sources de coût, le coût total est la somme des coûts. La dérivation, qui mesure une \emph{vitesse de variation}, se comporte de la même façon : multiplier une fonction par un nombre multiplie sa vitesse de variation par ce même nombre, et la vitesse de variation d'une somme est la somme des vitesses. On dit que la dérivation est \cle{linéaire}.

\begin{propriete}[Linéarité de la dérivation]
Soient $f$ et $g$ deux fonctions dérivables sur un intervalle $I$, et $\alpha$, $\beta$ deux nombres réels. Alors la fonction $\alpha f + \beta g$ est dérivable sur $I$ et :
\[ (\alpha f + \beta g)' = \alpha f' + \beta g'. \]
En particulier : $(f+g)' = f' + g'$ et $(\alpha f)' = \alpha f'$.
\end{propriete}

\begin{experience}[Démonstration guidée : la dérivée d'une somme]
Nous démontrons que $(f+g)'(a) = f'(a) + g'(a)$ à partir de la définition du nombre dérivé. Suivez les étapes.

\begin{enumerate}
\item \textbf{Écrire le taux d'accroissement de $f+g$ entre $a$ et $a+h$.}
\[ \tau(h) = \frac{(f+g)(a+h) - (f+g)(a)}{h} = \frac{f(a+h) + g(a+h) - f(a) - g(a)}{h}. \]
\item \textbf{Regrouper astucieusement les termes} en séparant ce qui concerne $f$ de ce qui concerne $g$ :
\[ \tau(h) = \underbrace{\frac{f(a+h) - f(a)}{h}}_{\text{taux de } f} + \underbrace{\frac{g(a+h) - g(a)}{h}}_{\text{taux de } g}. \]
\item \textbf{Passer à la limite.} Quand $h$ tend vers $0$, le premier taux tend vers $f'(a)$ (car $f$ est dérivable en $a$) et le second vers $g'(a)$. La limite d'une somme étant la somme des limites :
\[ \lim_{h \to 0} \tau(h) = f'(a) + g'(a). \]
\end{enumerate}
La fonction $f+g$ est donc dérivable en $a$ et $(f+g)'(a) = f'(a) + g'(a)$. Le cas de $(\alpha f)' = \alpha f'$ se traite de même, en factorisant $\alpha$ dans le taux d'accroissement. \hfill $\blacksquare$
\end{experience}

\begin{exemplebox}[Application directe]
Soient $f(x) = x^2$ et $g(x) = x^3$. On sait que $f'(x) = 2x$ et $g'(x) = 3x^2$. Alors :
\[ (3f - 2g)'(x) = 3f'(x) - 2g'(x) = 3 \times 2x - 2 \times 3x^2 = 6x - 6x^2. \]
\end{exemplebox}

\courssub{Dérivée de la fonction puissance}

Avant d'aller plus loin, rappelons la règle fondamentale qui permet de dériver les monômes, briques de base des polynômes.

\begin{propriete}[Dérivée de $x^n$]
Pour tout entier relatif $n$ non nul, la fonction $x \mapsto x^n$ est dérivable sur $\mathbb{R}$ si $n > 0$ (sur $\mathbb{R}^*$ si $n < 0$), et :
\[ (x^n)' = n\,x^{n-1}. \]
Cas particuliers à connaître par cœur : $(x)' = 1$, $(x^2)' = 2x$, $(x^3)' = 3x^2$, et pour $n = -1$ : $\left(\dfrac{1}{x}\right)' = -\dfrac{1}{x^2}$.
\end{propriete}

\begin{aretenir}[Idée de la preuve pour $n$ entier positif]
Le taux d'accroissement de $x^n$ en $a$ s'écrit $\dfrac{(a+h)^n - a^n}{h}$. En développant $(a+h)^n$ (formule du binôme de Newton), le terme $a^n$ s'élimine, le terme suivant est $n\,a^{n-1}h$, et tous les autres termes contiennent au moins $h^2$. Après simplification par $h$, il reste $n\,a^{n-1}$ plus des termes qui tendent vers $0$. D'où la formule.
\end{aretenir}

\begin{methode}[Dériver un polynôme terme à terme]
Pour dériver un polynôme, on combine la linéarité et la règle de la puissance :
\begin{enumerate}
\item On \textbf{décompose} le polynôme en une somme de monômes de la forme $a_k x^k$.
\item On \textbf{dérive chaque monôme} séparément : $(a_k x^k)' = a_k \times k\,x^{k-1}$ (le coefficient $a_k$ est conservé, l'exposant descend d'un cran).
\item On \textbf{additionne} les résultats. Attention : un terme constant a pour dérivée $0$.
\end{enumerate}
\end{methode}

\begin{exemplebox}[Exemple chiffré fondamental]
Dériver $f(x) = 2x^3 - 5x + 4$.
\begin{align*}
f'(x) &= (2x^3)' - (5x)' + (4)' \\
&= 2 \times 3x^2 - 5 \times 1 + 0 \\
&= 6x^2 - 5.
\end{align*}
Le terme constant $4$ disparaît : une fonction constante ne varie pas, sa vitesse de variation est nulle.
\end{exemplebox}

\courssub{Règle du produit}

\textbf{Intuition.} Un atelier fabrique des pièces : le chiffre d'affaires est le \emph{produit} du prix unitaire par la quantité vendue. Si le prix augmente et la quantité augmente, la variation du chiffre d'affaires dépend des \emph{deux} variations : il faut tenir compte de l'augmentation du prix (appliquée à la quantité actuelle) \emph{et} de l'augmentation de la quantité (appliquée au prix actuel). C'est exactement le contenu de la règle du produit.

\begin{propriete}[Dérivée d'un produit]
Soient $f$ et $g$ deux fonctions dérivables sur un intervalle $I$. Alors le produit $fg$ est dérivable sur $I$ et :
\[ (fg)' = f'g + fg'. \]
\end{propriete}

\begin{attention}[Piège classique]
La dérivée d'un produit n'est \textbf{pas} le produit des dérivées : $(fg)' \neq f'g'$ en général ! Par exemple, avec $f(x) = x$ et $g(x) = x$ : on a $(fg)(x) = x^2$, donc $(fg)'(x) = 2x$, alors que $f'(x)g'(x) = 1 \times 1 = 1$. Vérifiez toujours que vous appliquez $(fg)' = f'g + fg'$.
\end{attention}

\begin{experience}[Démonstration guidée : la règle du produit]
Nous démontrons que $(fg)'(a) = f'(a)g(a) + f(a)g'(a)$. L'astuce consiste à \textbf{ajouter et soustraire} la même quantité $f(a+h)g(a)$.

\begin{enumerate}
\item \textbf{Écrire le taux d'accroissement du produit :}
\[ \tau(h) = \frac{f(a+h)g(a+h) - f(a)g(a)}{h}. \]
\item \textbf{Ajouter et soustraire $f(a+h)g(a)$ au numérateur} (cela ne change rien, puisqu'on ajoute $0$) :
\[ \tau(h) = \frac{f(a+h)g(a+h) - f(a+h)g(a) + f(a+h)g(a) - f(a)g(a)}{h}. \]
\item \textbf{Factoriser par paquets de deux termes :}
\[ \tau(h) = f(a+h) \times \frac{g(a+h) - g(a)}{h} + g(a) \times \frac{f(a+h) - f(a)}{h}. \]
\item \textbf{Passer à la limite} quand $h$ tend vers $0$ :
\begin{itemize}
\item $f(a+h) \to f(a)$ (une fonction dérivable est continue) ;
\item $\dfrac{g(a+h)-g(a)}{h} \to g'(a)$ et $\dfrac{f(a+h)-f(a)}{h} \to f'(a)$ ;
\item $g(a)$ reste fixe.
\end{itemize}
On obtient : $\displaystyle\lim_{h \to 0} \tau(h) = f(a)g'(a) + g(a)f'(a)$. \hfill $\blacksquare$
\end{enumerate}
\end{experience}

\begin{exoresolu}[Dériver un produit de deux façons]
Énoncé. Soit $f(x) = (x^2 + 1)(x^3 - 2x)$. Calculer $f'(x)$ de deux manières : en développant d'abord, puis avec la règle du produit.
\tcblower
Solution détaillée.

\textbf{Méthode 1 : développement puis dérivation terme à terme.}
\[ f(x) = x^5 - 2x^3 + x^3 - 2x = x^5 - x^3 - 2x, \]
donc $f'(x) = 5x^4 - 3x^2 - 2$.

\textbf{Méthode 2 : règle du produit.} Posons $u(x) = x^2 + 1$ et $v(x) = x^3 - 2x$. Alors $u'(x) = 2x$ et $v'(x) = 3x^2 - 2$.
\begin{align*}
f'(x) &= u'(x)v(x) + u(x)v'(x) \\
&= 2x(x^3 - 2x) + (x^2 + 1)(3x^2 - 2) \\
&= 2x^4 - 4x^2 + 3x^4 - 2x^2 + 3x^2 - 2 \\
&= 5x^4 - 3x^2 - 2.
\end{align*}
Les deux méthodes donnent le même résultat, ce qui conforte la règle du produit. Notez que pour un produit compliqué (par exemple avec une racine carrée), le développement serait impossible : la règle du produit devient alors indispensable.
\end{exoresolu}

\begin{popfigure}
\begin{center}
\begin{tikzpicture}[
  grow=down,
  level distance=1.5cm,
  level 1/.style={sibling distance=7cm},
  level 2/.style={sibling distance=3.6cm},
  every node/.style={draw, rounded corners=2pt, fill=popblueL, inner sep=3pt, font=\small},
  edge from parent/.style={draw=popdark, -{Stealth}, thick}
]
\node[fill=poporangeL] {$f = u \times v$ \;:\; $(x^2+1)(x^3-2x)$}
  child { node {$u = x^2 + 1$}
    child { node[fill=popgreenL] {$u' = 2x$} }
  }
  child { node {$v = x^3 - 2x$}
    child { node[fill=popgreenL] {$v' = 3x^2 - 2$} }
  };
\node[draw=none, fill=none, font=\small\itshape, text=popdark] at (0,-4.6)
  {Règle du produit : $f' = u'v + uv' = 2x(x^3-2x) + (x^2+1)(3x^2-2) = 5x^4 - 3x^2 - 2$};
\end{tikzpicture}
\end{center}
\legende{Arbre de dérivation de $f(x) = (x^2+1)(x^3-2x)$ : on identifie le produit, on dérive chaque facteur par linéarité (feuilles vertes), puis on assemble avec la règle $f' = u'v + uv'$.}
\end{popfigure}

\courssub{Règle du quotient}

\textbf{Intuition.} En gestion, on s'intéresse souvent à des \emph{rapports} : coût moyen (coût total divisé par la quantité), rendement, vitesse moyenne. Quand le numérateur et le dénominateur varient tous les deux, il faut une règle spécifique.

\begin{propriete}[Dérivée d'un quotient]
Soient $f$ et $g$ deux fonctions dérivables sur un intervalle $I$, avec $g(x) \neq 0$ pour tout $x$ de $I$. Alors le quotient $\dfrac{f}{g}$ est dérivable sur $I$ et :
\[ \left(\frac{f}{g}\right)' = \frac{f'g - fg'}{g^2}. \]
\end{propriete}

\begin{attention}[Conditions d'application et ordre des termes]
Deux pièges à éviter :
\begin{itemize}
\item \textbf{Condition $g \neq 0$} : avant de dériver un quotient, vérifiez que le dénominateur ne s'annule pas sur l'intervalle étudié.
\item \textbf{Ordre au numérateur} : c'est $f'g - fg'$ (la dérivée du \emph{haut} fois le bas, \emph{moins} le haut fois la dérivée du bas). Inverser l'ordre change le signe du résultat ! Une aide-mémoire : « la dérivée du haut descend, moins la dérivée du bas qui monte, le tout sur le bas au carré ».
\end{itemize}
\end{attention}

\begin{experience}[Démonstration guidée : la règle du quotient]
Nous démontrons que $\left(\frac{f}{g}\right)'(a) = \frac{f'(a)g(a) - f(a)g'(a)}{g(a)^2}$ pour $g(a) \neq 0$.

\begin{enumerate}
\item \textbf{Écrire le taux d'accroissement du quotient} et mettre au même dénominateur :
\[ \tau(h) = \frac{\frac{f(a+h)}{g(a+h)} - \frac{f(a)}{g(a)}}{h} = \frac{f(a+h)g(a) - f(a)g(a+h)}{h\,g(a+h)g(a)}. \]
\item \textbf{Ajouter et soustraire $f(a)g(a)$ au numérateur} :
\[ \tau(h) = \frac{\big(f(a+h)g(a) - f(a)g(a)\big) - \big(f(a)g(a+h) - f(a)g(a)\big)}{h\,g(a+h)g(a)}. \]
\item \textbf{Factoriser} $g(a)$ dans le premier paquet et $f(a)$ dans le second :
\[ \tau(h) = \frac{g(a)\frac{f(a+h)-f(a)}{h} - f(a)\frac{g(a+h)-g(a)}{h}}{g(a+h)g(a)}. \]
\item \textbf{Passer à la limite} $h \to 0$ :
\begin{itemize}
\item $g(a+h) \to g(a)$ (continuité de $g$) ;
\item Les taux d'accroissement tendent vers $f'(a)$ et $g'(a)$.
\end{itemize}
On obtient : $\displaystyle\lim_{h \to 0} \tau(h) = \frac{g(a)f'(a) - f(a)g'(a)}{g(a)^2}$. \hfill $\blacksquare$
\end{enumerate}
\end{experience}

\begin{exoresolu}[Dériver un quotient]
Énoncé. Soit $f(x) = \dfrac{3x + 1}{x - 2}$ pour $x \neq 2$. Calculer $f'(x)$.
\tcblower
Solution détaillée. Posons $u(x) = 3x + 1$ et $v(x) = x - 2$. Alors $u'(x) = 3$ et $v'(x) = 1$. La condition $v(x) \neq 0$ est satisfaite sur $\mathbb{R} \setminus \{2\}$.
\begin{align*}
f'(x) &= \frac{u'(x)v(x) - u(x)v'(x)}{[v(x)]^2} \\
&= \frac{3(x - 2) - (3x + 1) \times 1}{(x - 2)^2} \\
&= \frac{3x - 6 - 3x - 1}{(x - 2)^2} \\
&= \frac{-7}{(x - 2)^2}.
\end{align*}
Comme $(x-2)^2 > 0$ pour tout $x \neq 2$, on en déduit au passage que $f'(x) < 0$ : la fonction $f$ est strictement décroissante sur chacun des intervalles de son ensemble de définition.
\end{exoresolu}

\courssub{Règle de la chaîne (fonctions composées)}

\textbf{Intuition.} Sur un chantier, une poulie double la force, et un treuil triple l'effet de la poulie : l'effet total est le \emph{produit} des effets, $2 \times 3 = 6$. De même, quand deux fonctions sont composées, les vitesses de variation se \emph{multiplient}.

\begin{definition}[Fonction composée]
Soient $g$ une fonction définie sur un intervalle $I$ et $f$ une fonction définie sur un intervalle contenant $g(I)$. La \cle{composée} de $g$ par $f$, notée $f \circ g$, est définie par :
\[ (f \circ g)(x) = f\big(g(x)\big). \]
On applique d'abord $g$, puis $f$ au résultat.
\end{definition}

\begin{propriete}[Dérivée d'une composée (règle de la chaîne)]
Si $g$ est dérivable sur $I$ et $f$ est dérivable sur $g(I)$, alors $f \circ g$ est dérivable sur $I$ et :
\[ (f \circ g)'(x) = f'\big(g(x)\big) \times g'(x). \]
On dérive « la fonction extérieure » en laissant l'intérieur inchangé, puis on multiplie par la dérivée de « la fonction intérieure ».
\end{propriete}

\begin{aretenir}[Cas particulier très utile : la puissance d'une fonction]
En prenant $f(u) = u^n$, on obtient pour tout entier $n \geq 1$ :
\[ \big(u^n\big)' = n\,u^{n-1} \times u'. \]
Par exemple : $\big((x^2+1)^3\big)' = 3(x^2+1)^2 \times 2x = 6x(x^2+1)^2$.
\end{aretenir}

\begin{attention}[L'erreur la plus fréquente au Baccalauréat technique]
Oublier le facteur $g'(x)$ dans la dérivation d'une composée ! Par exemple, pour dériver $h(x) = (x^2 + 1)^3$ :
\begin{itemize}
\item \textbf{Faux} : $h'(x) = 3(x^2+1)^2$ (on a oublié de multiplier par la dérivée de l'intérieur).
\item \textbf{Correct} : $h'(x) = 3(x^2+1)^2 \times 2x = 6x(x^2+1)^2$.
\end{itemize}
Réflexe à acquérir : dès que vous voyez « quelque chose élevé à une puissance » où ce quelque chose contient $x$, demandez-vous : « ai-je bien multiplié par la dérivée de l'intérieur ? »
\end{attention}

\begin{experience}[Démonstration guidée : la règle de la chaîne]
Nous justifions $(f \circ g)'(a) = f'(g(a))\,g'(a)$ sous l'hypothèse que $g(a+h) \neq g(a)$ pour $h$ assez petit (le cas général se traite par une fonction auxiliaire).

\begin{enumerate}
\item \textbf{Écrire le taux d'accroissement de la composée} et multiplier par $1 = \frac{g(a+h)-g(a)}{g(a+h)-g(a)}$ :
\[ \tau(h) = \frac{f(g(a+h)) - f(g(a))}{h} = \frac{f(g(a+h)) - f(g(a))}{g(a+h)-g(a)} \times \frac{g(a+h)-g(a)}{h}. \]
\item \textbf{Identifier les deux facteurs.} Le second facteur est le taux d'accroissement de $g$ en $a$. Pour le premier, posons $k = g(a+h)-g(a)$. Comme $g$ est continue, $k \to 0$ quand $h \to 0$. Ce facteur s'écrit $\frac{f(g(a)+k) - f(g(a))}{k}$, c'est le taux d'accroissement de $f$ en $g(a)$.
\item \textbf{Passer à la limite.} Quand $h \to 0$ :
\begin{itemize}
\item $\frac{g(a+h)-g(a)}{h} \to g'(a)$ ;
\item $\frac{f(g(a)+k) - f(g(a))}{k} \to f'(g(a))$.
\end{itemize}
Le produit tend vers $f'(g(a))\,g'(a)$. \hfill $\blacksquare$
\end{enumerate}
\end{experience}

\begin{exoresolu}[Dériver une fonction composée]
Énoncé. Soit $h(x) = (2x - 5)^4$. Calculer $h'(x)$, puis déterminer $h'(3)$.
\tcblower
Solution détaillée. On reconnaît une composée : la fonction intérieure est $g(x) = 2x - 5$, la fonction extérieure est $f(u) = u^4$.

Étape 1 : dérivées des deux maillons. $g'(x) = 2$ et $f'(u) = 4u^3$.

Étape 2 : application de la règle de la chaîne.
\[ h'(x) = f'\big(g(x)\big) \times g'(x) = 4(2x-5)^3 \times 2 = 8(2x-5)^3. \]

Étape 3 : valeur numérique. $h'(3) = 8(2 \times 3 - 5)^3 = 8 \times 1^3 = 8$.

Interprétation : au point d'abscisse $3$, la tangente à la courbe de $h$ a pour coefficient directeur $8$.
\end{exoresolu}

\courssub{Synthèse : choisir la bonne règle}

Face à une fonction à dériver, la première question à se poser est : \emph{quelle est la structure de la fonction ?} Le tableau suivant résume la démarche.

\begin{center}
\begin{tabularx}{\linewidth}{|l|X|X|}
\hline
\rowcolor{popblueL} \textbf{Structure} & \textbf{Règle à appliquer} & \textbf{Exemple} \\
\hline
Somme / différence & Linéarité : $(\alpha f + \beta g)' = \alpha f' + \beta g'$ & $(2x^3 - 5x + 4)' = 6x^2 - 5$ \\
\hline
Produit & $(fg)' = f'g + fg'$ & $\big(x^2(x^3-2x)\big)'$ \\
\hline
Quotient & $\left(\dfrac{f}{g}\right)' = \dfrac{f'g - fg'}{g^2}$, $g \neq 0$ & $\left(\dfrac{3x+1}{x-2}\right)'$ \\
\hline
Composée & $(f \circ g)' = (f' \circ g) \times g'$ & $\big((2x-5)^4\big)' = 8(2x-5)^3$ \\
\hline
\end{tabularx}
\end{center}

\begin{aretenir}[Tableau récapitulatif des dérivées usuelles (fonctions algébriques)]
Les formules suivantes, conséquences directes des règles précédentes, doivent être connues par cœur pour le Probatoire / Baccalauréat technique.
\begin{center}
\begin{tabular}{|c|c|c|}
\hline
\rowcolor{popblueL} \textbf{Fonction $f(x)$} & \textbf{Domaine} & \textbf{Dérivée $f'(x)$} \\
\hline
$k$ (constante) & $\mathbb{R}$ & $0$ \\
\hline
$x^n$ ($n \in \mathbb{Z}^*$) & $\mathbb{R}$ ($n>0$), $\mathbb{R}^*$ ($n<0$) & $n x^{n-1}$ \\
\hline
$\sqrt{x} = x^{1/2}$ & $[0, +\infty[$ & $\dfrac{1}{2\sqrt{x}}$ (sur $]0,+\infty[$) \\
\hline
$\dfrac{1}{x} = x^{-1}$ & $\mathbb{R}^*$ & $-\dfrac{1}{x^2}$ \\
\hline
$u(x)^n$ ($n \in \mathbb{N}^*$) & Selon $u$ & $n\,u(x)^{n-1} \times u'(x)$ \\
\hline
\end{tabular}
\end{center}
\end{aretenir}

\begin{exercice}[À vous de jouer]
\begin{enumerate}
\item Dériver $f(x) = 4x^3 - 7x^2 + x - 9$.
\item Dériver $g(x) = (x + 3)(x^2 - 1)$ avec la règle du produit, puis vérifier en développant.
\item Dériver $h(x) = \dfrac{x^2}{x + 1}$ pour $x \neq -1$.
\item Dériver $k(x) = (3x^2 + 1)^5$.
\end{enumerate}
\end{exercice}

\begin{corrige}[Exercice : À vous de jouer]
\begin{enumerate}
\item $f'(x) = 12x^2 - 14x + 1$ (dérivation terme à terme ; le constant $-9$ disparaît).
\item Avec $u = x+3$, $v = x^2 - 1$ : $g'(x) = 1 \times (x^2 - 1) + (x+3) \times 2x = x^2 - 1 + 2x^2 + 6x = 3x^2 + 6x - 1$. Vérification : $g(x) = x^3 + 3x^2 - x - 3$, donc $g'(x) = 3x^2 + 6x - 1$. Identique.
\item $h'(x) = \dfrac{2x(x+1) - x^2 \times 1}{(x+1)^2} = \dfrac{x^2 + 2x}{(x+1)^2} = \dfrac{x(x+2)}{(x+1)^2}$.
\item Règle de la chaîne : $k'(x) = 5(3x^2+1)^4 \times 6x = 30x(3x^2+1)^4$.
\end{enumerate}
\end{corrige}

\begin{savaistu}[Lagrange, Leibniz et les notations de la dérivée]
La notation $f'(x)$ que nous utilisons est due au mathématicien franco-italien \textbf{Joseph-Louis Lagrange} (1736--1813), qui voyait la dérivée comme une nouvelle fonction « dérivée » de la première — d'où le nom ! Son contemporain allemand \textbf{Gottfried Wilhelm Leibniz} (1646--1716), co-inventeur du calcul différentiel avec Newton, utilisait la notation $\dfrac{dy}{dx}$, qui évoque le rapport de deux petites variations. La notation de Leibniz rend la règle de la chaîne presque évidente : si $y$ dépend de $u$ et $u$ dépend de $x$, alors $\dfrac{dy}{dx} = \dfrac{dy}{du} \times \dfrac{du}{dx}$ — on « simplifie » par $du$ comme avec des fractions. Attention toutefois : ce n'est qu'une aide-mémoire, $\dfrac{dy}{dx}$ n'est pas un vrai quotient ! Ces deux notations, nées d'une rivalité historique entre écoles anglaise et allemande, coexistent encore aujourd'hui dans tous les manuels du monde.
\end{savaistu}

\coursec{Tableau des dérivées usuelles et premières primitives}

Après avoir maîtrisé les \textbf{règles de dérivation usuelles} (linéarité, produit, quotient, composée) dans la section précédente, nous disposons maintenant d'une « boîte à outils » complète pour dériver la plupart des fonctions rencontrées en Terminale technique. Mais en technologie, en économie de la construction ou en étude de mouvements, il arrive souvent que l'on connaisse la \emph{vitesse} (la dérivée) et que l'on cherche la \emph{position} (la fonction de départ), ou que l'on connaisse le \emph{coût marginal} et que l'on veuille retrouver le \emph{coût total}. C'est précisément le rôle de la \textbf{primitive} : l'opération inverse de la dérivation. Cette section établit le tableau de correspondance qui permet de passer sans effort de la fonction à sa dérivée… et inversement.

\courssub{Le tableau de correspondance : dériver et primitiver}

L'idée centrale est simple : \textbf{dériver} et \textbf{primitiver} sont deux lectures d'un même tableau. Si l'on sait que la dérivée de $x^n$ est $n x^{n-1}$, alors on sait immédiatement qu'une primitive de $n x^{n-1}$ est $x^n$ (à une constante près). En Terminale technique, le programme officiel borne les \emph{primitives} aux fonctions \textbf{polynômes} (puissances entières naturelles), tandis que les \emph{dérivées usuelles} s'étendent aux puissances rationnelles utiles pour modéliser des situations concrètes (surfaces, volumes, lois de résistance des matériaux, décharges de condensateur, etc.).

\begin{propriete}[Tableau des dérivées usuelles (puissances rationnelles) et premières primitives (polynômes)]
Soit $n \in \mathbb{Q}$ (entier relatif ou fraction simple comme $\frac{1}{2}, -\frac{1}{2}$) pour la colonne de gauche, et $n \in \mathbb{N}$ pour la colonne des primitives (programme officiel). Soit $I$ un intervalle ne contenant pas les valeurs interdites (division par zéro, racine d'un négatif). On a le tableau de correspondance suivant :
\end{propriete}

\begin{center}
\begin{tabularx}{\linewidth}{|>{\columncolor{popblueL}}c|X|X|}\hline
\textbf{Fonction $f(x)$ sur $I$} & \textbf{Dérivée $f'(x)$} & \textbf{Primitive $F(x)$ de $f(x)$ ($F'=f$)} \\ \hline
$k$ (constante, $k \in \mathbb{R}$) & $0$ & $kx + C$ \\ \hline
$x^n \quad (n \in \mathbb{N},\ n \neq 0)$ & $n x^{n-1}$ & $\dfrac{x^{n+1}}{n+1} + C$ \\ \hline
$\sqrt{x} = x^{1/2}$ & $\dfrac{1}{2\sqrt{x}} = \frac{1}{2}x^{-1/2}$ & $\dfrac{2}{3}x^{3/2} + C = \frac{2}{3}x\sqrt{x} + C$ \\ \hline
$\dfrac{1}{x} = x^{-1}$ & $-\dfrac{1}{x^2} = -x^{-2}$ & $\ln|x| + C$ \textit{(culture générale, hors programme primitives)} \\ \hline
$\dfrac{1}{\sqrt{x}} = x^{-1/2}$ & $-\dfrac{1}{2}x^{-3/2}$ & $2\sqrt{x} + C = 2x^{1/2} + C$ \\ \hline
$\dfrac{1}{x^2} = x^{-2}$ & $-\dfrac{2}{x^3} = -2x^{-3}$ & $-\dfrac{1}{x} + C = -x^{-1} + C$ \\ \hline
\end{tabularx}
\end{center}

\begin{attention}[Lecture dans les deux sens et pièges classiques]
\begin{itemize}
    \item \textbf{Colonne du milieu (Dériver) :} On part de la ligne de gauche, on applique la formule « exposant $\times$ $x^{\text{exposant}-1}$ ». Valable pour $n \in \mathbb{Q}$ sur l'intervalle adéquat.
    \item \textbf{Colonne de droite (Primitiver) :} \textbf{Restriction programme :} On ne primitivise \textbf{que les puissances entières naturelles} ($n \in \mathbb{N}$) dans le cadre des polynômes. La règle est : \textbf{on augmente l'exposant de $1$ et on divise par le nouvel exposant}. Les lignes $\sqrt{x}$, $1/x$, $1/\sqrt{x}$, $1/x^2$ sont données pour la \emph{dérivée} (colonne du milieu) ; leurs primitives (colonne de droite) ne sont \textbf{pas} exigibles au Probatoire pour l'instant (sauf cas particulier $1/x \to \ln|x|$ en culture).
    \item \textbf{Domaine de définition :} Pour $\sqrt{x}$ et $\ln|x|$, l'intervalle $I$ doit être choisi soigneusement ($]0,+\infty[$ pour $\sqrt{x}$ et $\ln x$ ; $]-\infty,0[$ pour $\ln(-x)$). Au Probatoire, on précise souvent « sur $]0,+\infty[$ » pour lever l'ambiguïté.
    \item \textbf{La constante $C$ :} Elle n'apparaît \textbf{jamais} dans la colonne des dérivées (la dérivée d'une constante est $0$), mais elle est \textbf{obligatoire} dans la colonne des primitives pour désigner l'ensemble des solutions.
\end{itemize}
\end{attention}

\courssub{La notion de primitive : définition et propriété fondamentale}

Avant de manipuler les polynômes, posons les bases rigoureuses. En classe de Terminale technique, la définition doit être connue par cœur et la propriété fondamentale admise (sa démonstration dépasse le cadre du programme, mais son intuition est essentielle).

\begin{definition}[Primitive d'une fonction sur un intervalle]
Soit $f$ une fonction définie sur un intervalle $I$ de $\mathbb{R}$. On dit qu'une fonction $F$, dérivable sur $I$, est une \textbf{primitive} de $f$ sur $I$ si et seulement si :
\[
\forall x \in I, \quad F'(x) = f(x).
\]
On note souvent $F = \int f(x) \, dx$ (intégrale indéfinie) pour désigner l'ensemble des primitives.
\end{definition}

\begin{propriete}[Existence et forme générale des primitives (Admise)]
\begin{enumerate}
    \item \textbf{Théorème fondamental de l'analyse (admis) :} Toute fonction \textbf{continue} sur un intervalle $I$ admet au moins une primitive sur $I$.
    \item \textbf{Toutes les primitives :} Si $F_0$ est une primitive de $f$ sur $I$, alors l'ensemble de \textbf{toutes} les primitives de $f$ sur $I$ est exactement :
    \[
    \{ F_0 + C \mid C \in \mathbb{R} \}.
    \]
    On note cet ensemble $\int f(x) \, dx = F_0(x) + C$.
\end{enumerate}
\end{propriete}

\begin{savaistu}[Le symbole $\int$ : un « S » pour « Somme »]
Le symbole d'intégrale $\int$ a été introduit par \textbf{Gottfried Wilhelm Leibniz} (1646--1716). C'est un « S » majuscule allongé, initiale de \emph{Summa} (somme en latin). Leibniz concevait l'intégrale comme une somme infinie de quantités infinitésimales ($f(x)dx$). En Terminale technique, quand vous calculez une primitive, vous retrouvez l'esprit de Leibniz : « sommer » les variations élémentaires pour reconstruire la quantité totale (surface, volume, travail, charge électrique).
\end{savaistu}

\courssub{Primitives des fonctions polynomiales : méthode et exemples}

Les fonctions polynômes sont les « briques de base » de la modélisation technique (trajectoires balistiques, courbes de coûts, profils de poutres). Heureusement, leur primitive s'obtient \textbf{terme à terme} grâce à la linéarité.

\begin{methode}[Trouver une primitive d'un polynôme $P(x)$]
Pour déterminer une primitive de $P(x) = a_n x^n + a_{n-1} x^{n-1} + \dots + a_1 x + a_0$ sur $\mathbb{R}$ :
\begin{enumerate}
    \item \textbf{Écrire le polynôme sous forme de somme de puissances} (développer si nécessaire : $(x+1)^2 = x^2+2x+1$).
    \item \textbf{Traiter chaque terme $a_k x^k$ séparément :}
    \begin{itemize}
        \item Augmenter l'exposant de $1$ : $k \to k+1$.
        \item Diviser le coefficient par le \textbf{nouvel exposant} : $\dfrac{a_k}{k+1}$.
        \item Écrire le terme primitivé : $\dfrac{a_k}{k+1} x^{k+1}$.
    \end{itemize}
    \item \textbf{Sommer les primitives obtenues} (primitive d'une somme = somme des primitives).
    \item \textbf{Ajouter la constante $+C$} (indispensable pour désigner \emph{toutes} les primitives).
    \item \textbf{Vérifier en dérivant} le résultat trouvé : on doit retomber exactement sur $P(x)$.
\end{enumerate}
\end{methode}

\begin{exoresolu}[Primitive d'un trinôme du second degré]
\textbf{Énoncé :} Déterminer l'ensemble des primitives de la fonction $f$ définie sur $\mathbb{R}$ par $f(x) = 3x^2 - 4x + 5$.

\textbf{Solution détaillée :}
\begin{enumerate}
    \item $f$ est un polynôme, elle est continue sur $\mathbb{R}$, elle admet donc des primitives sur $\mathbb{R}$ (Théorème fondamental).
    \item On primitivise terme à terme en appliquant la règle « exposant $+1$, division par nouvel exposant » :
    \begin{align*}
        \text{Primitive de } 3x^2 &: \quad \frac{3}{2+1} x^{2+1} = \frac{3}{3}x^3 = x^3. \\
        \text{Primitive de } -4x = -4x^1 &: \quad \frac{-4}{1+1} x^{1+1} = \frac{-4}{2}x^2 = -2x^2. \\
        \text{Primitive de } 5 = 5x^0 &: \quad \frac{5}{0+1} x^{0+1} = 5x.
    \end{align*}
    \item On somme ces primitives et on ajoute $+C$ :
    \[
    F(x) = x^3 - 2x^2 + 5x + C, \quad C \in \mathbb{R}.
    \]
    \item \textbf{Vérification (obligatoire au Probatoire) :} On dérive $F$ :
    \[
    F'(x) = 3x^2 - 4x + 5 = f(x). \quad \c{c}a y est !
    \]
\end{enumerate}
\textbf{Conclusion :} L'ensemble des primitives de $f$ sur $\mathbb{R}$ est $\{ x^3 - 2x^2 + 5x + C \mid C \in \mathbb{R} \}$.
\end{exoresolu}

\begin{attention}[Les 3 erreurs fatales à ne pas commettre]
\begin{enumerate}
    \item \textbf{Oublier le $+C$ :} Sans $C$, vous ne donnez qu'\emph{une} primitive, pas \emph{l'ensemble} des primitives. Au Probatoire, l'oubli de $+C$ fait perdre la moitié des points (ou la totalité si la question porte sur « l'ensemble des primitives »).
    \item \textbf{Confondre dériver et primitiver :} Devant $3x^2$, certains écrivent $6x$ (dérivée) au lieu de $x^3$ (primitive). \textbf{Règle mnémotechnique :} Pour primitiver, l'exposant \textbf{monte} ($2 \to 3$) ; pour dériver, l'exposant \textbf{descend} ($2 \to 1$).
    \item \textbf{Oublier de diviser par le nouvel exposant :} Écrire $3x^3$ au lieu de $x^3$ pour la primitive de $3x^2$. La vérification par dérivation ($9x^2 \neq 3x^2$) permet de s'en apercevoir immédiatement.
\end{enumerate}
\end{attention}

\courssub{Interprétation graphique : une famille de courbes décalées verticalement}

La constante $C$ n'est pas un artifice algébrique : elle a une signification géométrique précise. Toutes les primitives d'une même fonction forment une \textbf{famille de courbes} qui se déduisent les unes des autres par une translation verticale. Cela signifie que la \emph{pente} (la dérivée) est identique en tout point d'abscisse $x$ pour toutes les courbes de la famille.

\begin{popfigure}
\begin{tikzpicture}[scale=0.8, domain=-1:2.2]
    % Axes
    \draw[->, thick, popdark] (-1.5,0) -- (2.7,0) node[right] {$x$};
    \draw[->, thick, popdark] (0,-6) -- (0,10) node[above] {$y$};
    % Graduations
    \foreach \x in {-1,1,2} \draw (\x,0.1) -- (\x,-0.1) node[below] {\small $\x$};
    \foreach \y in {-5,-2,2,4,6,8} \draw (0.1,\y) -- (-0.1,\y) node[left] {\small $\y$};
    
    % Fonction f(x) = 3x^2 - 4x + 5 (parabole, dérivée)
    \draw[poporange, very thick, samples=200, domain=-1:2.2] plot (\x, {3*\x*\x - 4*\x + 5});
    \node[poporange, right] at (2.2, {3*2.2*2.2 - 4*2.2 + 5}) {$f(x)=3x^2-4x+5$};
    
    % Primitive F(x) = x^3 - 2x^2 + 5x + C pour C = -2, 0, 2
    \draw[popblue, thick, samples=200, domain=-1:2.2] plot (\x, {\x*\x*\x - 2*\x*\x + 5*\x - 2});
    \node[popblue, right] at (2.2, {2.2*2.2*2.2 - 2*2.2*2.2 + 5*2.2 - 2}) {$F_{-2}(x)$};
    
    \draw[popblue, thick, dashed, samples=200, domain=-1:2.2] plot (\x, {\x*\x*\x - 2*\x*\x + 5*\x + 0});
    \node[popblue, right] at (2.2, {2.2*2.2*2.2 - 2*2.2*2.2 + 5*2.2 + 0}) {$F_{0}(x)$};
    
    \draw[popblue, thick, dotted, samples=200, domain=-1:2.2] plot (\x, {\x*\x*\x - 2*\x*\x + 5*\x + 2});
    \node[popblue, right] at (2.2, {2.2*2.2*2.2 - 2*2.2*2.2 + 5*2.2 + 2}) {$F_{2}(x)$};
    
    % Tangentes parallèles en x=1 pour illustrer même dérivée
    % f(1) = 4. F_C(1) = 4 + C.
    % Tangente en (1, 4+C) de pente 4 : y - (4+C) = 4(x-1) -> y = 4x + C.
    \draw[popgreen, thick, domain=-0.5:2] plot (\x, {4*\x + 0}); % Tangente pour C=0
    \node[popgreen, above left] at (2, 8) {Tangentes de pente $f(1)=4$};
    
    % Points de contact sur les F (x=1)
    \foreach \c/\style in {-2/solid, 0/dashed, 2/dotted} {
        \pgfmathsetmacro{\yc}{\c + 4} % F(1) = 4 + C
        \fill[popblue] (1,\yc) circle (2.5pt);
    }
    % Ligne verticale pointillée x=1
    \draw[popmuted, thin, dashed] (1,-6) -- (1,10);
\end{tikzpicture}
\legende{Famille des primitives $F_C(x)=x^3-2x^2+5x+C$ de $f(x)=3x^2-4x+5$. En $x=1$, les trois courbes ont la même pente $f(1)=4$ (tangentes parallèles vertes). La courbe orange représente la dérivée $f$.}
\end{popfigure}

\begin{aretenir}[Bilan : Dérivée $\leftrightarrow$ Primitive]
\begin{itemize}
    \item Le tableau des dérivées usuelles se lit dans les deux sens (attention : primitives exigibles seulement pour puissances entières $n \in \mathbb{N}$).
    \item Une primitive de $f$ sur $I$ est une fonction $F$ telle que $F'=f$ sur $I$.
    \item Si $F$ est une primitive, alors $F+C$ ($C\in\mathbb{R}$) sont \textbf{toutes} les primitives.
    \item Pour un polynôme : on primitivise terme à terme (exposant $+1$, division par nouvel exposant) et on ajoute $+C$.
    \item \textbf{Toujours vérifier en dérivant.}
\end{itemize}
\end{aretenir}

La section suivante, « Primitives des fonctions polynomiales », approfondira les cas particuliers (produits, quotients, développements nécessaires) et l'application aux problèmes concrets : calcul d'aires, de volumes, de travaux mécaniques et de charges électriques.

\coursec{Rappels sur la dérivée et interprétation géométrique}

\courssub{Nombre dérivé et tangente}
Soit $f$ une fonction définie sur un intervalle $I$ et $a \in I$. Si la limite
\[ \lim_{x \to a} \frac{f(x) - f(a)}{x - a} \]
existe et est finie, on dit que $f$ est \cle{dérivable} en $a$. Cette limite est appelée \cle{nombre dérivé} de $f$ en $a$ et se note $f'(a)$.

\begin{propriete}[Interprétation géométrique]
Le nombre dérivé $f'(a)$ est le \textbf{coefficient directeur} de la tangente à la courbe $\mathcal{C}_f$ au point $A(a\,;\,f(a))$. L'équation de cette tangente est :
\[ y = f'(a)(x - a) + f(a). \]
\end{propriete}

\begin{attention}[Conditions de dérivabilité]
Une fonction n'est pas dérivable en $a$ si :
\begin{itemize}
\item elle n'est pas continue en $a$ ;
\item sa courbe présente un \textbf{pic} (demi-tangentes de pentes opposées, ex. $x \mapsto |x|$ en $0$) ;
\item sa courbe a une \textbf{tangente verticale} (pente infinie, ex. $x \mapsto \sqrt[3]{x}$ en $0$).
\end{itemize}
\end{attention}

\courssub{Fonction dérivée}
Si $f$ est dérivable en tout point de $I$, on définit la \cle{fonction dérivée} $f' : I \to \mathbb{R}$ qui associe à $x$ le nombre dérivé $f'(x)$. Les notations usuelles sont $f'(x)$, $y'$, $\dfrac{\mathrm{d}y}{\mathrm{d}x}$.

\begin{exemplebox}[Calcul d'un nombre dérivé et équation de tangente]
Soit $f(x) = x^2 - 2x + 3$. Calculer $f'(1)$ et l'équation de la tangente en $x=1$.
\begin{align*}
f'(x) &= 2x - 2 \quad \text{(règle de la puissance, vue ci-après)} \\
f'(1) &= 0. \\
\text{Tangente en } (1,\,2) :\quad y &= 0 \cdot (x - 1) + 2 = 2.
\end{align*}
La tangente est horizontale (extremum).
\end{exemplebox}

\coursec{Règles de dérivation usuelles}

\begin{propriete}[Linéarité]
Soient $u$ et $v$ dérivables sur $I$, $\lambda, \mu \in \mathbb{R}$.
\[ (\lambda u + \mu v)' = \lambda u' + \mu v'. \]
\end{propriete}

\begin{propriete}[Produit et quotient]
Soient $u, v$ dérivables sur $I$.
\begin{itemize}
\item \textbf{Produit :} $(uv)' = u'v + uv'$.
\item \textbf{Quotient :} $\left(\dfrac{u}{v}\right)' = \dfrac{u'v - uv'}{v^2}$, pour $v(x) \neq 0$.
\end{itemize}
\end{propriete}

\begin{propriete}[Fonction composée (règle de la chaîne)]
Soient $u$ dérivable sur $J$ et $v$ dérivable sur $I$ avec $v(I) \subset J$. Alors $u \circ v$ est dérivable sur $I$ et :
\[ (u \circ v)'(x) = v'(x) \cdot u'(v(x)). \]
\end{propriete}

\begin{methode}[Dériver une fonction composée $u(v(x))$]
\begin{enumerate}
\item Identifier la fonction \textbf{extérieure} $u$ et la fonction \textbf{intérieure} $v$.
\item Calculer $v'(x)$.
\item Calculer $u'(X)$ (variable $X$), puis remplacer $X$ par $v(x)$ pour obtenir $u'(v(x))$.
\item Multiplier : $v'(x) \times u'(v(x))$.
\end{enumerate}
\end{methode}

\begin{attention}[Erreurs fréquentes sur la chaîne]
\begin{itemize}
\item Oublier de dériver l'intérieur $v(x)$ (facteur $v'(x)$ manquant).
\item Dériver l'extérieur en gardant $x$ au lieu de $v(x)$.
\item Confondre $(u \circ v)'$ et $u' \circ v'$.
\end{itemize}
\end{attention}

\coursec{Tableau des dérivées usuelles et premières primitives}

\begin{propriete}[Tableau de référence (à connaître par cœur)]
\begin{center}
\begin{tabularx}{\linewidth}{|l|X|X|X|}\hline
\rowcolor{popblueL}
\textbf{Fonction $f(x)$} & \textbf{Domaine $D$} & \textbf{Dérivée $f'(x)$} & \textbf{Une primitive $F(x)$ sur $D$} \\ \hline
$k$ (constante) & $\mathbb{R}$ & $0$ & $kx$ \\ \hline
$x^n$ ($n \in \mathbb{N}$) & $\mathbb{R}$ & $n x^{n-1}$ & $\dfrac{x^{n+1}}{n+1}$ \\ \hline
$x^\alpha$ ($\alpha \in \mathbb{R}$) & $]0,+\infty[$ & $\alpha x^{\alpha-1}$ & $\dfrac{x^{\alpha+1}}{\alpha+1}$ ($\alpha \neq -1$) \\ \hline
$\dfrac{1}{x}$ & $\mathbb{R}^*$ & $-\dfrac{1}{x^2}$ & $\ln|x|$ \\ \hline
$\sqrt{x}$ & $[0,+\infty[$ & $\dfrac{1}{2\sqrt{x}}$ & $\dfrac{2}{3}x^{3/2}$ \\ \hline
$\ln|x|$ & $\mathbb{R}^*$ & $\dfrac{1}{x}$ & $x\ln|x| - x$ \\ \hline
$\mathrm{e}^x$ & $\mathbb{R}$ & $\mathrm{e}^x$ & $\mathrm{e}^x$ \\ \hline
$\sin x$ & $\mathbb{R}$ & $\cos x$ & $-\cos x$ \\ \hline
$\cos x$ & $\mathbb{R}$ & $-\sin x$ & $\sin x$ \\ \hline
\end{tabularx}
\end{center}
\end{propriete}

\begin{aretenir}[Opérations sur les primitives]
Si $F$ est une primitive de $f$ et $G$ une primitive de $g$ sur un intervalle $I$, alors :
\begin{itemize}
\item $F + G$ est primitive de $f + g$ ;
\item $\lambda F$ ($\lambda \in \mathbb{R}$) est primitive de $\lambda f$ ;
\item $\dfrac{1}{a}F(ax+b)$ ($a \neq 0$) est primitive de $f(ax+b)$.
\end{itemize}
\end{aretenir}

\coursec{Primitives des fonctions polynomiales}

Dans la section précédente, nous avons dressé le tableau des dérivées usuelles et nous avons commencé à « remonter » le calcul : connaissant une fonction $f$, chercher une fonction $F$ dont la dérivée est $f$. Cette démarche, appelée \cle{recherche de primitives}, est l'opération inverse de la dérivation. Nous allons maintenant la systématiser pour les \cle{fonctions polynomiales}, qui sont les fonctions les plus fréquentes dans les applications techniques : coût de production, distance parcourue, aire d'une tôle, volume d'une pièce, etc.

\courssub{Existence des primitives d'un polynôme}

Commençons par une question naturelle : toute fonction admet-elle une primitive ? La réponse générale est délicate, mais pour les polynômes, elle est simple et rassurante.

\begin{propriete}[Existence des primitives d'un polynôme (admise)]
Toute fonction \cle{polynomiale} est \textbf{continue} sur $\mathbb{R}$. Or toute fonction continue sur un intervalle admet des primitives sur cet intervalle. Donc \textbf{tout polynôme admet des primitives sur $\mathbb{R}$}.
\end{propriete}

Ce résultat est admis au programme de Terminale technique : il nous autorise à chercher des primitives de polynômes sans jamais nous poser la question de leur existence. De plus, on sait déjà que si $F$ est une primitive de $f$, alors toutes les primitives de $f$ sont les fonctions $F + C$, où $C$ est une constante réelle quelconque. On parle de la \cle{primitive générale}.

\begin{attention}[Ne pas oublier la constante $C$]
Deux primitives d'une même fonction diffèrent d'une constante. Quand on demande « une primitive » ou « la primitive générale », la réponse doit contenir « $+\,C$ ». Oublier $C$ est une erreur classique au Probatoire et au Baccalauréat technique. En revanche, si l'on donne une condition supplémentaire (par exemple $F(0) = 5$), on peut déterminer $C$ et obtenir \textbf{une seule} primitive.
\end{attention}

\courssub{Primitive d'un monôme : la règle fondamentale}

Tout repose sur une observation simple. Nous savons dériver une puissance :
\[ \left(x^{n}\right)' = n\,x^{n-1}. \]
Pour « remonter », il faut faire l'inverse : augmenter l'exposant de $1$, puis compenser le facteur qui apparaît en dérivant. En effet :
\[ \left(\frac{x^{n+1}}{n+1}\right)' = \frac{(n+1)x^{n}}{n+1} = x^{n}. \]

\begin{experience}[Découverte de la règle de la primitive d'un monôme]
On cherche une primitive de $x^2$. On teste des candidats de la forme $k x^3$.
\begin{itemize}
\item $(x^3)' = 3x^2$ $\to$ trop grand (facteur 3).
\item $\left(\frac{x^3}{2}\right)' = \frac{3x^2}{2}$ $\to$ toujours pas $x^2$.
\item $\left(\frac{x^3}{3}\right)' = \frac{3x^2}{3} = x^2$ $\to$ \textbf{Correct !}
\end{itemize}
On a divisé par le \textbf{nouvel} exposant (3). Généralisons : pour $x^n$, on essaie $\frac{x^{n+1}}{n+1}$.
\end{experience}

\begin{propriete}[Primitive d'un monôme]
Pour tout entier naturel $n$, une primitive de la fonction $x \mapsto x^{n}$ sur $\mathbb{R}$ est :
\[ x \mapsto \frac{x^{n+1}}{n+1}. \]
En particulier :
\begin{itemize}
\item une primitive de $1$ (c'est-à-dire $x^0$) est $x$ ;
\item une primitive de $x$ est $\dfrac{x^{2}}{2}$ ;
\item une primitive de $x^{2}$ est $\dfrac{x^{3}}{3}$ ;
\item une primitive de $x^{3}$ est $\dfrac{x^{4}}{4}$.
\end{itemize}
\end{propriete}

La démonstration est immédiate et formatrice : il suffit de dériver le résultat proposé et de vérifier qu'on retombe sur le monôme de départ. C'est exactement ce que nous avons fait ci-dessus.

\begin{attention}[L'erreur fréquente : diviser par l'ancien exposant]
Beaucoup d'élèves écrivent $\displaystyle\int x^{2}\,\mathrm{d}x = \frac{x^{3}}{2} + C$. C'est \textbf{faux} : on divise par le \textbf{nouvel} exposant, pas par l'ancien. La bonne réponse est :
\[ \int x^{2}\,\mathrm{d}x = \frac{x^{3}}{3} + C. \]
Vérification : $\left(\dfrac{x^{3}}{3}\right)' = \dfrac{3x^{2}}{3} = x^{2}$. En revanche $\left(\dfrac{x^{3}}{2}\right)' = \dfrac{3x^{2}}{2} \neq x^{2}$. \textbf{Réflexe de survie} : après avoir trouvé une primitive, dérivez-la mentalement pour contrôler.
\end{attention}

\courssub{Primitive générale d'un polynôme et notation de l'intégrale indéfinie}

La recherche de primitives possède deux propriétés de linéarité, héritées de la dérivation :
\begin{itemize}
\item une primitive d'une \textbf{somme} est la somme des primitives ;
\item une primitive de $k \times f$ (où $k$ est une constante) est $k$ fois une primitive de $f$.
\end{itemize}

Un polynôme étant une somme de monômes, on peut donc le primitiver \textbf{terme à terme}.

\begin{propriete}[Primitive générale d'un polynôme]
Soit $P(x) = a_{0} + a_{1}x + a_{2}x^{2} + \cdots + a_{n}x^{n} = \displaystyle\sum_{k=0}^{n} a_{k}\,x^{k}$ un polynôme. Alors la primitive générale de $P$ sur $\mathbb{R}$ est :
\[ \int P(x)\,\mathrm{d}x = a_{0}x + a_{1}\frac{x^{2}}{2} + a_{2}\frac{x^{3}}{3} + \cdots + a_{n}\frac{x^{n+1}}{n+1} + C = \sum_{k=0}^{n} a_{k}\,\frac{x^{k+1}}{k+1} + C, \]
où $C$ est une constante réelle arbitraire.
\end{propriete}

\begin{definition}[Intégrale indéfinie]
On appelle \cle{intégrale indéfinie} de $f$, et on note $\displaystyle\int f(x)\,\mathrm{d}x$, l'ensemble de toutes les primitives de $f$. Le symbole $\int$ se lit « intégrale » (ou « somme »), et $\mathrm{d}x$ rappelle que la variable est $x$. Ainsi :
\[ \int P(x)\,\mathrm{d}x = \sum_{k=0}^{n} a_{k}\,\frac{x^{k+1}}{k+1} + C. \]
\end{definition}

\begin{methode}[Intégrer un polynôme en trois étapes]
Pour trouver la primitive générale d'un polynôme :
\begin{enumerate}
\item \textbf{Traiter chaque terme séparément}, en gardant son coefficient.
\item Pour chaque monôme $a\,x^{k}$ : \textbf{augmenter l'exposant de 1} ($x^{k}$ devient $x^{k+1}$), puis \textbf{diviser par le nouvel exposant} : on obtient $a\,\dfrac{x^{k+1}}{k+1}$.
\item \textbf{Ajouter la constante} $+C$ à la fin, puis \textbf{vérifier} en dérivant le résultat : on doit retrouver le polynôme de départ.
\end{enumerate}
\end{methode}

\begin{exemplebox}[Exemple chiffré détaillé]
Calculons $\displaystyle\int \left(4x^{3} - 2x + 1\right)\mathrm{d}x$.
\begin{align*}
\int \left(4x^{3} - 2x + 1\right)\mathrm{d}x &= 4\times\frac{x^{4}}{4} - 2\times\frac{x^{2}}{2} + 1\times x + C \\
&= x^{4} - x^{2} + x + C.
\end{align*}
\textbf{Vérification} : $\left(x^{4} - x^{2} + x + C\right)' = 4x^{3} - 2x + 1$. C'est bien le polynôme de départ.
\end{exemplebox}

\textbf{Cas particulier : la fonction nulle.} Quelle est la primitive de la fonction nulle $f(x) = 0$ ? On cherche les fonctions dont la dérivée est nulle sur $\mathbb{R}$ : ce sont exactement les \textbf{fonctions constantes}. Donc :
\[ \int 0\,\mathrm{d}x = C, \quad C \in \mathbb{R}. \]
Ce cas, en apparence trivial, est très utile : il justifie que deux primitives d'une même fonction diffèrent d'une constante (leur différence a une dérivée nulle, donc est constante).

\courssub{De la primitive à l'aire : l'intégrale définie}

Pourquoi les techniciens et les ingénieurs s'intéressent-ils tant aux primitives ? Parce qu'elles permettent de calculer des \cle{aires}, et donc des quantités concrètes : surface de matériau, distance parcourue, énergie consommée, quantité produite.

\textbf{Intuition.} Considérons la fonction $f(x) = 2x + 1$ sur $[0\,;3]$. L'aire comprise entre la courbe, l'axe des abscisses et les droites $x = 0$ et $x = 3$ peut être approchée par une somme d'aires de rectangles (on parle de \emph{rectangles de Riemann}) : plus les rectangles sont fins, meilleure est l'approximation. À la limite, on obtient l'aire exacte, appelée \cle{intégrale définie} de $f$ entre $0$ et $3$, notée $\displaystyle\int_{0}^{3} f(x)\,\mathrm{d}x$.

\begin{popfigure}
\begin{tikzpicture}
\begin{axis}[
  width=0.9\linewidth, height=7cm,
  axis lines=middle,
  xlabel={$x$}, ylabel={$y$},
  xmin=-0.3, xmax=3.5, ymin=-0.5, ymax=8,
  xtick={0,0.5,1,1.5,2,2.5,3},
  ytick={1,3,5,7},
  tick label style={font=\small},
  label style={font=\small},
  clip=false
]
% Rectangles de Riemann (sommes à gauche), pas de 0.5
\foreach \i in {0,0.5,1,1.5,2,2.5}{
  \addplot[draw=popblue, fill=popblue!30, opacity=0.7] coordinates {(\i,0) (\i,{2*\i+1}) (\i+0.5,{2*\i+1}) (\i+0.5,0)} \closedcycle;
}
% Aire exacte sous la courbe
\addplot[draw=none, fill=poporange!40, opacity=0.4, domain=0:3, samples=50] {2*x+1} \closedcycle;
% Courbe
\addplot[popcurve, domain=-0.2:3.3, samples=100, thick] {2*x+1};
\node[font=\small, popdark] at (axis cs:2.1,6.6) {$f(x)=2x+1$};
\end{axis}
\end{tikzpicture}
\legende{Aire sous la courbe de $f(x)=2x+1$ entre $x=0$ et $x=3$, approchée par des rectangles de Riemann (en bleu). En raffinant le découpage, la somme des aires des rectangles tend vers l'aire exacte (en orange).}
\end{popfigure}

\begin{propriete}[Théorème fondamental de l'analyse (admis)]
Soit $f$ une fonction continue et positive sur $[a\,;b]$, et $F$ une primitive de $f$ sur $[a\,;b]$. Alors l'aire sous la courbe de $f$ entre $a$ et $b$, exprimée en unités d'aire, vaut :
\[ \int_{a}^{b} f(x)\,\mathrm{d}x = F(b) - F(a). \]
On note souvent $F(b) - F(a) = \left[F(x)\right]_{a}^{b}$. Ce résultat, admis au programme, est le pont entre les primitives et les aires.
\end{propriete}

\begin{exemplebox}[Calcul d'aire avec une primitive]
Calculons l'aire sous la courbe de $f(x) = 2x + 1$ entre $0$ et $3$. Une primitive de $f$ est $F(x) = x^{2} + x$. Donc :
\[ \int_{0}^{3} (2x+1)\,\mathrm{d}x = \left[x^{2} + x\right]_{0}^{3} = (9 + 3) - (0 + 0) = 12. \]
L'aire vaut $12$ unités d'aire. \textbf{Contrôle géométrique} : la région est un trapèze de hauteurs $f(0)=1$ et $f(3)=7$, de largeur $3$ : $\mathcal{A} = \dfrac{(1+7)\times 3}{2} = 12$. Les deux méthodes concordent.
\end{exemplebox}

\begin{exoresolu}[Primitive déterminée par une condition]
\textbf{Énoncé.} Déterminer la primitive $F$ de $f(x) = 6x^{2} - 4x$ sur $\mathbb{R}$ telle que $F(1) = 3$.
\tcblower
\textbf{Solution détaillée.}
\textbf{Étape 1 : primitive générale.} On primitive terme à terme :
\[ F(x) = 6\times\frac{x^{3}}{3} - 4\times\frac{x^{2}}{2} + C = 2x^{3} - 2x^{2} + C. \]
\textbf{Étape 2 : utilisation de la condition.} $F(1) = 2 - 2 + C = C$. Or $F(1) = 3$, donc $C = 3$.
\textbf{Étape 3 : conclusion.} $F(x) = 2x^{3} - 2x^{2} + 3$.
\textbf{Vérification} : $F'(x) = 6x^{2} - 4x = f(x)$ et $F(1) = 2 - 2 + 3 = 3$. Correct.
\end{exoresolu}

\begin{exoresolu}[Vérification par dérivation]
\textbf{Énoncé.} Un élève affirme que $\displaystyle\int \left(3x^{2} + 2x - 5\right)\mathrm{d}x = x^{3} + x^{2} - 5x + C$. A-t-il raison ?
\tcblower
\textbf{Solution détaillée.} Pour vérifier une primitive, on la dérive :
\[ \left(x^{3} + x^{2} - 5x + C\right)' = 3x^{2} + 2x - 5. \]
On retombe exactement sur le polynôme intégré : l'élève a \textbf{raison}. Chaque terme respecte la règle : $3x^{2} \to 3\times\dfrac{x^{3}}{3} = x^{3}$ ; $2x \to 2\times\dfrac{x^{2}}{2} = x^{2}$ ; $-5 \to -5x$.
\end{exoresolu}

\coursec{Résolution de problèmes concrets : modélisation et interprétation}

\courssub{Application : reconstituer un coût total à partir du coût marginal}

En gestion de production (comptabilité, atelier, exploitation agricole), on appelle \cle{coût marginal} le coût supplémentaire engendré par la production d'une unité de plus. On le modélise par la dérivée $C'(x)$ du coût total $C(x)$, où $x$ est la quantité produite. Reconstituer le coût total à partir du coût marginal est donc un problème de \textbf{primitive}.

\begin{exoresolu}[Du coût marginal au coût total]
\textbf{Énoncé.} Un atelier de menuiserie à Douala produit $x$ chaises par semaine. Le coût marginal, en milliers de francs CFA, est donné par $C'(x) = 0{,}6x^{2} - 4x + 15$. Les coûts fixes (loyer, machines) s'élèvent à $50$ milliers de francs, c'est-à-dire $C(0) = 50$. Déterminer le coût total $C(x)$, puis le coût de production de $10$ chaises.
\tcblower
\textbf{Solution détaillée.}
\textbf{Étape 1 : primitive générale du coût marginal.}
\[ C(x) = 0{,}6\times\frac{x^{3}}{3} - 4\times\frac{x^{2}}{2} + 15x + K = 0{,}2x^{3} - 2x^{2} + 15x + K. \]
\textbf{Étape 2 : détermination de la constante.} $C(0) = K = 50$ (les coûts fixes). Donc :
\[ C(x) = 0{,}2x^{3} - 2x^{2} + 15x + 50. \]
\textbf{Étape 3 : application numérique.} Pour $x = 10$ :
\[ C(10) = 0{,}2\times 1000 - 2\times 100 + 15\times 10 + 50 = 200 - 200 + 150 + 50 = 200. \]
La production de $10$ chaises coûte $200$ milliers de francs CFA, soit $200\,000$ FCFA.
\textbf{Vérification} : $C'(x) = 0{,}6x^{2} - 4x + 15$, on retrouve bien le coût marginal.
\end{exoresolu}

\courssub{Application : de la vitesse à la position (cinématique)}

En physique, la \textbf{vitesse} est la dérivée de la \textbf{position} par rapport au temps : $v(t) = x'(t)$. Par conséquent, la primitive de la vitesse redonne la position, à une constante près — cette constante étant la position initiale.

\begin{exoresolu}[Position à partir de la vitesse]
\textbf{Énoncé.} Un véhicule part de l'origine ($x(0)=0$) avec une vitesse $v(t) = 3t^{2} - 2t$ (en m/s, $t$ en secondes). Déterminer sa position $x(t)$ et sa position à $t=3$ s.
\tcblower
\textbf{Solution.}
Primitive générale : $x(t) = 3\frac{t^{3}}{3} - 2\frac{t^{2}}{2} + C = t^{3} - t^{2} + C$.
Condition initiale : $x(0) = C = 0$. Donc $x(t) = t^{3} - t^{2}$.
À $t=3$ : $x(3) = 27 - 9 = 18$ m.
\end{exoresolu}

\begin{savaistu}[Primitives et physique : de la vitesse à la position]
C'est exactement le même mécanisme que pour le coût marginal et le coût total : les primitives permettent de « remonter » d'une variation (vitesse, coût marginal, accélération) à la quantité elle-même (position, coût total, vitesse). C'est l'un des outils fondateurs de la mécanique depuis Newton et Leibniz, inventeurs du calcul intégral au XVII\textsuperscript{e} siècle. Au Cameroun, cet outil est utilisé en génie civil (calcul de volumes de terrassement), en électricité (charge d'un condensateur $Q = \int i\,dt$), en gestion des stocks, etc.
\end{savaistu}

\begin{exercice}[À vous de jouer]
\begin{enumerate}
\item Déterminer la primitive générale de $P(x) = 5x^{4} - 3x^{2} + 7$, puis vérifier par dérivation.
\item Déterminer la primitive $F$ de $f(x) = 4x - 1$ telle que $F(2) = 10$.
\item Calculer l'aire sous la courbe de $f(x) = 3x^{2}$ entre $x = 1$ et $x = 2$.
\item Un coût marginal est donné par $C'(x) = 2x + 10$. Les coûts fixes sont de $100$. Trouver $C(x)$ et $C(5)$.
\end{enumerate}
\end{exercice}

\begin{corrige}[Exercices]
\begin{enumerate}
\item $F(x) = 5\times\dfrac{x^{5}}{5} - 3\times\dfrac{x^{3}}{3} + 7x + C = x^{5} - x^{3} + 7x + C$. Vérification : $F'(x) = 5x^{4} - 3x^{2} + 7 = P(x)$.
\item Primitive générale : $F(x) = 2x^{2} - x + C$. Condition : $F(2) = 8 - 2 + C = 6 + C = 10$, donc $C = 4$. Ainsi $F(x) = 2x^{2} - x + 4$.
\item Une primitive de $3x^{2}$ est $x^{3}$. Donc $\displaystyle\int_{1}^{2} 3x^{2}\,\mathrm{d}x = \left[x^{3}\right]_{1}^{2} = 8 - 1 = 7$ unités d'aire.
\item $C(x) = x^{2} + 10x + K$. $C(0) = K = 100$. $C(x) = x^{2} + 10x + 100$. $C(5) = 25 + 50 + 100 = 175$.
\end{enumerate}
\end{corrige}

\begin{aretenir}[L'essentiel sur les primitives de polynômes]
\begin{itemize}
\item Tout polynôme, étant continu sur $\mathbb{R}$, admet des primitives sur $\mathbb{R}$.
\item Règle du monôme : $\displaystyle\int x^{n}\,\mathrm{d}x = \frac{x^{n+1}}{n+1} + C$ — on \textbf{augmente l'exposant} et on \textbf{divise par le nouvel exposant}.
\item On primitive un polynôme \textbf{terme à terme} : $\displaystyle\int \sum_{k=0}^{n} a_{k}x^{k}\,\mathrm{d}x = \sum_{k=0}^{n} a_{k}\frac{x^{k+1}}{k+1} + C$.
\item La primitive de la fonction nulle est une constante.
\item Aire sous la courbe (théorème admis) : $\displaystyle\int_{a}^{b} f(x)\,\mathrm{d}x = F(b) - F(a)$.
\item \textbf{Réflexe de contrôle} : dériver la primitive trouvée doit redonner la fonction de départ.
\end{itemize}
\end{aretenir}

\coursec{Résolution de problèmes concrets : modélisation et interprétation}

Après avoir maîtrisé le calcul des dérivées et des primitives des fonctions polynomiales, nous abordons maintenant l'objectif final de cette leçon : \textbf{utiliser ces outils pour résoudre des problèmes réels}. En entreprise, en gestion de production ou en économie, les quantités évoluent souvent selon des lois polynomiales. La dérivée nous donne le \cle{taux de variation instantané} (coût marginal, productivité marginale) et l'annulation de la dérivée permet de trouver les \cle{optimums} (bénéfice maximal, coût minimal). La primitive, quant à elle, permet de \cle{reconstituer une quantité totale} à partir de sa variation (coût total à partir du coût marginal). Cette section met en œuvre la démarche complète exigée au Probatoire : modélisation, calcul, interprétation graphique et conclusion rédigée.

\courssub{Modélisation économique : coût total, revenu et bénéfice}

Dans une unité de production (usine, atelier, exploitation agricole), on modélise souvent les grandeurs économiques par des fonctions polynomiales de la variable $x$ (quantité produite ou vendue, en unités ou en tonnes).

\begin{definition}[Fonctions économiques usuelles]
Soit $x \in [0, x_{\max}]$ la quantité produite.
\begin{itemize}
    \item \textbf{Coût total $C_T(x)$} : somme des coûts fixes ($C_F$, indépendants de $x$) et des coûts variables ($C_V(x)$). Souvent modélisé par un polynôme de degré 2 ou 3 : $C_T(x) = ax^3 + bx^2 + cx + C_F$.
    \item \textbf{Revenu total $R_T(x)$} : prix de vente unitaire $p(x)$ multiplié par la quantité. Si le prix est constant ($p$), $R_T(x) = px$ (fonction affine). Si le prix baisse avec la quantité (loi de la demande), $R_T(x)$ est un polynôme de degré 2 (exemple : $R_T(x) = -dx^2 + px$).
    \item \textbf{Bénéfice $B(x)$} : différence entre le revenu et le coût total.
    \[ B(x) = R_T(x) - C_T(x) \]
    C'est une fonction polynôme dont on cherche souvent le \textbf{maximum}.
\end{itemize}
\end{definition}

\begin{exemplebox}[Modélisation d'une unité de transformation de manioc]
Une coopérative transforme du manioc en farine. Les études de marché et de coûts donnent :
\begin{itemize}
    \item Coûts fixes (location, assurance) : $C_F = 100$ (en milliers de FCFA).
    \item Coût variable : $C_V(x) = 2x^2 + 20x$ (main d'œuvre, énergie, matière première).
    \item Prix de vente unitaire en fonction de la demande : $p(x) = 80 - 2x$ (en milliers de FCFA par tonne).
\end{itemize}
\emph{Modélisation :}
\[ C_T(x) = C_V(x) + C_F = 2x^2 + 20x + 100 \]
\[ R_T(x) = x \cdot p(x) = x(80 - 2x) = -2x^2 + 80x \]
\[ B(x) = R_T(x) - C_T(x) = (-2x^2 + 80x) - (2x^2 + 20x + 100) = -4x^2 + 60x - 100 \]
Le domaine réaliste est $x \in [0, 40]$ (au-delà de $x = 40$, le prix $p(x) = 80 - 2x$ devient négatif, ce qui n'a pas de sens économique).
\end{exemplebox}

\courssub{Coût marginal : dérivée du coût total et interprétation}

Le \cle{coût marginal} $C_m(x)$ est le coût engendré par la production d'une unité supplémentaire au voisinage de $x$. Mathématiquement, c'est la dérivée du coût total.

\begin{propriete}[Coût marginal]
Si $C_T$ est dérivable sur $]0, x_{\max}[$, le coût marginal en $x$ est :
\[ C_m(x) = C_T'(x) \]
\emph{Interprétation concrète :} $C_m(x_0) \approx C_T(x_0+1) - C_T(x_0)$. C'est le coût approximatif de la $(x_0+1)$-ième unité produite.
\end{propriete}

\begin{attention}[Erreur fréquente : coût marginal et coût moyen]
Ne pas confondre :
\begin{itemize}
    \item \textbf{Coût marginal} $C_m(x) = C_T'(x)$ : variation instantanée (dérivée).
    \item \textbf{Coût moyen} $C_{moy}(x) = \dfrac{C_T(x)}{x}$ : coût total divisé par la quantité produite.
\end{itemize}
\emph{Exemple :} si $C_T(x) = x^2 + 10x + 100$, alors $C_m(x) = 2x + 10$ donc $C_m(10) = 30$, tandis que $C_{moy}(10) = \dfrac{100 + 100 + 100}{10} = 30$. Ils coïncident ici par hasard, mais ce sont des concepts distincts. Le coût marginal sert à décider de la production de l'unité \emph{supplémentaire} ; le coût moyen sert à fixer le prix de vente pour couvrir l'ensemble des charges.
\end{attention}

\courssub{Recherche de l'optimum : maximisation du bénéfice}

L'objectif de l'entreprise est souvent de maximiser le bénéfice $B(x)$. Comme $B$ est une fonction polynôme (donc dérivable) sur un intervalle fermé $[0, x_{\max}]$, le maximum, s'il est atteint à l'intérieur de l'intervalle, se trouve en un point où la dérivée s'annule en changeant de signe.

\begin{methode}[Recherche du bénéfice maximal]
\begin{enumerate}
    \item \textbf{Écrire} la fonction bénéfice $B(x) = R_T(x) - C_T(x)$ et préciser son domaine $D$ (souvent $[0, x_{\max}]$).
    \item \textbf{Dériver} : calculer $B'(x)$.
    \item \textbf{Résoudre} $B'(x) = 0$ dans $D$. On obtient les \cle{valeurs critiques} $x_1, x_2, \dots$
    \item \textbf{Étudier le signe de $B'(x)$} à l'aide d'un tableau de variations :
        \begin{itemize}
            \item si $B'(x)$ passe de $+$ à $-$ en $x_1$ $\rightarrow$ \textbf{maximum local} en $x_1$ ;
            \item si $B'(x)$ passe de $-$ à $+$ en $x_1$ $\rightarrow$ minimum local en $x_1$.
        \end{itemize}
    \item \textbf{Comparer} les valeurs de $B$ aux extrema locaux et aux bornes du domaine ($0$ et $x_{\max}$).
    \item \textbf{Conclure} en langage courant : quantité optimale, bénéfice maximal, vérification de la faisabilité (stock, capacité de production).
\end{enumerate}
\end{methode}

\begin{exoresolu}[Maximisation du bénéfice : exemple type Probatoire]
\textbf{Énoncé :} On considère la fonction bénéfice (en milliers de FCFA) :
\[ B(x) = -2x^2 + 40x - 100 \quad \text{pour } x \in [0, 25] \]
où $x$ est la quantité produite en tonnes. Déterminer la production qui maximise le bénéfice et la valeur de ce maximum.
\tcblower
\textbf{Solution détaillée :}
\begin{enumerate}
    \item \textbf{Domaine :} $D = [0, 25]$. $B$ est une fonction polynôme, donc dérivable sur $\mathbb{R}$, en particulier sur $[0, 25]$.
    \item \textbf{Dérivée :} $B'(x) = -4x + 40$.
    \item \textbf{Annulation :} $B'(x) = 0 \Leftrightarrow -4x + 40 = 0 \Leftrightarrow x = 10$.
        Or $10 \in [0, 25]$ : la valeur critique est bien dans le domaine.
    \item \textbf{Signe de $B'(x)$ et tableau de variations :} $B'(x)$ est une fonction affine de coefficient directeur $-4 < 0$, donc $B'(x) > 0$ pour $x < 10$ et $B'(x) < 0$ pour $x > 10$.
    \begin{center}
    \begin{tabular}{|c|ccccc|}\hline
    $x$ & $0$ & & $10$ & & $25$ \\ \hline
    $B'(x)$ & & $+$ & $0$ & $-$ & \\ \hline
    $B(x)$ & $-100$ & $\nearrow$ & $100$ & $\searrow$ & $-350$ \\ \hline
    \end{tabular}
    \end{center}
    \emph{Vérification des valeurs aux bornes :} $B(0) = -100$ et $B(25) = -2(625) + 40(25) - 100 = -1250 + 1000 - 100 = -350$.
    \item \textbf{Valeur du maximum :} $B(10) = -2(10)^2 + 40(10) - 100 = -200 + 400 - 100 = 100$.
    \item \textbf{Conclusion :} l'entreprise doit produire \textbf{10 tonnes} pour obtenir un bénéfice maximal de \textbf{100 milliers de FCFA}, soit \textbf{100 000 FCFA}.
\end{enumerate}
\end{exoresolu}

\courssub{Utilisation des primitives : retrouver le coût total à partir du coût marginal}

Réciproquement, si l'on connaît le coût marginal $C_m(x)$ (souvent issu de l'analyse des coûts variables unitaires) et les coûts fixes $C_F$, on retrouve le coût total par recherche d'une \textbf{primitive}.

\begin{propriete}[Coût total et primitive du coût marginal]
Si $C_m(x) = C_T'(x)$ sur un intervalle, alors le coût total est une primitive de $C_m$ :
\[ C_T(x) = P(x) + K \]
où $P$ est une primitive quelconque de $C_m$. La constante $K$ est déterminée par la condition initiale : $C_T(0) = C_F$ (coûts fixes), ce qui donne $K = C_F - P(0)$.
\emph{Interprétation :} la quantité $C_T(x) - C_F = P(x) - P(0)$ représente le \textbf{coût variable total} pour produire $x$ unités ; géométriquement, c'est l'aire sous la courbe du coût marginal entre $0$ et $x$.
\end{propriete}

\begin{experience}[Découverte : aire sous la courbe du coût marginal]
On reprend l'exemple du manioc : $C_T(x) = 2x^2 + 20x + 100$.
\begin{enumerate}
    \item Calculer $C_m(x) = C_T'(x)$.
    \item Tracer la droite $y = C_m(x)$ sur $[0, 10]$.
    \item Calculer l'aire du trapèze situé sous cette droite entre $0$ et $10$.
    \item Comparer cette aire à $C_V(10) = C_T(10) - C_F$.
\end{enumerate}
\textbf{Observation :} $C_m(x) = 4x + 20$. L'aire sous la droite entre $0$ et $10$ est l'aire d'un trapèze de bases $C_m(0) = 20$ et $C_m(10) = 60$, et de hauteur $10$ :
\[ \mathcal{A} = \frac{20 + 60}{2} \times 10 = 400. \]
Or $C_V(10) = 2(10)^2 + 20(10) = 200 + 200 = 400$. \textbf{L'égalité est exacte :} l'aire sous la courbe du coût marginal entre $0$ et $x$ est égale au coût variable total $C_V(x)$. On retrouve ce résultat par les primitives : une primitive de $C_m(x) = 4x + 20$ est $P(x) = 2x^2 + 20x$, et $P(10) - P(0) = 400 - 0 = 400$.
\end{experience}

\courssub{Analyse graphique : lecture de la pente et de l'aire}

La géométrie donne une vision intuitive des résultats algébriques : la dérivée se lit comme la \cle{pente de la tangente}, et la primitive comme une \cle{aire sous la courbe}.

\begin{popfigure}
\begin{tikzpicture}[scale=0.55, >=Stealth]
    % Axes
    \draw[->] (-0.5,0) -- (13,0) node[right] {$x$ (tonnes)};
    \draw[->] (0,-4) -- (0,11) node[above] {$B(x)$ (dizaines de milliers de FCFA)};
    % Grille légère
    \draw[popline, step=1] (0,-3.5) grid (12.5,10.5);
    % Courbe Bénéfice B(x) = -2x^2 + 40x - 100, échelle y = B/10
    \draw[popcurve, thick, domain=0:12.5, samples=200, variable=\x]
        plot ({\x}, {-0.2*\x*\x + 4*\x - 10});
    \node[popdark, anchor=south west] at (0.3, 9.2) {$B(x) = -2x^2 + 40x - 100$};
    % Point max x=10, B=100 -> y=10
    \fill[poporange] (10, 10) circle (3pt);
    \node[above, popdark] at (10, 10.2) {$M(10\,;\,100)$};
    \draw[dashed, popmuted] (10, 0) node[below] {$10$} -- (10, 10);
    % Tangente au maximum (horizontale)
    \draw[popblue, dashed, thick] (7, 10) -- (13, 10) node[right, popblue] {tangente horizontale ($B'(10)=0$)};
    % Points aux bornes
    \fill[popdark] (0, -10) circle (0pt);
    \fill[popgreen] (0, -10) circle (0pt);
\end{tikzpicture}
\legende{Courbe du bénéfice $B(x)$ (ordonnées divisées par 10). Le maximum est atteint en $x = 10$ : la tangente y est horizontale, ce qui traduit $B'(10) = 0$.}
\end{popfigure}

\begin{popfigure}
\begin{tikzpicture}[scale=0.7, >=Stealth]
    % Axes pour le coût marginal
    \draw[->] (-0.5,0) -- (13,0) node[right] {$x$};
    \draw[->] (0,-0.5) -- (0,7) node[above] {$C_m(x)$};
    \draw[popline, step=1] (0,0) grid (12.5,6.5);
    % Aire de 0 à 10 sous Cm(x) = 4x + 20, échelle y = Cm/10
    \fill[poporange!20, opacity=0.6] (0,0) -- (0,2) -- plot[domain=0:10, samples=50, variable=\x] ({\x}, {0.4*\x + 2}) -- (10,0) -- cycle;
    % Courbe Cm(x) = 4x + 20 -> y = 0.4x + 2
    \draw[popcurve, thick, poporange, domain=0:12.5, samples=100, variable=\x]
        plot ({\x}, {0.4*\x + 2});
    \node[poporange!80!black, anchor=south west] at (0.5, 5.6) {$C_m(x) = 4x + 20$};
    \node[poporange!80!black] at (5, 1.5) {Aire $= 400$};
    \draw[dashed, popmuted] (10, 0) node[below] {$10$} -- (10, 6);
    \draw[dashed, popmuted] (0, 2) node[left] {$20$} -- (10, 2);
    \draw[dashed, popmuted] (0, 6) node[left] {$60$} -- (10, 6);
\end{tikzpicture}
\legende{Courbe du coût marginal $C_m(x) = 4x + 20$ (ordonnées divisées par 10). L'aire colorée du trapèze entre $0$ et $10$ vaut $400$ : c'est le coût variable total pour produire 10 tonnes.}
\end{popfigure}

\begin{savaistu}[Application au Cameroun]
Les entreprises camerounaises de transformation agricole (cacao, café, manioc, plantain) et les PME du secteur secondaire (menuiserie, soudure, confection) utilisent ces modèles polynomiaux pour :
\begin{itemize}
    \item fixer le \textbf{prix de vente optimal} en fonction de la courbe de demande ;
    \item déterminer le \textbf{seuil de rentabilité} (point mort : $B(x) = 0$) ;
    \item planifier la \textbf{capacité de production} : ne pas produire au-delà de la quantité qui maximise le bénéfice si la demande ne suit pas ;
    \item négocier les \textbf{coûts fixes} (loyer, remboursement de crédit) en connaissant le volume minimal à atteindre pour être rentable.
\end{itemize}
La maîtrise de la dérivée et de la primitive est donc une compétence \emph{professionnelle} directe pour un technicien.
\end{savaistu}

\courssub{Synthèse : la démarche complète de rédaction (type Probatoire)}

Pour toute question de modélisation, la copie doit suivre une structure rigoureuse.

\begin{aretenir}[Structure de la réponse attendue]
\begin{enumerate}
    \item \textbf{Modélisation :} définir les variables ($x$, $C_T$, $R_T$, $B$), donner les expressions littérales, préciser le domaine $D$ (contraintes réelles : $x \geq 0$, capacité maximale, prix positif).
    \item \textbf{Étude mathématique :} calcul de la dérivée ($B'$ ou $C_m$), résolution de l'équation $B'(x) = 0$, tableau de variations (signe de $B'$).
    \item \textbf{Interprétation des résultats :} traduire $x_{opt}$ en « quantité à produire », $B_{max}$ en « bénéfice maximal », $C_m(x)$ en « coût de l'unité supplémentaire ».
    \item \textbf{Vérification de la cohérence :} $x_{opt} \in D$ ? Bénéfice positif ? Capacité de production respectée ? Si $x$ désigne un nombre d'objets entiers, comparer les deux entiers encadrant la valeur critique.
    \item \textbf{Conclusion :} phrase complète répondant à la question posée, avec les unités.
\end{enumerate}
\end{aretenir}

\begin{attention}[Pièges fréquents au Probatoire]
\begin{itemize}
    \item Oublier de vérifier que la valeur critique appartient au domaine $D$ : si $x_1 \notin D$, le maximum est atteint à une borne de l'intervalle.
    \item Conclure avec une valeur approchée non entière ($x \approx 41,67$ chaises) sans tester les entiers voisins.
    \item Confondre le maximum de la fonction (valeur de $B$) et le point où il est atteint (valeur de $x$).
    \item Oublier les unités dans la conclusion (tonnes, FCFA, milliers de FCFA).
\end{itemize}
\end{attention}

\begin{exercice}[Entraînement : production de chaises]
Un atelier de menuiserie produit $x$ chaises par semaine ($x \in [0, 50]$).
Le coût total (en FCFA) est $C_T(x) = x^2 + 50x + 2000$.
Le prix de vente unitaire (en FCFA) est $p(x) = 300 - 2x$.
\begin{enumerate}
    \item Exprimer le revenu $R_T(x)$ et le bénéfice $B(x)$.
    \item Déterminer le nombre de chaises à produire pour maximiser le bénéfice hebdomadaire.
    \item Calculer le coût marginal pour cette production optimale et l'interpréter.
    \item En utilisant une primitive, vérifier que le coût variable total pour cette production optimale est bien égal à l'aire sous la courbe du coût marginal entre $0$ et cette production.
\end{enumerate}
\end{exercice}

\begin{corrige}[Exercice : production de chaises]
\begin{enumerate}
    \item $R_T(x) = x \cdot p(x) = x(300 - 2x) = -2x^2 + 300x$.
    \[ B(x) = R_T(x) - C_T(x) = (-2x^2 + 300x) - (x^2 + 50x + 2000) = -3x^2 + 250x - 2000. \]
    \item $B'(x) = -6x + 250$. \quad $B'(x) = 0 \Leftrightarrow x = \dfrac{250}{6} = \dfrac{125}{3} \approx 41,67$.
    $B'$ est affine de coefficient directeur $-6 < 0$, d'où le tableau de variations sur $[0, 50]$ :
    \begin{center}
    \begin{tabular}{|c|ccccc|}\hline
    $x$ & $0$ & & $\frac{125}{3}$ & & $50$ \\ \hline
    $B'(x)$ & & $+$ & $0$ & $-$ & \\ \hline
    $B(x)$ & $-2000$ & $\nearrow$ & $\approx 3208$ & $\searrow$ & $3000$ \\ \hline
    \end{tabular}
    \end{center}
    \emph{Vérification :} $B(50) = -3(2500) + 250(50) - 2000 = -7500 + 12500 - 2000 = 3000$.
    Le maximum théorique est atteint en $x \approx 41,67$. Comme on ne produit pas une fraction de chaise, on compare les entiers voisins :
    \begin{align*}
    B(41) &= -3(41)^2 + 250(41) - 2000 = -5043 + 10250 - 2000 = 3207, \\
    B(42) &= -3(42)^2 + 250(42) - 2000 = -5292 + 10500 - 2000 = 3208.
    \end{align*}
    \textbf{Production optimale : 42 chaises par semaine}, pour un bénéfice maximal de \textbf{3208 FCFA}.
    \item $C_m(x) = C_T'(x) = 2x + 50$, donc $C_m(42) = 2(42) + 50 = 134$ FCFA.
    \emph{Interprétation :} la production de la 43\up{e} chaise coûterait environ 134 FCFA supplémentaires.
    \item Le coût variable est $C_V(x) = x^2 + 50x$. Une primitive de $C_m(x) = 2x + 50$ est $P(x) = x^2 + 50x$, et $P(0) = 0$, donc $C_V(x) = P(x) - P(0)$.
    Aire sous la courbe de $C_m$ entre $0$ et $42$ :
    \[ P(42) - P(0) = 42^2 + 50 \times 42 = 1764 + 2100 = 3864. \]
    Or $C_V(42) = 42^2 + 50 \times 42 = 3864$. \textbf{L'égalité est vérifiée :} l'aire sous la courbe du coût marginal entre $0$ et $42$ est bien égale au coût variable total, soit 3864 FCFA.
\end{enumerate}
\end{corrige}

\newpage

\coursec{Activité d'intégration}

\courssub{Objectifs de l'activité}
Mobiliser l'ensemble des savoirs et savoir-faire de la leçon « Dérivée et primitives » dans une situation professionnelle authentique : modéliser l'activité économique d'une petite unité de transformation de manioc, déterminer les fonctions économiques associées (revenu, bénéfice), utiliser la dérivée pour analyser la variation du coût marginal et optimiser le bénéfice, exploiter le lien dérivation/primitives pour vérifier la cohérence du modèle, et communiquer les résultats sous forme graphique et rédactionnelle adaptée à un décideur non mathématicien.

\courssub{Situation}
Dans la région du Centre, à proximité de Yaoundé, une coopérative de jeunes agriculteurs a lancé une petite unité de transformation de manioc en farine (fariña, gari, farine de manioc pour le \emph{water fufu} et le \emph{bobolo}). L'atelier fonctionne de manière artisanale mais structurée : achat des tubercules aux producteurs locaux, épluchage, lavage, râpage, pressage, tamisage, torréfaction et conditionnement en sacs de 50 kg.

Le gérant, M. \textsc{Fouda}, a collecté des données sur les coûts de production mensuels (main-d'œuvre, énergie, maintenance, achat matière première, transport) pour différents volumes de production. Après analyse statistique, le service comptabilité a modélisé le \textbf{coût total mensuel} $C(x)$ (en milliers de francs CFA) en fonction de la quantité $x$ (en tonnes de farine produite) par la fonction polynôme du troisième degré :
\[
C(x) = 0,5x^{3} - 3x^{2} + 10x + 100 \qquad \text{avec } x \in [0\,;\, 10].
\]
Le domaine de définition $[0\,;\, 10]$ correspond à la capacité maximale actuelle de l'atelier (10 tonnes/mois).

Le prix de vente unitaire de la tonne de farine a été fixé, après étude de marché, à \textbf{50 milliers de francs CFA} (soit 50\,000 FCFA/tonne), prix compétitif sur le marché local de Mfoundi et des environs.

M. \textsc{Fouda} souhaite prendre des décisions éclairées pour le prochain trimestre : quelle quantité produire pour maximiser le bénéfice ? Quel est le coût marginal à ce niveau ? L'atelier est-il rentable ? Il demande à son équipe technique (les élèves) de réaliser une étude complète et de lui présenter un rapport clair, chiffré et graphique.

\courssub{Consignes}
\emph{Travail à réaliser en groupe de 3 ou 4 élèves, restitution orale et écrite.}

\begin{enumerate}
    \item \textbf{Modélisation du revenu et du bénéfice.}
    \begin{enumerate}
        \item Exprimer la fonction revenu total $R(x)$ (en milliers de FCFA) en fonction de $x$.
        \item Déduire l'expression de la fonction bénéfice $B(x) = R(x) - C(x)$. Simplifier.
        \item Préciser le domaine de définition commun aux fonctions $C$, $R$ et $B$ dans le contexte.
    \end{enumerate}

    \item \textbf{Analyse du coût marginal.}
    \begin{enumerate}
        \item Calculer la dérivée $C'(x)$ de la fonction coût total. Rappeler l'interprétation économique de $C'(x)$ (coût marginal).
        \item Calculer $C'(4)$ et $C'(8)$. Interpréter concrètement ces deux valeurs pour le gérant.
        \item Étudier les variations de $C'(x)$ sur $[0\,;\, 10]$. Que peut-on dire de l'évolution du coût marginal lorsque la production augmente ?
    \end{enumerate}

    \item \textbf{Optimisation du bénéfice.}
    \begin{enumerate}
        \item Calculer la dérivée $B'(x)$ de la fonction bénéfice.
        \item Résoudre l'équation $B'(x) = 0$ sur $[0\,;\, 10]$. Combien de solutions ? Laquelle est économiquement pertinente ?
        \item Construire le tableau de variations de $B$ sur $[0\,;\, 10]$.
        \item Déterminer la quantité $x_{\text{opt}}$ qui maximise le bénéfice et la valeur du bénéfice maximal $B(x_{\text{opt}})$.
        \item Vérifier que ce maximum est bien un maximum global sur l'intervalle.
    \end{enumerate}

    \item \textbf{Lien dérivation -- primitives (vérification).}
    \begin{enumerate}
        \item Déterminer une primitive $F(x)$ de la fonction dérivée $C'(x)$ sur $\mathbb{R}$ qui s'annule en $0$ ($F(0)=0$).
        \item Calculer $F(x)$ pour $x = x_{\text{opt}}$.
        \item Comparer $F(x_{\text{opt}})$ à $C(x_{\text{opt}}) - C(0)$. Commenter le résultat au regard de la relation entre une fonction et sa dérivée.
    \end{enumerate}

    \item \textbf{Représentation graphique et analyse visuelle.}
    \begin{enumerate}
        \item Sur un même repère (papier millimétré ou logiciel type GeoGebra/Excel), tracer les courbes représentatives $\mathcal{C}_C$, $\mathcal{C}_R$, $\mathcal{C}_B$ et $\mathcal{C}_{C'}$.
        \item Choisir des unités graphiques adaptées (ex : 1 cm pour 1 tonne en abscisses, 1 cm pour 50 milliers FCFA en ordonnées pour $C,R,B$ ; 1 cm pour 10 milliers FCFA/tonne pour $C'$).
        \item Repérer graphiquement : le seuil de rentabilité, le point de bénéfice maximal, la zone de perte, la zone de profit.
        \item Commenter la position relative des courbes $\mathcal{C}_R$ et $\mathcal{C}_C$ au voisinage de l'optimum.
    \end{enumerate}

    \item \textbf{Synthèse et recommandation.}
    Rédiger une note de synthèse d'une demi-page, adressée à M. \textsc{Fouda}, résumant les résultats chiffrés clés (quantité optimale, bénéfice maximal, coût marginal à l'optimum, seuil de rentabilité) et formulant une recommandation claire sur le volume de production à viser pour le trimestre prochain. Utiliser un langage professionnel accessible.
\end{enumerate}

\courssub{Ressources à mobiliser}
\begin{itemize}
    \item \textbf{Dérivées usuelles :} $(x^n)' = nx^{n-1}$, $(k)' = 0$ ($k$ constante), $(u+v)' = u'+v'$, $(ku)' = ku'$.
    \item \textbf{Primitives des fonctions polynômes :} une primitive de $x^n$ sur $\mathbb{R}$ est $\dfrac{x^{n+1}}{n+1}$ pour $n \in \mathbb{N}$ ; une primitive d'une constante $k$ est $kx$ ; on primitive terme à terme.
    \item \textbf{Recherche d'extremum :} si $f$ est dérivable sur $[a;b]$, tout extremum atteint en un point intérieur $c$ vérifie $f'(c)=0$. Si $f'$ s'annule en $c$ en changeant de signe de $+$ en $-$, alors $f$ admet un maximum en $c$.
    \item \textbf{Relation fonction -- primitive :} si $F$ est une primitive de $f$ sur $[a;b]$, alors $F(b)-F(a)$ mesure l'accumulation de $f$ entre $a$ et $b$. Ici on utilise : $C(x) - C(0) = F(x) - F(0)$ où $F$ est une primitive de $C'$.
    \item \textbf{Notions économiques :} Revenu $R(x) = p \times x$ ; Bénéfice $B(x) = R(x) - C(x)$ ; Coût marginal $C'(x)$ ; Seuil de rentabilité : plus petite valeur de $x$ telle que $B(x)=0 \iff R(x)=C(x)$.
\end{itemize}

\courssub{Démarche et corrigé commenté}
\begin{exoresolu}{Corrigé détaillé de l'activité d'intégration}
\underline{\textbf{1. Modélisation du revenu et du bénéfice}}

\begin{enumerate}
    \item Le prix de vente unitaire est $p = 50$ (milliers FCFA/tonne). Le revenu total pour $x$ tonnes vendues est :
    \[
    R(x) = p \times x = 50x.
    \]
    \item Le bénéfice est la différence Revenu -- Coût :
    \begin{align*}
    B(x) &= R(x) - C(x) \\
         &= 50x - (0,5x^{3} - 3x^{2} + 10x + 100) \\
         &= -0,5x^{3} + 3x^{2} + 40x - 100.
    \end{align*}
    \item Le domaine d'étude est imposé par la capacité de l'atelier : $x \in [0\,;\, 10]$. Les trois fonctions $C$, $R$, $B$ sont des polynômes, donc définies sur $\mathbb{R}$, mais l'étude économique se restreint à $[0\,;\, 10]$.
\end{enumerate}

\underline{\textbf{2. Analyse du coût marginal}}

\begin{enumerate}
    \item Dérivée de $C(x) = 0,5x^{3} - 3x^{2} + 10x + 100$, terme à terme avec $(x^n)'=nx^{n-1}$ :
    \[
    C'(x) = 3 \times 0,5 x^{2} - 2 \times 3 x + 10 = 1,5x^{2} - 6x + 10.
    \]
    \textbf{Interprétation :} $C'(x)$ représente le \textbf{coût marginal}, c'est-à-dire le coût approximatif de production d'une tonne supplémentaire lorsque la production est déjà de $x$ tonnes. Unité : milliers FCFA/tonne.
    \item Calculs :
    \begin{align*}
    C'(4) &= 1,5 \times 16 - 6 \times 4 + 10 = 24 - 24 + 10 = 10. \\
    C'(8) &= 1,5 \times 64 - 6 \times 8 + 10 = 96 - 48 + 10 = 58.
    \end{align*}
    \textbf{Interprétation :} à 4 tonnes produites, la tonne supplémentaire coûte environ 10\,000 FCFA ; à 8 tonnes, elle coûte environ 58\,000 FCFA, c'est-à-dire \emph{plus cher que le prix de vente} (50\,000 FCFA). Le coût marginal augmente fortement : l'atelier approche de sa capacité maximale, les heures supplémentaires, l'usure des machines et la logistique renchérissent le coût unitaire additionnel.
    \item Variations de $C'(x) = 1,5x^{2} - 6x + 10$ sur $[0\,;\, 10]$.
    $C'$ est un polynôme, donc dérivable, et $(C')'(x) = 3x - 6$.
    $(C')'(x) = 0 \iff x = 2$ ; $(C')'(x) < 0$ pour $x < 2$ et $(C')'(x) > 0$ pour $x > 2$.
    Valeurs : $C'(0) = 10$, $C'(2) = 1,5 \times 4 - 12 + 10 = 4$, $C'(10) = 150 - 60 + 10 = 100$.
    \begin{center}
    \begin{tabular}{|c|ccccc|}\hline
    $x$ & $0$ & & $2$ & & $10$ \\ \hline
    $(C')'(x)$ & & $-$ & $0$ & $+$ & \\ \hline
    $C'(x)$ & $10$ & $\searrow$ & $4$ & $\nearrow$ & $100$ \\ \hline
    \end{tabular}
    \end{center}
    Le coût marginal décroît d'abord (économies d'échelle, meilleure répartition des coûts fixes) jusqu'à 2 tonnes (minimum de 4 milliers FCFA/tonne), puis croît fortement (déséconomies d'échelle, saturation de l'atelier).
\end{enumerate}

\underline{\textbf{3. Optimisation du bénéfice}}

\begin{enumerate}
    \item $B(x) = -0,5x^{3} + 3x^{2} + 40x - 100$, donc :
    \[
    B'(x) = -1,5x^{2} + 6x + 40.
    \]
    \item Résolution de $B'(x) = 0 \iff -1,5x^{2} + 6x + 40 = 0$.
    Multiplions les deux membres par $-2$ pour simplifier : $3x^{2} - 12x - 80 = 0$.
    Discriminant : $\Delta = (-12)^{2} - 4 \times 3 \times (-80) = 144 + 960 = 1104 > 0$ : deux racines réelles.
    $\sqrt{\Delta} = \sqrt{1104} = \sqrt{16 \times 69} = 4\sqrt{69} \approx 33,23$.
    \[
    x_1 = \frac{12 - 4\sqrt{69}}{6} = 2 - \frac{2}{3}\sqrt{69} \approx -3,54 \quad (\text{rejetée, hors } [0;10])
    \]
    \[
    x_2 = \frac{12 + 4\sqrt{69}}{6} = 2 + \frac{2}{3}\sqrt{69} \approx 7,54.
    \]
    Seule $x_{\text{opt}} = 2 + \dfrac{2}{3}\sqrt{69} \approx 7,54$ tonnes est économiquement pertinente.
    \item Tableau de variations de $B$ sur $[0\,;\, 10]$.
    $B'$ est un trinôme du second degré de coefficient dominant $-1,5 < 0$, de racines $x_1 < 0$ et $x_2 \approx 7,54$. Donc $B'(x) > 0$ sur $[0\,;\, x_2[$ (entre les racines) et $B'(x) < 0$ sur $]x_2\,;\, 10]$.
    Valeurs : $B(0) = -100$ ; $B(10) = -0,5 \times 1000 + 3 \times 100 + 400 - 100 = -500 + 300 + 400 - 100 = 100$.
    Pour $x_2 \approx 7,54$ : $x_2^{2} \approx 56,82$ et $x_2^{3} \approx 428,3$, d'où
    \[
    B(x_2) \approx -0,5 \times 428,3 + 3 \times 56,82 + 40 \times 7,54 - 100 \approx -214,1 + 170,5 + 301,5 - 100 \approx 157,9.
    \]
    \begin{center}
    \begin{tabular}{|c|ccccc|}\hline
    $x$ & $0$ & & $x_2 \approx 7,54$ & & $10$ \\ \hline
    $B'(x)$ & & $+$ & $0$ & $-$ & \\ \hline
    $B(x)$ & $-100$ & $\nearrow$ & $\approx 157,9$ & $\searrow$ & $100$ \\ \hline
    \end{tabular}
    \end{center}
    \item Le bénéfice maximal est d'environ \textbf{158 milliers de FCFA} (soit environ 158\,000 FCFA), atteint pour une production de \textbf{7,54 tonnes} par mois.
    \item Vérification du maximum global : $B$ est croissante sur $[0\,;\, x_2]$ puis décroissante sur $[x_2\,;\, 10]$, donc le maximum de $B$ sur $[0\,;\, 10]$ est atteint en $x_2$. De plus $B(x_2) \approx 157,9 > B(10) = 100 > B(0) = -100$ : c'est bien le maximum global sur l'intervalle.
\end{enumerate}

\underline{\textbf{4. Lien dérivation -- primitives (vérification)}}

\begin{enumerate}
    \item $C'(x) = 1,5x^{2} - 6x + 10$. On primitive terme à terme avec la formule $\dfrac{x^{n+1}}{n+1}$ :
    \[
    F(x) = 1,5 \times \frac{x^{3}}{3} - 6 \times \frac{x^{2}}{2} + 10x = 0,5x^{3} - 3x^{2} + 10x.
    \]
    Vérification : $F'(x) = 1,5x^{2} - 6x + 10 = C'(x)$ et $F(0) = 0$ : $F$ est bien la primitive de $C'$ qui s'annule en $0$.
    \item Pour $x_{\text{opt}} \approx 7,54$ :
    \[
    F(7,54) = 0,5 \times 428,3 - 3 \times 56,82 + 10 \times 7,54 \approx 214,1 - 170,5 + 75,4 \approx 119,0.
    \]
    \item D'autre part :
    \[
    C(x_{\text{opt}}) - C(0) = \left(0,5x_{\text{opt}}^{3} - 3x_{\text{opt}}^{2} + 10x_{\text{opt}} + 100\right) - 100 = 0,5x_{\text{opt}}^{3} - 3x_{\text{opt}}^{2} + 10x_{\text{opt}} = F(x_{\text{opt}}).
    \]
    \textbf{Commentaire :} on retrouve exactement l'égalité $C(x) - C(0) = F(x) - F(0)$, conséquence du fait que $F$ est une primitive de $C'$. La partie variable du coût total (le \textbf{coût variable}) est la primitive du coût marginal qui s'annule en $0$. Le terme constant $C(0) = 100$ représente les \textbf{coûts fixes} (loyer, assurances, amortissements) qui ne dépendent pas du volume produit et qui disparaissent lors de la dérivation.
\end{enumerate}

\underline{\textbf{5. Représentation graphique et analyse visuelle}}

\begin{popfigure}
\begin{tikzpicture}[scale=0.75, every node/.style={font=\small}]
% Axes
\draw[->] (-0.5,0) -- (11,0) node[right] {$x$ (tonnes)};
\draw[->] (0,-2.5) -- (0,6.5) node[above] {$y$ (échelles réduites)};

% Graduations
\foreach \x in {1,...,10} \draw[popline] (\x,-0.1) -- (\x,0.1) node[below,yshift=-2pt] {\x};
\foreach \y in {-2,-1,1,2,3,4,5,6} \draw[popline] (-0.1,\y) -- (0.1,\y) node[left,xshift=-2pt] {\y};

% Courbes normalisées :
% C/50 = 0.01x^3 - 0.06x^2 + 0.2x + 2
% R/50 = x
% B/50 = -0.01x^3 + 0.06x^2 + 0.8x - 2
% C'/10 = 0.15x^2 - 0.6x + 1
\draw[popblue, thick, domain=0:10, samples=100, variable=\x] plot ({\x}, {0.01*\x^3 - 0.06*\x^2 + 0.2*\x + 2}) node[right, pos=0.92] {$C/50$};
\draw[popgreen, thick, domain=0:6.2, samples=20] plot ({\x}, {\x}) node[right, pos=0.95] {$R/50$};
\draw[poporange, thick, domain=0:10, samples=100, variable=\x] plot ({\x}, {-0.01*\x^3 + 0.06*\x^2 + 0.8*\x - 2}) node[right, pos=0.55] {$B/50$};
\draw[poppurple, thick, dashed, domain=0:6, samples=50, variable=\x] plot ({\x}, {0.15*\x^2 - 0.6*\x + 1}) node[right, pos=0.95] {$C'/10$};

% Points clés
\fill[popblue] (0,2) circle (2pt) node[above left] {$C(0)=100$};
\fill[poporange] (0,-2) circle (2pt) node[below left] {$B(0)=-100$};

% Seuil de rentabilité : B=0 pour x ~ 2.26 ; R/50 = 2.26
\fill[red] (2.26, 2.26) circle (2.5pt) node[above right] {Seuil de rentabilité};

% Optimum B max : x=7.54, B/50 ~ 3.16
\fill[popgold] (7.54, 3.16) circle (3pt) node[above] {Max de $B$};

% Min C' : x=2, C'/10 = 0.4
\fill[poppurple] (2, 0.4) circle (2pt) node[below] {Min de $C'$};

\end{tikzpicture}
\legende{Représentation schématique des courbes $C/50$, $R/50$, $B/50$ et $C'/10$ (ordonnées divisées par 50 pour $C,R,B$ et par 10 pour $C'$, afin de tout faire tenir sur un même graphique). En vraie grandeur, on utilise deux axes d'ordonnées.}
\end{popfigure}

\begin{attention}[Conseil pour le tracé réel]
Sur papier ou logiciel, utilisez \textbf{deux axes d'ordonnées} :
\begin{itemize}
    \item axe de gauche (milliers FCFA) pour $C$, $R$, $B$ : 1 cm $\rightarrow$ 50 milliers FCFA ;
    \item axe de droite (milliers FCFA/tonne) pour $C'$ : 1 cm $\rightarrow$ 20 milliers FCFA/t.
\end{itemize}
Cela évite d'écraser la courbe $C'$, qui prend des valeurs beaucoup plus petites que $C$, $R$, $B$.
\end{attention}

\textbf{Commentaires graphiques attendus :}
\begin{itemize}
    \item \textbf{Seuil de rentabilité :} les courbes $\mathcal{C}_R$ et $\mathcal{C}_C$ se coupent une seule fois sur $[0\,;\,10]$, en $x \approx 2,26$ tonnes (résolution approchée de $B(x)=0$ : $B(2) = -4 + 12 + 80 - 100 = -12 < 0$ et $B(2,5) \approx -7,8 + 18,75 + 100 - 100 = 10,9 > 0$). En dessous de 2,26 tonnes, $\mathcal{C}_C$ est au-dessus de $\mathcal{C}_R$ : \textbf{zone de perte} ($B<0$). Au-dessus, et jusqu'à la capacité maximale de 10 tonnes, $\mathcal{C}_R$ reste au-dessus de $\mathcal{C}_C$ : \textbf{zone de profit} ($B>0$). En effet $B(10) = 100 > 0$ : l'atelier reste bénéficiaire même à pleine capacité, mais avec un bénéfice moindre qu'à l'optimum.
    \item \textbf{Optimum :} le maximum de $\mathcal{C}_B$ correspond à l'abscisse où les tangentes à $\mathcal{C}_R$ et $\mathcal{C}_C$ sont parallèles : $B'(x)=0 \iff R'(x)=C'(x) \iff 50 = C'(x)$. Vérification : $C'(7,54) \approx 1,5 \times 56,82 - 6 \times 7,54 + 10 \approx 85,2 - 45,2 + 10 = 50$. C'est la \textbf{règle d'or du producteur} : à l'optimum, revenu marginal = coût marginal.
    \item \textbf{Coût marginal :} la courbe $\mathcal{C}_{C'}$ passe par un minimum en $x=2$ (4 milliers FCFA/t), puis coupe la droite d'équation $y = 50$ (prix de vente) en $x \approx 7,54$ : au-delà, produire une tonne de plus coûte plus cher que ce qu'elle rapporte.
\end{itemize}

\underline{\textbf{6. Synthèse et recommandation (exemple de note pour M. Fouda)}}

\begin{center}
\begin{tabularx}{\linewidth}{|X|}\hline
\rowcolor{popblueL}
\textbf{NOTE DE SYNTHÈSE -- OPTIMISATION DE LA PRODUCTION DE FARINE DE MANIOC} \\ \hline
\textbf{Objet :} volume de production optimal pour le trimestre à venir. \\[4pt]
\textbf{Contexte :} coût total $C(x)=0,5x^3-3x^2+10x+100$ (milliers FCFA), prix de vente 50\,000 FCFA/t, capacité maximale 10 t/mois. \\[4pt]
\textbf{Résultats clés :}
\begin{itemize}
    \item \textbf{Seuil de rentabilité :} environ \textbf{2,3 tonnes/mois}. En dessous, l'atelier est en perte ; au-dessus, il est bénéficiaire sur toute la plage jusqu'à 10 tonnes.
    \item \textbf{Production optimale :} \textbf{7,54 tonnes/mois} (soit environ 151 sacs de 50 kg).
    \item \textbf{Bénéfice maximal mensuel :} environ \textbf{158\,000 FCFA}.
    \item \textbf{Coût marginal à l'optimum :} 50\,000 FCFA/tonne, exactement égal au prix de vente.
    \item \textbf{Coûts fixes mensuels :} 100\,000 FCFA (couverts dès que la production dépasse le seuil de rentabilité).
\end{itemize}
\textbf{Recommandation :}
viser une production mensuelle de \textbf{7,5 tonnes} (arrondi pratique). Cela assure le bénéfice maximal tout en restant en deçà de la capacité maximale : à 10 tonnes, le bénéfice retombe à 100\,000 FCFA car le coût marginal dépasse le prix de vente. Ne jamais descendre sous 2,3 tonnes (zone de perte). Prévoir un stock de sécurité de matière première pour tenir ce rythme.
\\[4pt]
\textit{Fait à Yaoundé, le \today. L'équipe technique.}
\\ \hline
\end{tabularx}
\end{center}
\end{exoresolu}

\begin{attention}[Erreurs fréquentes à éviter]
\begin{itemize}
    \item \textbf{Erreur de signe dans le bénéfice :} $B(x) = R(x) - C(x)$ impose de changer le signe de \emph{tous} les termes de $C$ : $-(0,5x^3 - 3x^2 + 10x + 100) = -0,5x^3 + 3x^2 - 10x - 100$, puis on ajoute $50x$ qui donne $+40x$.
    \item \textbf{Oublier le domaine :} une racine de $B'(x)=0$ hors de $[0\,;\,10]$ (ici $x_1 \approx -3,54$) doit être rejetée et ce rejet doit être \emph{justifié} dans la rédaction.
    \item \textbf{Confondre coût moyen et coût marginal :} le coût marginal est la \emph{dérivée} $C'(x)$, pas le quotient $\dfrac{C(x)}{x}$.
    \item \textbf{Constante d'intégration :} une primitive est définie à une constante près ; la condition $F(0)=0$ permet de choisir celle qui correspond au coût variable. Ne pas oublier de vérifier $F' = C'$.
    \item \textbf{Unités :} $C$, $R$, $B$ sont en milliers de FCFA ; $C'$ est en milliers de FCFA \emph{par tonne}. Mélanger les unités fausse les interprétations et le graphique.
\end{itemize}
\end{attention}

\courssub{Bilan de l'activité}
Cette activité d'intégration a permis de mettre en évidence, dans un contexte professionnel camerounais réaliste, l'articulation complète des notions du chapitre :
\begin{itemize}
    \item La \textbf{modélisation} par des fonctions polynômes (degré 3 pour le coût, degré 1 pour le revenu) issue de données techniques.
    \item L'usage de la \textbf{dérivée} comme outil d'analyse : coût marginal (dérivée du coût), optimisation du bénéfice (annulation de la dérivée, tableau de variations, étude du signe de $B'$).
    \item La compréhension du lien \textbf{dérivée -- primitive} : le coût variable est une primitive du coût marginal, les coûts fixes étant la constante $C(0)$ qui disparaît par dérivation.
    \item La \textbf{lecture graphique} multi-courbes (échelles doubles) pour valider les résultats analytiques : seuil de rentabilité, maximum de bénéfice, égalité $R'(x)=C'(x)$ à l'optimum.
    \item La \textbf{communication professionnelle} : traduction des résultats mathématiques (valeurs exactes, approximations, unités) en recommandations décisionnelles claires pour un gestionnaire non spécialiste.
\end{itemize}
L'élève a ainsi mobilisé l'ensemble des savoir-faire visés : calculer, justifier (tableau de variations, rejet d'une racine hors domaine), interpréter, modéliser, représenter, communiquer.

\newpage

\coursec{Méthodes et exercices résolus}

\coursec{Fonctions : Dérivée et primitives}

\courssub{Rappels sur la dérivée et interprétation géométrique}

\begin{definition}[Nombre dérivé et tangente]
Soit $f$ une fonction définie sur un intervalle $I$ et $a \in I$. Si le taux d'accroissement $\dfrac{f(x) - f(a)}{x - a}$ admet une limite finie lorsque $x \to a$, on dit que $f$ est \cle{dérivable} en $a$. Cette limite est le \cle{nombre dérivé} de $f$ en $a$, noté $f'(a)$.
\[ f'(a) = \lim_{x \to a} \frac{f(x) - f(a)}{x - a} = \lim_{h \to 0} \frac{f(a+h) - f(a)}{h}. \]
\emph{Interprétation géométrique :} $f'(a)$ est le coefficient directeur de la tangente à la courbe $\mathcal{C}_f$ au point $A(a\,;\,f(a))$.
\end{definition}

\begin{propriete}[Tableau des dérivées usuelles -- À connaître par cœur]
Les fonctions suivantes sont dérivables sur leur ensemble de définition :
\begin{center}
\begin{tabularx}{\linewidth}{|l|X|l|}\hline
\rowcolor{popblueL}
\textbf{Fonction $f(x)$} & \textbf{Dérivée $f'(x)$} & \textbf{Domaine} \\ \hline
$k$ (constante) & $0$ & $\mathbb{R}$ \\
$x^n$ ($n \in \mathbb{N}^*$) & $n x^{n-1}$ & $\mathbb{R}$ \\
$\dfrac{1}{x} = x^{-1}$ & $-\dfrac{1}{x^2}$ & $\mathbb{R}^*$ \\
$\sqrt{x} = x^{1/2}$ & $\dfrac{1}{2\sqrt{x}}$ & $]0\,;\,+\infty[$ \\ \hline
\end{tabularx}
\end{center}
\emph{Règles de linéarité :} $(u+v)' = u' + v'$ \quad $(ku)' = k\,u'$ \quad ($k \in \mathbb{R}$).
\end{propriete}

\courssub{Règles de dérivation usuelles}

\begin{methode}[Dériver une fonction polynomiale ou usuelle]
Pour calculer la dérivée d'une fonction :
\begin{enumerate}
\item \textbf{Identifier} la structure : somme, produit par une constante, puissance entière ou rationnelle simple.
\item \textbf{Choisir la formule adaptée} dans le tableau des dérivées usuelles.
\item \textbf{Dériver terme à terme} en appliquant la linéarité : la dérivée d'une somme est la somme des dérivées.
\item \textbf{Simplifier} le résultat (réduire les coefficients, écrire sous forme fractionnaire si nécessaire).
\item \textbf{Vérifier} mentalement : la dérivation fait baisser le degré d'un polynôme de 1.
\end{enumerate}
\end{methode}

\begin{aretenir}[Réflexes à retenir pour la dérivation]
$(k)' = 0$ (constante) ; $(ax+b)' = a$ ; $(x^2)' = 2x$ ; $(x^3)' = 3x^2$ ; $(x^n)' = n x^{n-1}$.
\end{aretenir}

\begin{exoresolu}[Dérivée d'un polynôme (type Probatoire)]
Soit $f$ la fonction définie sur $\mathbb{R}$ par $f(x) = 2x^3 - 5x^2 + 3x - 7$. Déterminer $f'(x)$ puis calculer $f'(2)$.
\tcblower
\textbf{Solution détaillée.}

La fonction $f$ est un polynôme, donc dérivable sur $\mathbb{R}$. On dérive terme à terme en utilisant $(x^n)' = n x^{n-1}$ et $(ku)' = k\,u'$ :
\begin{align*}
f'(x) &= (2x^3)' - (5x^2)' + (3x)' - (7)' \\
&= 2 \times 3x^2 - 5 \times 2x + 3 \times 1 - 0 \\
&= 6x^2 - 10x + 3.
\end{align*}
On calcule ensuite $f'(2)$ en remplaçant $x$ par $2$ :
\[ f'(2) = 6 \times 2^2 - 10 \times 2 + 3 = 24 - 20 + 3 = 7. \]
\textbf{Conclusion :} $f'(x) = 6x^2 - 10x + 3$ et $f'(2) = 7$.
\end{exoresolu}

\courssub{Dérivation d'un produit et d'un quotient}

\begin{methode}[Utiliser les formules de produit et de quotient]
\begin{enumerate}
\item \textbf{Repérer} les deux facteurs $u$ et $v$ (ou le numérateur et le dénominateur).
\item \textbf{Calculer séparément} $u'$ et $v'$.
\item \textbf{Appliquer la formule} :
      \begin{itemize}
      \item Produit : $(uv)' = u'v + uv'$.
      \item Quotient : $\left(\dfrac{u}{v}\right)' = \dfrac{u'v - uv'}{v^2}$ \quad ($v \neq 0$).
      \end{itemize}
\item \textbf{Développer et réduire} le numérateur ; ne jamais développer le dénominateur d'un quotient (il reste sous forme factorisée ou carrée).
\end{enumerate}
\end{methode}

\begin{attention}[Piège classique : produit et quotient]
La dérivée d'un produit n'est \textbf{pas} le produit des dérivées : $(uv)' \neq u' \times v'$.
De même, $\left(\dfrac{u}{v}\right)' \neq \dfrac{u'}{v'}$.
Respecter scrupuleusement l'ordre des termes au numérateur du quotient : $u'v - uv'$.
\end{attention}

\begin{exoresolu}[Dérivée d'un produit]
Soit $g$ la fonction définie sur $\mathbb{R}$ par $g(x) = (x^2 + 1)(3x - 2)$. Calculer $g'(x)$.
\tcblower
\textbf{Solution détaillée.}

On pose $u(x) = x^2 + 1$ et $v(x) = 3x - 2$. Alors $u'(x) = 2x$ et $v'(x) = 3$.

On applique la formule du produit $(uv)' = u'v + uv'$ :
\begin{align*}
g'(x) &= u'(x)\,v(x) + u(x)\,v'(x) \\
&= 2x(3x - 2) + (x^2 + 1) \times 3 \\
&= 6x^2 - 4x + 3x^2 + 3 \\
&= 9x^2 - 4x + 3.
\end{align*}
\textbf{Vérification} (en développant d'abord) : $g(x) = 3x^3 - 2x^2 + 3x - 2$, donc $g'(x) = 9x^2 - 4x + 3$. Les deux méthodes donnent le même résultat.

\textbf{Conclusion :} $g'(x) = 9x^2 - 4x + 3$.
\end{exoresolu}

\courssub{Équation de la tangente à une courbe}

\begin{methode}[Équation de la tangente en un point]
Pour déterminer l'équation de la tangente à la courbe de $f$ au point d'abscisse $a$ :
\begin{enumerate}
\item \textbf{Calculer} $f(a)$ : c'est l'ordonnée du point de contact $A(a\,;\,f(a))$.
\item \textbf{Calculer} $f'(x)$ puis $f'(a)$ : c'est le \cle{coefficient directeur} de la tangente.
\item \textbf{Écrire la formule} : $y = f'(a)(x - a) + f(a)$.
\item \textbf{Développer et réduire} pour obtenir l'équation sous la forme $y = mx + p$.
\end{enumerate}
\end{methode}

\begin{exoresolu}[Tangente à une parabole]
Soit $f(x) = x^2 - 4x + 5$ et $\mathcal{C}_f$ sa courbe représentative. Déterminer une équation de la tangente à $\mathcal{C}_f$ au point $A$ d'abscisse $3$.
\tcblower
\textbf{Solution détaillée.}

\textbf{Étape 1 :} ordonnée du point de contact.
\[ f(3) = 3^2 - 4 \times 3 + 5 = 9 - 12 + 5 = 2. \]
Donc $A(3\,;\,2)$.

\textbf{Étape 2 :} nombre dérivé. $f'(x) = 2x - 4$, donc $f'(3) = 2 \times 3 - 4 = 2$.

\textbf{Étape 3 :} on applique la formule $y = f'(a)(x - a) + f(a)$ avec $a = 3$ :
\[ y = 2(x - 3) + 2. \]

\textbf{Étape 4 :} on développe : $y = 2x - 6 + 2$, soit $y = 2x - 4$.

\textbf{Conclusion :} la tangente à $\mathcal{C}_f$ au point $A(3\,;\,2)$ a pour équation $y = 2x - 4$.
\end{exoresolu}

\courssub{Primitives des fonctions polynomiales}

\begin{definition}[Primitive d'une fonction]
Soit $f$ une fonction continue sur un intervalle $I$. Une fonction $F$ dérivable sur $I$ est une \cle{primitive} de $f$ sur $I$ si $F'(x) = f(x)$ pour tout $x \in I$.
Si $F$ est une primitive de $f$, alors toutes les primitives de $f$ sur $I$ sont les fonctions $F + C$ avec $C \in \mathbb{R}$.
\end{definition}

\begin{propriete}[Tableau des primitives usuelles -- Correspondance avec les dérivées]
\begin{center}
\begin{tabularx}{\linewidth}{|l|X|l|}\hline
\rowcolor{popgreenL}
\textbf{Fonction $f(x)$} & \textbf{Primitive $F(x)$ (sur l'intervalle convenable)} & \textbf{Condition} \\ \hline
$k$ (constante) & $k x + C$ & $\mathbb{R}$ \\
$x^n$ ($n \in \mathbb{N}$) & $\dfrac{x^{n+1}}{n+1} + C$ & $\mathbb{R}$ ($n \neq -1$) \\
$x^{-1} = \dfrac{1}{x}$ & $\ln|x| + C$ & $\mathbb{R}^*$ (Hors programme Terminale Tech, donné pour culture) \\
$x^{1/2} = \sqrt{x}$ & $\dfrac{2}{3}x^{3/2} + C$ & $[0\,;\,+\infty[$ \\ \hline
\end{tabularx}
\end{center}
\emph{Règles de linéarité :} primitive de $u+v$ = primitive de $u$ + primitive de $v$ ; primitive de $k u$ = $k \times$ primitive de $u$.
\end{propriete}

\begin{methode}[Trouver une primitive d'un polynôme]
Pour déterminer une primitive $F$ d'une fonction polynomiale $f$ :
\begin{enumerate}
\item \textbf{Utiliser la formule} : une primitive de $x^n$ est $\dfrac{x^{n+1}}{n+1}$ (on \emph{augmente le degré de 1} et on \emph{divise par le nouveau degré}).
\item \textbf{Primitiver terme à terme} : une primitive d'une somme est la somme des primitives ; une primitive de $k \times f$ est $k \times F$.
\item \textbf{Ajouter une constante} $C$ : si $F$ est une primitive de $f$, alors toutes les primitives sont $F + C$ ($C \in \mathbb{R}$).
\item \textbf{Vérifier} en dérivant : on doit retrouver $F'(x) = f(x)$.
\end{enumerate}
\end{methode}

\begin{aretenir}[Formules de primitives à retenir (Polynômes)]
Primitive de $a$ : $ax + C$ ; primitive de $x$ : $\dfrac{x^2}{2} + C$ ; primitive de $x^2$ : $\dfrac{x^3}{3} + C$ ; primitive de $x^n$ : $\dfrac{x^{n+1}}{n+1} + C$.
\end{aretenir}

\begin{exoresolu}[Primitives d'un polynôme]
Soit $f$ la fonction définie sur $\mathbb{R}$ par $f(x) = 6x^2 - 4x + 5$.
\begin{enumerate}
\item Déterminer l'ensemble des primitives de $f$ sur $\mathbb{R}$.
\item Déterminer la primitive $F$ de $f$ qui s'annule en $1$.
\end{enumerate}
\tcblower
\textbf{Solution détaillée.}

\textbf{1.} On primitive terme à terme avec la formule : une primitive de $x^n$ est $\dfrac{x^{n+1}}{n+1}$.
\begin{align*}
F(x) &= 6 \times \dfrac{x^3}{3} - 4 \times \dfrac{x^2}{2} + 5 \times x + C \\
&= 2x^3 - 2x^2 + 5x + C, \quad C \in \mathbb{R}.
\end{align*}
\textbf{Vérification :} $F'(x) = 6x^2 - 4x + 5 = f(x)$. La réponse est correcte.

\textbf{2.} On cherche $C$ tel que $F(1) = 0$ :
\begin{align*}
F(1) &= 2 \times 1^3 - 2 \times 1^2 + 5 \times 1 + C = 0 \\
&\iff 2 - 2 + 5 + C = 0 \\
&\iff C = -5.
\end{align*}
\textbf{Conclusion :} la primitive de $f$ qui s'annule en $1$ est $F(x) = 2x^3 - 2x^2 + 5x - 5$.
\end{exoresolu}

\courssub{Applications concrètes : modélisation et interprétation}

\begin{methode}[Modéliser une situation technique ou économique]
Face à un problème concret (coût, bénéfice, vitesse, distance, remplissage) :
\begin{enumerate}
\item \textbf{Lire attentivement} l'énoncé et identifier la grandeur étudiée et sa variable (quantité, temps, prix).
\item \textbf{Traduire} la question en langage mathématique :
      \begin{itemize}
      \item « taux de variation instantané », « coût marginal », « vitesse », « débit » $\to$ \textbf{dérivée}.
      \item « retrouver la grandeur totale à partir de son taux », « quantité à partir du débit » $\to$ \textbf{primitive}.
      \end{itemize}
\item \textbf{Calculer} avec les formules usuelles.
\item \textbf{Interpréter} le résultat dans le contexte, avec l'\cle{unité} adaptée (FCFA, m, L, unités par heure...), et rédiger une phrase de conclusion.
\end{enumerate}
\end{methode}

\begin{exoresolu}[Coût marginal dans un atelier de Douala]
Un atelier de confection à Douala produit des chemises. Le coût total de production, en milliers de FCFA, de $x$ dizaines de chemises est modélisé par :
\[ C(x) = x^3 - 6x^2 + 15x + 10, \quad x \in [0\,;\,8]. \]
On appelle \emph{coût marginal} la dérivée $C'(x)$ : il approche le coût de production d'une dizaine de chemises supplémentaire.
\begin{enumerate}
\item Calculer le coût marginal $C'(x)$.
\item Calculer $C'(4)$ et interpréter le résultat.
\item Le coût marginal est-il minimal pour une production de 20 chemises ? Justifier.
\end{enumerate}
\tcblower
\textbf{Solution détaillée.}

\textbf{1.} $C$ est un polynôme, dérivable sur $[0\,;\,8]$ :
\[ C'(x) = 3x^2 - 12x + 15. \]

\textbf{2.} $C'(4) = 3 \times 16 - 12 \times 4 + 15 = 48 - 48 + 15 = 15$.

\textbf{Interprétation :} lorsque l'atelier produit 40 chemises ($x = 4$ dizaines), produire une dizaine de chemises supplémentaire coûte environ \num{15000} FCFA (15 milliers de FCFA).

\textbf{3.} Une production de 20 chemises correspond à $x = 2$. Étudions le minimum de $C'$ sur $[0\,;\,8]$. $C'$ est un polynôme du second degré ; sa dérivée est $(C')'(x) = 6x - 12$, qui s'annule en $x = 2$, en étant négative avant 2 et positive après. Donc $C'$ admet un minimum en $x = 2$, qui vaut :
\[ C'(2) = 3 \times 4 - 12 \times 2 + 15 = 12 - 24 + 15 = 3. \]
\textbf{Conclusion :} oui, le coût marginal est minimal pour une production de 20 chemises ; ce coût marginal minimal vaut \num{3000} FCFA par dizaine de chemises supplémentaire.
\end{exoresolu}

\begin{methode}[Problème inverse : de la dérivée à la fonction (Primitive avec condition initiale)]
Lorsqu'on connaît le taux de variation d'une grandeur et une valeur de cette grandeur :
\begin{enumerate}
\item \textbf{Identifier} la fonction dérivée connue $f = F'$ (vitesse, taux de production, recette marginale, débit...).
\item \textbf{Primitiver} $f$ pour obtenir $F(x) + C$.
\item \textbf{Utiliser la condition initiale} (valeur connue de $F$ en un point) pour déterminer la constante $C$.
\item \textbf{Conclure} en donnant l'expression complète de $F$ et répondre à la question posée avec les unités.
\end{enumerate}
\end{methode}

\begin{exoresolu}[Remplissage d'un réservoir d'eau]
Un réservoir d'un chantier à Yaoundé est rempli par une pompe. Le débit (en litres par minute) à l'instant $t$ (en minutes) est modélisé par :
\[ d(t) = 6t + 4, \quad t \in [0\,;\,10]. \]
Au bout de 2 minutes, le réservoir contient déjà 28 litres. On note $V(t)$ le volume d'eau (en litres) dans le réservoir à l'instant $t$, de sorte que $V'(t) = d(t)$.
\begin{enumerate}
\item Déterminer l'expression de $V(t)$.
\item Quel volume contient le réservoir au bout de 10 minutes ?
\end{enumerate}
\tcblower
\textbf{Solution détaillée.}

\textbf{1.} $V$ est une primitive de $d$. On primitive terme à terme :
\[ V(t) = 6 \times \dfrac{t^2}{2} + 4t + C = 3t^2 + 4t + C, \quad C \in \mathbb{R}. \]
On utilise la condition $V(2) = 28$ :
\begin{align*}
V(2) = 28 &\iff 3 \times 4 + 4 \times 2 + C = 28 \\
&\iff 12 + 8 + C = 28 \\
&\iff C = 8.
\end{align*}
Donc $V(t) = 3t^2 + 4t + 8$.

\textbf{Vérification :} $V'(t) = 6t + 4 = d(t)$ et $V(2) = 12 + 8 + 8 = 28$. Correct.

\textbf{2.} $V(10) = 3 \times 100 + 4 \times 10 + 8 = 300 + 40 + 8 = 348$.

\textbf{Conclusion :} au bout de 10 minutes, le réservoir contient 348 litres d'eau.
\end{exoresolu}

\courssub{Erreurs fréquentes et conseils pour le Probatoire}

\begin{attention}[Les 5 erreurs qui coûtent des points au Probatoire]
\begin{enumerate}
\item \textbf{Oublier que la dérivée d'une constante est nulle} : écrire $(7)' = 7$ ou $(-5)' = -5$ est faux ; toute constante a pour dérivée $0$.
\item \textbf{Dériver un produit terme à terme} : $(uv)' \neq u'v'$. Il faut obligatoirement écrire $(uv)' = u'v + uv'$. De même pour le quotient : $\left(\dfrac{u}{v}\right)' = \dfrac{u'v - uv'}{v^2}$, en n'oubliant pas le carré au dénominateur.
\item \textbf{Confondre dérivation et primitivation} : pour une primitive de $x^n$, on \emph{augmente} l'exposant et on \emph{divise} par $n+1$ : $\dfrac{x^{n+1}}{n+1}$. Écrire $n x^{n-1}$ pour une primitive est une inversion grave. Réflexe de sécurité : \textbf{toujours vérifier en redérivant}.
\item \textbf{Oublier la constante d'intégration $C$} : l'ensemble des primitives s'écrit $F(x) + C$. Si une condition est donnée (par exemple $F(1) = 0$), elle sert à calculer $C$ ; sans cette étape, la réponse est incomplète.
\item \textbf{Ne pas rédiger la conclusion ni les justifications} : au Probatoire, il faut annoncer la formule utilisée, justifier la dérivabilité (« $f$ est un polynôme, donc dérivable sur $\mathbb{R}$ »), et terminer par une phrase de conclusion avec l'unité dans les problèmes concrets. Une réponse sans phrase de conclusion perd des points de présentation.
\end{enumerate}
\end{attention}

\newpage

\coursec{Exercices}

\coursec{Fonctions : Dérivée et primitives}

\courssub{Rappels sur la dérivée et interprétation géométrique}

\begin{definition}[Nombre dérivé]
Soit $f$ une fonction définie sur un intervalle $I$ et $a \in I$. Si le taux d'accroissement $\dfrac{f(x)-f(a)}{x-a}$ admet une limite finie lorsque $x \to a$, on dit que $f$ est \textbf{dérivable} en $a$. Cette limite est appelée \textbf{nombre dérivé} de $f$ en $a$ et se note $f'(a)$.
\[ f'(a) = \lim_{x \to a} \frac{f(x)-f(a)}{x-a} = \lim_{h \to 0} \frac{f(a+h)-f(a)}{h}. \]
\end{definition}

\begin{propriete}[Interprétation géométrique]
Si $f$ est dérivable en $a$, la droite tangente $\mathcal{T}$ à la courbe $\mathcal{C}_f$ au point $A(a\,;\,f(a))$ existe et son coefficient directeur est $f'(a)$. Son équation réduite est :
\[ \mathcal{T} : y = f'(a)(x-a) + f(a). \]
\end{propriete}

\begin{exemplebox}[Tangente à une courbe]
Soit $f(x)=x^2$ et $A(1\,;\,1)$. $f'(x)=2x$, donc $f'(1)=2$.
L'équation de la tangente en $A$ est $y = 2(x-1)+1 = 2x-1$.
\end{exemplebox}

\courssub{Règles de dérivation usuelles}

\begin{propriete}[Linéarité de la dérivation]
Soient $u$ et $v$ deux fonctions dérivables sur un intervalle $I$ et $\lambda, \mu \in \mathbb{R}$.
La fonction $\lambda u + \mu v$ est dérivable sur $I$ et :
\[ (\lambda u + \mu v)' = \lambda u' + \mu v'. \]
\end{propriete}

\begin{propriete}[Formules de dérivation usuelles (à connaître)]
Soient $u$ et $v$ dérivables sur $I$.
\begin{itemize}
    \item \textbf{Produit :} $(uv)' = u'v + uv'$.
    \item \textbf{Quotient :} $\left(\dfrac{u}{v}\right)' = \dfrac{u'v - uv'}{v^2}$ \quad ($v(x) \neq 0$ sur $I$).
    \item \textbf{Composée (chaîne) :} Si $f$ est dérivable sur $J$ et $u(I) \subset J$, alors $f \circ u$ est dérivable sur $I$ et $(f \circ u)' = (f' \circ u) \times u'$.
\end{itemize}
\end{propriete}

\begin{methode}[Dériver une fonction]
\begin{enumerate}
    \item Identifier la forme de la fonction (somme, produit, quotient, composée, puissance).
    \item Appliquer la règle correspondante.
    \item Simplifier l'expression obtenue (factoriser si possible pour l'étude de signe).
    \item Préciser l'ensemble de définition de la dérivée (même que $f$ sauf pour $\sqrt{x}$, $\ln x$, quotients).
\end{enumerate}
\end{methode}

\courssub{Tableau des dérivées usuelles et premières primitives}

\begin{aretenir}[Tableau de référence (Programme Terminale Technique)]
\begin{center}
\begin{tabularx}{\linewidth}{|c|>{\centering\arraybackslash}X|>{\centering\arraybackslash}X|>{\centering\arraybackslash}X|}\hline
\rowcolor{popblueL}
\textbf{Fonction $f(x)$} & \textbf{Domaine $D_f$} & \textbf{Dérivée $f'(x)$} & \textbf{Primitive $F(x)$ (sur $D_f$)} \\ \hline
$k$ (constante) & $\mathbb{R}$ & $0$ & $kx + C$ \\ \hline
$x^n$ ($n \in \mathbb{N}^*$) & $\mathbb{R}$ & $n x^{n-1}$ & $\dfrac{x^{n+1}}{n+1} + C$ \\ \hline
$x^\alpha$ ($\alpha \in \mathbb{Q}$, $\alpha \neq -1$) & $]0,+\infty[$ \text{ ou } $\mathbb{R}^*$ & $\alpha x^{\alpha-1}$ & $\dfrac{x^{\alpha+1}}{\alpha+1} + C$ \\ \hline
$\sqrt{x} = x^{1/2}$ & $[0,+\infty[$ & $\dfrac{1}{2\sqrt{x}}$ & $\dfrac{2}{3}x^{3/2} + C$ \\ \hline
$\dfrac{1}{x} = x^{-1}$ & $\mathbb{R}^*$ & $-\dfrac{1}{x^2}$ & \textbf{Hors programme} (Logarithme) \\ \hline
\end{tabularx}
\end{center}
\textbf{Attention :} La primitive de $x \mapsto \frac{1}{x}$ est $x \mapsto \ln|x|$, non au programme de Terminale Technique. On ne demande que les primitives de fonctions \textbf{polynômes} ou \textbf{puissances $x^\alpha$ ($\alpha \in \mathbb{Q}, \alpha \neq -1$)}.
\end{aretenir}

\courssub{Primitives des fonctions polynomiales}

\begin{definition}[Primitive]
Soit $f$ une fonction définie sur un intervalle $I$. On appelle \textbf{primitive de $f$ sur $I$} toute fonction $F$ dérivable sur $I$ telle que $F' = f$ sur $I$.
\end{definition}

\begin{propriete}[Ensemble des primitives]
Si $F$ est une primitive de $f$ sur $I$, alors l'ensemble des primitives de $f$ sur $I$ est $\{F + C \mid C \in \mathbb{R}\}$. On note $\int f(x)\,dx = F(x) + C$.
\end{propriete}

\begin{propriete}[Linéarité de la recherche de primitives]
Soient $f, g$ admettant des primitives sur $I$ et $\lambda, \mu \in \mathbb{R}$.
\[ \int (\lambda f(x) + \mu g(x))\,dx = \lambda \int f(x)\,dx + \mu \int g(x)\,dx. \]
\end{propriete}

\begin{methode}[Déterminer les primitives d'un polynôme / fonction puissance]
\begin{enumerate}
    \item Écrire la fonction comme somme de termes de la forme $a x^\alpha$ ($\alpha \in \mathbb{Q}, \alpha \neq -1$).
    \item Appliquer la formule $\int x^\alpha\,dx = \dfrac{x^{\alpha+1}}{\alpha+1} + C$ terme à terme.
    \item Ajouter la constante $C$ \textbf{une seule fois} à la fin.
    \item Vérifier en dérivant le résultat.
\end{enumerate}
\end{methode}

\begin{exoresolu}[Primitives d'une fonction rationnelle/racine]
Déterminer les primitives sur $]0,+\infty[$ de $f(x) = \dfrac{3}{x^2} - 2\sqrt{x} + 5$.
\tcblower
$f(x) = 3x^{-2} - 2x^{1/2} + 5$.
$\int f(x)\,dx = 3\dfrac{x^{-1}}{-1} - 2\dfrac{x^{3/2}}{3/2} + 5x + C = -\dfrac{3}{x} - \dfrac{4}{3}x^{3/2} + 5x + C$.
Vérification : dérivée de $-\frac{3}{x}$ est $\frac{3}{x^2}$ ; dérivée de $-\frac{4}{3}x^{3/2}$ est $-2x^{1/2}$ ; dérivée de $5x$ est $5$. Correct.

\end{exoresolu}

\courssub{Résolution de problèmes concrets : modélisation et interprétation}

\begin{methode}[Modélisation par une fonction coût / bénéfice]
\begin{enumerate}
    \item \textbf{Identifier les variables :} $x$ = quantité produite/vendue ; $C(x)$ = coût total ; $R(x)$ = revenu total ; $B(x)=R(x)-C(x)$ = bénéfice.
    \item \textbf{Coût marginal :} $C_m(x) = C'(x)$. Coût approximatif de la $(x+1)$-ième unité.
    \item \textbf{Coût moyen :} $C_{moy}(x) = \dfrac{C(x)}{x}$ ($x>0$). Étudier ses variations pour trouver l'optimum de production.
    \item \textbf{Optimisation :} Étudier le signe de la dérivée ($B'(x)$, $C_m'(x)$, $C_{moy}'(x)$) $\to$ Tableau de variations $\to$ Maximum/Minimum.
    \item \textbf{Interpréter :} Revenir au contexte (unités, rentabilité, seuils).
\end{enumerate}
\end{methode}

\begin{exemplebox}[Coût marginal et coût moyen]
$C(x) = 0,5x^2 + 20x + 500$ (en kFCFA).
$C_m(x) = x + 20$. $C_{moy}(x) = 0,5x + 20 + \frac{500}{x}$.
$C_{moy}'(x) = 0,5 - \frac{500}{x^2} = 0 \iff x^2 = 1000 \iff x = 10\sqrt{10} \approx 31,6$.
Le coût moyen est minimal pour une production d'environ 31,6 tonnes.
\end{exemplebox}

\courssub{Je vérifie mes connaissances}

\begin{exercice}[QCM : Dérivées usuelles]
Parmi les quatre propositions suivantes, une seule est exacte. La recopier en justifiant brièvement.
\begin{enumerate}
    \item La dérivée de la fonction $f(x) = \sqrt{x}$ sur $]0,+\infty[$ est $f'(x) = \dfrac{1}{2\sqrt{x}}$.
    \item La dérivée de la fonction $g(x) = \dfrac{1}{x^2}$ sur $\mathbb{R}^*$ est $g'(x) = \dfrac{2}{x^3}$.
    \item La dérivée de la fonction $h(x) = x^{3/2}$ sur $[0,+\infty[$ est $h'(x) = \dfrac{2}{3}x^{1/2}$.
    \item La dérivée de la fonction $k(x) = 5x^4 - 3x^{-1}$ sur $\mathbb{R}^*$ est $k'(x) = 20x^3 - 3x^{-2}$.
\end{enumerate}
\end{exercice}

\begin{exercice}[Vrai ou Faux ? Justifier]
Dire si les affirmations suivantes sont vraies ou fausses. Justifier chaque réponse par un calcul ou un contre-exemple.
\begin{enumerate}
    \item Si $F$ est une primitive de $f$ sur $\mathbb{R}$, alors $F+5$ est aussi une primitive de $f$ sur $\mathbb{R}$.
    \item La fonction $f(x) = |x|$ est dérivable sur $\mathbb{R}$.
    \item Une primitive de la fonction $x \mapsto 3x^2 - 4x + 2$ sur $\mathbb{R}$ est $x \mapsto x^3 - 2x^2 + 2x$.
    \item Si $f'(x) > 0$ sur un intervalle $I$, alors la fonction $f$ est strictement croissante sur $I$.
\end{enumerate}
\end{exercice}

\begin{exercice}[Définitions et notations]
Compléter les phrases suivantes en utilisant le vocabulaire précis du cours.
\begin{enumerate}
    \item Soit $f$ une fonction dérivable sur un intervalle $I$. Le nombre $f'(a)$ représente géométriquement \emph{le coefficient directeur de la tangente à la courbe $\mathcal{C}_f$ au point d'abscisse $a$}.
    \item On appelle \emph{primitive} de la fonction $f$ sur $I$ toute fonction $F$ dérivable sur $I$ telle que $F' = f$ sur $I$.
    \item L'ensemble des primitives de $f$ sur $I$ se note \emph{$\int f(x)\,dx$} et s'écrit sous la forme \emph{$F(x) + C$} où $C$ est une constante réelle.
    \item La dérivée de la fonction constante $x \mapsto k$ ($k \in \mathbb{R}$) est \emph{la fonction nulle $x \mapsto 0$}.
\end{enumerate}
\end{exercice}

\begin{exercice}[Calculs directs de dérivées et primitives]
Sans rédiger de démonstration, donner le résultat demandé.
\begin{enumerate}
    \item Dériver $f(x) = 7x^3 - 4x^2 + 8x - 12$ sur $\mathbb{R}$.
    \item Dériver $g(x) = \dfrac{5}{x^3} - 2\sqrt{x} + 4$ sur $]0,+\infty[$.
    \item Trouver une primitive sur $\mathbb{R}$ de $h(x) = 12x^2 - 10x + 3$.
    \item Trouver une primitive sur $]0,+\infty[$ de $k(x) = \dfrac{4}{x^2} + 3\sqrt{x}$.
\end{enumerate}
\end{exercice}

\courssub{Je m'entraîne}

\begin{exercice}[Application des règles de dérivation]
Soit la fonction $f$ définie sur $\mathbb{R}$ par $f(x) = (2x^2 - 3)(x^3 + 4x)$.
\begin{enumerate}
    \item Développer $f(x)$ sous forme polynomiale.
    \item Calculer $f'(x)$ à partir de l'expression développée.
    \item Vérifier le résultat en utilisant la formule de dérivation d'un produit $(uv)' = u'v + uv'$.
    \item Déterminer l'équation de la tangente à la courbe $\mathcal{C}_f$ au point d'abscisse $1$.
\end{enumerate}
\end{exercice}

\begin{exercice}[Recherche de primitives : fonctions puissances]
On considère les fonctions suivantes (domaines adaptés pour éviter $1/x$) :
\begin{align*}
    f(x) &= \frac{6}{x^4} - \frac{2}{x^2} + 5 \quad \text{sur } ]0,+\infty[ \\
    g(x) &= 3\sqrt{x} - \frac{4}{\sqrt{x}} + 2x \quad \text{sur } ]0,+\infty[ \\
    h(x) &= (x+2)(2x-1) \quad \text{sur } \mathbb{R}
\end{align*}
Pour chacune d'elles, déterminer l'ensemble de ses primitives sur son ensemble de définition.
\end{exercice}

\begin{exercice}[Vérification et détermination de constante]
Soit la fonction $f$ définie sur $\mathbb{R}$ par $f(x) = 6x^2 - 8x + 5$ et la fonction $F$ définie sur $\mathbb{R}$ par $F(x) = 2x^3 - 4x^2 + 5x - 7$.
\begin{enumerate}
    \item Vérifier que $F$ est une primitive de $f$ sur $\mathbb{R}$.
    \item En déduire l'ensemble des primitives de $f$ sur $\mathbb{R}$.
    \item Déterminer l'unique primitive $G$ de $f$ sur $\mathbb{R}$ telle que $G(1) = 10$.
\end{enumerate}
\end{exercice}

\begin{exercice}[Modélisation simple : Coût marginal]
Dans une usine de transformation de bois, le coût total de production (en milliers de FCFA) pour $x$ tonnes de bois traitées est modélisé par la fonction $C$ définie sur $[0, +\infty[$ par $C(x) = 0,5x^2 + 20x + 500$.
\begin{enumerate}
    \item Interpréter concrètement le nombre $C(0)$.
    \item Déterminer la fonction coût marginal $C_m$. Que représente $C_m(x)$ pour l'industriel ?
    \item Calculer $C_m(10)$ et interpréter le résultat.
    \item Déterminer la fonction coût moyen $C_{moy}$. Étudier ses variations sur $]0, +\infty[$.
\end{enumerate}
\end{exercice}

\courssub{Je me prépare au Probatoire}

\begin{exercice}[Étude complète d'une fonction polynôme : Bénéfice d'une entreprise]
Une petite entreprise de fabrication de meubles a modélisé son bénéfice mensuel (en millions de FCFA) en fonction du nombre $x$ de centaines de meubles produits et vendus par la fonction $B$ définie sur $[0, +\infty[$ par :
\[ B(x) = -x^3 + 6x^2 + 15x - 20. \]
\begin{enumerate}
    \item Calculer la dérivée $B'(x)$ et factoriser l'expression obtenue.
    \item Résoudre l'inéquation $B'(x) > 0$ sur $[0, +\infty[$.
    \item Dresser le tableau de variations de $B$ sur $[0, +\infty[$.
    \item Quelle production (en centaines de meubles) permet de maximiser le bénéfice ? Calculer ce bénéfice maximal.
    \item L'entreprise ne peut pas produire plus de 800 meubles par mois. Quel est le bénéfice maximal réalisable dans cette contrainte ?
    \item Déterminer l'ensemble des primitives de la fonction $B$ sur $\mathbb{R}$. Parmi elles, laquelle s'annule en $x=0$ ?
\end{enumerate}
\end{exercice}

\begin{exercice}[Coût total, coût marginal et optimisation]
Le coût total $C_T(x)$ (en milliers de FCFA) pour la production de $x$ centaines d'unités d'un produit est donné par la fonction $C_T$ définie sur $[0, +\infty[$ par :
\[ C_T(x) = 2x^3 - 15x^2 + 60x + 125. \]
\begin{enumerate}
    \item Déterminer l'expression de la fonction coût marginal $C_m(x)$.
    \item Étudier les variations de $C_m$ sur $[0, +\infty[$. Pour quelle production le coût marginal est-il minimal ? Calculer ce coût marginal minimal.
    \item Déterminer l'expression de la fonction coût moyen $C_{moy}(x)$.
    \item Résoudre l'équation $C_m(x) = C_{moy}(x)$ sur $]0, +\infty[$. Interpréter graphiquement et économiquement le(s) résultat(s) obtenu(s).
    \item Trouver l'ensemble des primitives de la fonction $C_m$ sur $\mathbb{R}$. Vérifier que $C_T$ en fait partie.
\end{enumerate}
\end{exercice}

\begin{exercice}[Primitive, constante d'intégration et condition initiale]
Soit la fonction $f$ définie sur $\mathbb{R}$ par $f(x) = 4x^3 - 6x^2 + 10x - 4$.
\begin{enumerate}
    \item Déterminer l'ensemble des primitives de $f$ sur $\mathbb{R}$.
    \item Soit $F$ la primitive de $f$ sur $\mathbb{R}$ telle que $F(1) = 3$. Déterminer l'expression de $F(x)$.
    \item On considère la fonction $G$ définie sur $\mathbb{R}$ par $G(x) = x^4 - 2x^3 + 5x^2 - 4x + k$ où $k$ est un réel.
    \begin{enumerate}
        \item Pour quelle(s) valeur(s) de $k$ la fonction $G$ est-elle une primitive de $f$ ?
        \item Pour $k = 2$, calculer $G(2) - G(0)$. Interpréter ce nombre en terme d'aire (sans calcul d'intégrale formelle, en utilisant la relation entre $G$ et $f$).
    \end{enumerate}
    \item Soit la courbe $\mathcal{C}$ d'équation $y = f(x)$. Déterminer l'abscisse du point de $\mathcal{C}$ où la tangente est parallèle à la droite d'équation $y = 10x - 5$.
\end{enumerate}
\end{exercice}

\newpage

\coursec{Corrigés des exercices}

\begin{corrige}[Exercice 1 — QCM : Dérivées usuelles]
\textbf{Rappel des propositions :}
\begin{enumerate}
    \item $f(x)=\sqrt{x} \implies f'(x)=\dfrac{1}{2\sqrt{x}}$
    \item $g(x)=\dfrac{1}{x^2} \implies g'(x)=\dfrac{2}{x^3}$
    \item $h(x)=x^{3/2} \implies h'(x)=\dfrac{2}{3}x^{1/2}$
    \item $k(x)=5x^4-\dfrac{3}{x} \implies k'(x)=20x^3-\dfrac{3}{x^2}$
\end{enumerate}
Seule la \textbf{proposition 1} est exacte.

\begin{enumerate}
    \item \textbf{Exacte.} $f(x)=\sqrt{x}=x^{1/2}$. D'après la formule $(x^\alpha)'=\alpha x^{\alpha-1}$ avec $\alpha=\dfrac12$ :
    \[ f'(x)=\frac12 x^{-1/2}=\frac{1}{2\sqrt{x}}. \]
    \item \textbf{Fausse.} $g(x)=\dfrac{1}{x^2}=x^{-2}$, donc $g'(x)=-2x^{-3}=-\dfrac{2}{x^3}$. Le signe proposé ($+$) est faux.
    \item \textbf{Fausse.} $h(x)=x^{3/2}$, donc $h'(x)=\dfrac32 x^{1/2}$. Le coefficient $\dfrac23$ est celui de la \emph{primitive} ($\int x^{3/2} = \frac{2}{5}x^{5/2}$), pas de la dérivée : confusion classique entre dérivation (exposant descend) et primitivation (exposant monte).
    \item \textbf{Fausse.} $k(x)=5x^4-3x^{-1}$, donc $k'(x)=20x^3-3(-1)x^{-2}=20x^3+3x^{-2}=20x^3+\dfrac{3}{x^2}$. Le signe du second terme est faux dans la proposition.
\end{enumerate}

\textbf{Point de vigilance :} Pour $x^\alpha$, la dérivée fait \emph{descendre} l'exposant ($\alpha x^{\alpha-1}$), la primitive le fait \emph{monter} ($\dfrac{x^{\alpha+1}}{\alpha+1}$, $\alpha\neq -1$).
\end{corrige}

\begin{attention}[Pièges classiques sur les puissances]
\begin{itemize}
    \item Oublier le signe moins quand l'exposant est négatif : $(x^{-n})' = -n x^{-n-1}$.
    \item Confondre le coefficient de la dérivée ($\alpha$) et celui de la primitive ($1/(\alpha+1)$).
    \item Ne pas réécrire les racines et fractions sous forme de puissances ($x^\alpha$) avant de dériver/primitiver.
\end{itemize}
\end{attention}

\begin{corrige}[Exercice 2 — Vrai ou Faux ? Justifier]
\begin{enumerate}
    \item \textbf{Vrai.} Soit $F$ une primitive de $f$ sur $I$. Pour toute constante $C\in\mathbb{R}$, $(F+C)' = F' + 0 = f$. Toute fonction $F+C$ est donc une primitive de $f$ sur $I$. L'ensemble des primitives est bien $F(x)+C,\ C\in\mathbb{R}$.
    \item \textbf{Faux.} La fonction $f(x)=|x|$ n'est pas dérivable en $0$. Le taux d'accroissement en $0$ est $\dfrac{|h|}{h}$ : il tend vers $1$ quand $h\to 0^+$ et vers $-1$ quand $h\to 0^-$. Les limites à droite et à gauche étant différentes, la limite n'existe pas : $f$ n'est pas dérivable en $0$ (point anguleux / cuspide).
    \item \textbf{Vrai.} Soit $F(x)=x^3-2x^2+2x$. Alors $F'(x)=3x^2-4x+2=f(x)$ pour tout $x\in\mathbb{R}$. $F$ est bien une primitive de $f$ sur $\mathbb{R}$.
    \item \textbf{Vrai.} C'est le théorème du sens de variation (ou théorème des accroissements finis) : si $f$ est dérivable sur un intervalle $I$ et $f'(x)>0$ pour tout $x\in I$, alors $f$ est strictement croissante sur $I$.
\end{enumerate}
\end{corrige}

\begin{corrige}[Exercice 3 — Définitions et notations]
\begin{enumerate}
    \item Le nombre $f'(a)$ représente géométriquement le \cle{coefficient directeur} de la \cle{tangente} à la courbe $\mathcal{C}_f$ au point d'abscisse $a$. L'équation de cette tangente est $y = f'(a)(x-a) + f(a)$.
    \item On appelle \cle{primitive} de $f$ sur un intervalle $I$ toute fonction $F$ dérivable sur $I$ telle que $F'=f$ sur $I$.
    \item L'ensemble des primitives de $f$ sur $I$ se note $\displaystyle\int f(x)\,dx$ et s'écrit $F(x)+C$, avec $C\in\mathbb{R}$ (constante d'intégration).
    \item La dérivée de la fonction constante $x\mapsto k$ ($k\in\mathbb{R}$) est la fonction nulle $x\mapsto 0$.
\end{enumerate}
\end{corrige}

\begin{corrige}[Exercice 4 — Calculs directs de dérivées et primitives]
\begin{enumerate}
    \item $f(x)=7x^3-4x^2+8x \implies f'(x)=21x^2-8x+8$.
    \item $g(x)=5x^{-3}-2x^{1/2}+4$, donc
    \[ g'(x)=5(-3)x^{-4}-2\times\frac12 x^{-1/2}+0 = -15x^{-4}-x^{-1/2} = -\frac{15}{x^4}-\frac{1}{\sqrt{x}}. \]
    \item $h(x)=12x^2-10x+3$. Une primitive de $h$ sur $\mathbb{R}$ est :
    \[ H(x)=\frac{12}{3}x^3-\frac{10}{2}x^2+3x = 4x^3-5x^2+3x \quad (\text{on peut ajouter toute constante } C). \]
    \item $k(x)=4x^{-2}+3x^{1/2}$ sur $]0,+\infty[$. Une primitive est
    \[ K(x)=4\,\frac{x^{-1}}{-1}+3\,\frac{x^{3/2}}{3/2} = -\frac{4}{x}+2x^{3/2} + C,\quad C\in\mathbb{R}. \]
\end{enumerate}
\end{corrige}

\begin{corrige}[Exercice 5 — Application des règles de dérivation (produit)]
Soit $f(x)=(2x^2-3)(x^3+4x)$.
\begin{enumerate}
    \item On développe pour éviter la règle du produit :
    \begin{align*}
    f(x)&=(2x^2-3)(x^3+4x)\\
    &=2x^2(x^3+4x)-3(x^3+4x)\\
    &=2x^5+8x^3-3x^3-12x\\
    &=2x^5+5x^3-12x.
    \end{align*}
    \item On dérive le polynôme obtenu :
    \[ f'(x)=10x^4+15x^2-12. \]
    \item Vérification par la règle du produit $(uv)'=u'v+uv'$ :
    Posons $u(x)=2x^2-3$ et $v(x)=x^3+4x$ ; alors $u'(x)=4x$ et $v'(x)=3x^2+4$.
    \begin{align*}
    (uv)'&=u'v+uv'\\
    &=4x(x^3+4x)+(2x^2-3)(3x^2+4)\\
    &=4x^4+16x^2+6x^4+8x^2-9x^2-12\\
    &=10x^4+15x^2-12.
    \end{align*}
    On retrouve bien le résultat de la question 2.
    \item Équation de la tangente $\mathcal{T}$ en $x=1$ :
    $f(1)=2+5-12=-5$ et $f'(1)=10+15-12=13$.
    \[ \mathcal{T}:\ y=f'(1)(x-1)+f(1)=13(x-1)-5=13x-18. \]
\end{enumerate}
\end{corrige}

\begin{corrige}[Exercice 6 — Recherche de primitives : fonctions puissances]
\begin{enumerate}
    \item $f(x)=6x^{-4}-2x^{-2}+5$ sur $]0,+\infty[$ :
    \[ \int f(x)\,dx=6\,\frac{x^{-3}}{-3}-2\,\frac{x^{-1}}{-1}+5x+C = -\frac{2}{x^3}+\frac{2}{x}+5x+C,\quad C\in\mathbb{R}. \]
    \item $g(x)=3x^{1/2}-4x^{-1/2}+2x$ sur $]0,+\infty[$ :
    \[ \int g(x)\,dx=3\,\frac{x^{3/2}}{3/2}-4\,\frac{x^{1/2}}{1/2}+x^2+C = 2x^{3/2}-8\sqrt{x}+x^2+C,\quad C\in\mathbb{R}. \]
    \item $h(x)=(x+2)(2x-1)=2x^2+3x-2$ sur $\mathbb{R}$ :
    \[ \int h(x)\,dx=\frac{2}{3}x^3+\frac{3}{2}x^2-2x+C,\quad C\in\mathbb{R}. \]
\end{enumerate}
\textbf{Point de vigilance :} Penser à développer avant de primitiver un produit de polynômes, et à ajouter la constante $C$ une seule fois à la fin.
\end{corrige}

\begin{corrige}[Exercice 7 — Vérification et détermination de constante]
Soit $f(x)=6x^2-8x+5$ sur $\mathbb{R}$ et $F(x)=2x^3-4x^2+5x-7$.
\begin{enumerate}
    \item $F'(x)=6x^2-8x+5=f(x)$ pour tout $x\in\mathbb{R}$. Donc $F$ est bien une primitive de $f$ sur $\mathbb{R}$.
    \item L'ensemble des primitives de $f$ sur $\mathbb{R}$ est :
    \[ \left\{\,x\mapsto 2x^3-4x^2+5x-7+C\ \middle|\ C\in\mathbb{R}\,\right\}=\left\{\,x\mapsto 2x^3-4x^2+5x+C\ \middle|\ C\in\mathbb{R}\,\right\}. \]
    (La constante $-7$ est absorbée dans la constante arbitraire $C$).
    \item Soit $G(x)=2x^3-4x^2+5x+C$ la primitive cherchée. La condition $G(1)=10$ donne :
    \[ 2-4+5+C=10 \iff 3+C=10 \iff C=7. \]
    Donc $G(x)=2x^3-4x^2+5x+7$.
\end{enumerate}
\end{corrige}

\begin{corrige}[Exercice 8 — Modélisation simple : Coût marginal et coût moyen]
On considère la fonction coût total $C(x)=0{,}5x^2+20x+500$ (en milliers de FCFA) pour une production de $x$ tonnes ($x>0$).
\begin{enumerate}
    \item $C(0)=500$ : ce sont les \textbf{coûts fixes} de l'usine (500 milliers de FCFA), c'est-à-dire les dépenses engagées même sans aucune production (loyer, entretien des machines, salaires fixes).
    \item $C_m(x)=C'(x)=x+20$. Le coût marginal $C_m(x)$ représente, pour l'industriel, le coût approximatif de production d'une tonne supplémentaire lorsque le niveau de production est $x$ tonnes.
    \item $C_m(10)=10+20=30$. Produire une tonne supplémentaire à partir de 10 tonnes coûte environ 30 milliers de FCFA.
    \item Coût moyen : $C_{moy}(x)=\dfrac{C(x)}{x}=0{,}5x+20+\dfrac{500}{x}$ pour $x>0$.
    \[ C_{moy}'(x)=0{,}5-\frac{500}{x^2}=\frac{0{,}5x^2-500}{x^2}. \]
    Comme $x^2>0$, le signe de $C_{moy}'$ est celui de $0{,}5x^2-500$, qui s'annule pour $x^2=1000$, soit $x=\sqrt{1000}=10\sqrt{10}\approx 31{,}6$ (valeur positive).
    \begin{center}
    \begin{tabular}{|c|ccccc|}\hline
    $x$ & $0$ & & $10\sqrt{10}$ & & $+\infty$ \\ \hline
    $C_{moy}'(x)$ & & $-$ & $0$ & $+$ & \\ \hline
    $C_{moy}$ & $+\infty$ & $\searrow$ & $10\sqrt{10}+20\approx 51{,}6$ & $\nearrow$ & $+\infty$ \\ \hline
    \end{tabular}
    \end{center}
    Le coût moyen est minimal pour une production d'environ $31{,}6$ tonnes, et vaut $C_{moy}(10\sqrt{10})=10\sqrt{10}+20\approx 51{,}6$ milliers de FCFA par tonne.
\end{enumerate}
\end{corrige}

\begin{corrige}[Exercice 9 — Étude complète d'une fonction polynôme : Bénéfice d'une entreprise]
\textbf{Barème indicatif (sur 10) :} Q1 : 1,5 pt ; Q2 : 1,5 pt ; Q3 : 2 pts ; Q4 : 1,5 pt ; Q5 : 1,5 pt ; Q6 : 2 pts.

Soit $B(x)=-x^3+6x^2+15x-20$ le bénéfice (en millions de FCFA) pour une production de $x$ centaines de meubles ($x\ge 0$).
\begin{enumerate}
    \item $B'(x)=-3x^2+12x+15=-3(x^2-4x-5)$. Le trinôme $x^2-4x-5$ admet pour racines $x=-1$ et $x=5$ (discriminant $\Delta=16+20=36$). Donc :
    \[ B'(x)=-3(x+1)(x-5). \]
    \item Sur $[0,+\infty[$, $x+1>0$ et $-3<0$, donc le signe de $B'(x)$ est opposé à celui de $(x-5)$.
    \[ B'(x)>0 \iff x-5<0 \iff x<5. \]
    Ainsi $B'(x)>0$ sur $[0,5[$ et $B'(x)<0$ sur $]5,+\infty[$.
    \item $B(0)=-20$ ; $B(5)=-125+150+75-20=80$ ; $\lim_{x\to+\infty}B(x)=-\infty$ (terme dominant $-x^3$).
    \begin{center}
    \begin{tabular}{|c|ccccc|}\hline
    $x$ & $0$ & & $5$ & & $+\infty$ \\ \hline
    $B'(x)$ & & $+$ & $0$ & $-$ & \\ \hline
    $B$ & $-20$ & $\nearrow$ & $80$ & $\searrow$ & $-\infty$ \\ \hline
    \end{tabular}
    \end{center}
    \item Le bénéfice est maximal pour $x=5$, soit une production de \textbf{500 meubles} par mois. Le bénéfice maximal est $B(5)=80$ millions de FCFA.
    \item La contrainte de capacité est $x\leq 8$ (800 meubles). Sur l'intervalle $[0,8]$, le maximum de $B$ est atteint en $x=5$ car $5\in[0,8]$ et $B$ décroît après 5. Le bénéfice maximal réalisable reste donc $B(5)=80$ millions de FCFA (la contrainte n'est pas active / ne change pas l'optimum).
    \item $\displaystyle\int B(x)\,dx=\int(-x^3+6x^2+15x-20)\,dx=-\frac{x^4}{4}+2x^3+\frac{15}{2}x^2-20x+C$, $C\in\mathbb{R}$.
    La primitive s'annulant en $0$ vérifie $F(0)=C=0$ :
    \[ F(x)=-\frac{x^4}{4}+2x^3+\frac{15}{2}x^2-20x. \]
\end{enumerate}

\textbf{Points de vigilance :} Factoriser $B'$ avant d'étudier son signe ; justifier le signe par le tableau de signes du trinôme ; convertir $x$ (centaines de meubles) en nombre réel de meubles dans la conclusion ; ne pas oublier la constante $C$ dans les primitives.
\end{corrige}

\begin{corrige}[Exercice 10 — Coût total, coût marginal et optimisation]
\textbf{Barème indicatif (sur 10) :} Q1 : 1 pt ; Q2 : 3 pts ; Q3 : 1 pt ; Q4 : 3 pts ; Q5 : 2 pts.

Soit $C_T(x)=2x^3-15x^2+60x+125$ le coût total (en milliers de FCFA) pour une production de $x$ centaines d'unités ($x>0$).
\begin{enumerate}
    \item Coût marginal : $C_m(x)=C_T'(x)=6x^2-30x+60$.
    \item On étudie les variations de $C_m$ pour trouver son minimum.
    $C_m'(x)=12x-30=12\left(x-\dfrac52\right)$.
    $C_m'(x)<0$ pour $x<\dfrac52$ et $C_m'(x)>0$ pour $x>\dfrac52$.
    \begin{center}
    \begin{tabular}{|c|ccccc|}\hline
    $x$ & $0$ & & $\dfrac52$ & & $+\infty$ \\ \hline
    $C_m'(x)$ & & $-$ & $0$ & $+$ & \\ \hline
    $C_m$ & $60$ & $\searrow$ & $22{,}5$ & $\nearrow$ & $+\infty$ \\ \hline
    \end{tabular}
    \end{center}
    Le coût marginal est minimal pour $x=\dfrac52$, soit une production de \textbf{250 unités}, et ce minimum vaut :
    \[ C_m\!\left(\tfrac52\right)=6\times\tfrac{25}{4}-30\times\tfrac52+60=37{,}5-75+60=22{,}5 \text{ milliers de FCFA par centaine d'unités (soit 225 FCFA/unité)}. \]
    \item Coût moyen : $C_{moy}(x)=\dfrac{C_T(x)}{x}=2x^2-15x+60+\dfrac{125}{x}$ pour $x>0$.
    \item Résolvons $C_m(x)=C_{moy}(x)$ sur $]0,+\infty[$ :
    \begin{align*}
    6x^2-30x+60&=2x^2-15x+60+\frac{125}{x}\\
    4x^2-15x-\frac{125}{x}&=0\\
    4x^3-15x^2-125&=0 \quad (\text{en multipliant par } x>0).
    \end{align*}
    On teste $x=5$ : $4\times125-15\times25-125=500-375-125=0$. Donc $x=5$ est solution.
    Factorisation : $4x^3-15x^2-125=(x-5)(4x^2+5x+25)$. Le trinôme $4x^2+5x+25$ a pour discriminant $\Delta=25-400=-375<0$ : pas d'autre racine réelle.
    \[ \boxed{x=5} \]
    \textbf{Interprétation :} Graphiquement, les courbes de $C_m$ et $C_{moy}$ se coupent au point d'abscisse $5$. Économiquement, pour une production de 500 unités, le coût marginal égale le coût moyen : c'est le point où le coût moyen cesse de décroître et commence à croître (optimum technique de production / minimum du coût moyen).
    \item $\displaystyle\int C_m(x)\,dx=\int(6x^2-30x+60)\,dx=2x^3-15x^2+60x+C$, $C\in\mathbb{R}$.
    Or $C_T(x)=2x^3-15x^2+60x+125$ : c'est bien la primitive de $C_m$ obtenue pour $C=125$ (les coûts fixes). Vérification : $C_T'(x)=6x^2-30x+60=C_m(x)$.
\end{enumerate}

\textbf{Points de vigilance :} Bien dériver $C_m$ (et non $C_T$) à la question 2 ; multiplier par $x$ en précisant $x>0$ à la question 4 ; citer la valeur du discriminant pour exclure d'autres solutions réelles.
\end{corrige}

\begin{corrige}[Exercice 11 — Primitive, constante d'intégration et condition initiale]
\textbf{Barème indicatif (sur 10) :} Q1 : 2 pts ; Q2 : 2 pts ; Q3 : 4 pts ; Q4 : 2 pts.

Soit $f(x)=4x^3-6x^2+10x-4$ sur $\mathbb{R}$.
\begin{enumerate}
    \item $\displaystyle\int f(x)\,dx=\int(4x^3-6x^2+10x-4)\,dx=x^4-2x^3+5x^2-4x+C$, $C\in\mathbb{R}$.
    \item Soit $F(x)=x^4-2x^3+5x^2-4x+C$. La condition $F(1)=3$ donne :
    \[ 1-2+5-4+C=3 \iff 0+C=3 \iff C=3. \]
    Donc $F(x)=x^4-2x^3+5x^2-4x+3$.
    \item Soit $G(x)=x^4-2x^3+5x^2-4x+k$ avec $k\in\mathbb{R}$.
    \begin{enumerate}
        \item $G'(x)=4x^3-6x^2+10x-4=f(x)$ pour tout réel $x$, et ce \textbf{quel que soit} $k$ (la dérivée d'une constante est nulle). Donc $G$ est une primitive de $f$ pour \textbf{toute valeur de $k\in\mathbb{R}$}.
        \item Pour $k=2$ : $G(x)=x^4-2x^3+5x^2-4x+2$.
        $G(2)=16-16+20-8+2=14$ et $G(0)=2$, donc $G(2)-G(0)=12$.
        \textbf{Interprétation :} Comme $G'=f$, le nombre $G(2)-G(0)$ représente la \textbf{variation de la primitive $G$} entre $0$ et $2$, c'est-à-dire l'\textbf{accumulation} de la grandeur $f$ sur $[0,2]$. C'est l'aire algébrique entre la courbe de $f$ et l'axe des abscisses sur $[0,2]$, notée plus tard $\displaystyle\int_0^2 f(x)\,dx=12$. (On note que $f$ change de signe sur $[0,2]$ car $f(0)=-4<0$ et $f(1)=4>0$, l'intégrale est donc une somme algébrique d'aires).
    \end{enumerate}
    \item La tangente en $a$ est parallèle à la droite $y=10x-5$ si et seulement si leurs coefficients directeurs sont égaux, c'est-à-dire $f'(a)=10$.
    \[ f'(x)=12x^2-12x+10. \]
    \[ 12x^2-12x+10=10 \iff 12x^2-12x=0 \iff 12x(x-1)=0 \iff x=0 \text{ ou } x=1. \]
    Les points de $\mathcal{C}_f$ d'abscisses $x=0$ et $x=1$ ont une tangente parallèle à la droite donnée.
\end{enumerate}

\textbf{Points de vigilance :} À la question 3.a, bien justifier que $k$ disparaît par dérivation ; à la question 4, écrire la condition « coefficients directeurs égaux » avant de résoudre ; pour l'interprétation de $G(b)-G(a)$, parler d'aire algébrique (intégrale) et non d'aire géométrique si $f$ change de signe.
\end{corrige}

\newpage

\coursec{Fiche bilan}

\begin{popfigure}
\begin{tikzpicture}[mindmap, grow cyclic, every node/.style={concept}, root concept/.append style={concept color=popblue, font=\bfseries\small}, level 1/.append style={level distance=3.6cm, sibling angle=51.4}, level 2/.append style={level distance=2.2cm, sibling angle=45, font=\scriptsize}]
\node[root concept]{Dérivées et primitives}
  child[concept color=poporange]{ node {Dérivée : définition}
    child { node {$f'(a)=\lim\limits_{h\to0}\dfrac{f(a+h)-f(a)}{h}$} }
    child { node {Tangente : $y=f'(a)(x-a)+f(a)$} }
  }
  child[concept color=popgreen]{ node {Opérations}
    child { node {$(u+v)'=u'+v'$} }
    child { node {$(ku)'=ku'$} }
    child { node {$(uv)'=u'v+uv'$} }
    child { node {$\left(\dfrac{u}{v}\right)'=\dfrac{u'v-uv'}{v^{2}}$} }
  }
  child[concept color=poppurple]{ node {Dérivées usuelles}
    child { node {$(x^{n})'=nx^{n-1}$} }
    child { node {$(\sqrt{x})'=\dfrac{1}{2\sqrt{x}}$} }
    child { node {$(\cos)'=-\sin$ ; $(\sin)'=\cos$} }
  }
  child[concept color=poppink]{ node {Sens de variation}
    child { node {$f'\geq0 \Rightarrow f$ croissante} }
    child { node {$f'\leq0 \Rightarrow f$ décroissante} }
  }
  child[concept color=popgold]{ node {Primitives}
    child { node {$F'=f$} }
    child { node {$F$ et $F+k$} }
  }
  child[concept color=popgreen!70!black]{ node {Primitives de polynômes}
    child { node {$x^{n}\to\dfrac{x^{n+1}}{n+1}$} }
    child { node {Terme à terme} }
  }
  child[concept color=popblue!60]{ node {Problèmes concrets}
    child { node {Vitesse, coût marginal} }
    child { node {Optimisation} }
  };
\end{tikzpicture}
\legende{Carte mentale de la leçon : dérivées et primitives.}
\end{popfigure}

\begin{aretenir}[L'essentiel de la leçon]
\textbf{Définitions clés}
\begin{itemize}
  \item \cle{Nombre dérivé} de $f$ en $a$ : $f'(a)=\lim\limits_{h\to0}\dfrac{f(a+h)-f(a)}{h}$ (quand cette limite existe et est finie). On dit alors que $f$ est \cle{dérivable} en $a$.
  \item \cle{Fonction dérivée} $f'$ : fonction qui à tout $x$ associe $f'(x)$, là où $f$ est dérivable.
  \item \cle{Tangente} à la courbe de $f$ au point d'abscisse $a$ : $y=f'(a)(x-a)+f(a)$ ; $f'(a)$ est son coefficient directeur.
  \item \cle{Primitive} de $f$ sur un intervalle $I$ : toute fonction $F$ dérivable sur $I$ telle que $F'=f$ sur $I$. Si $F$ est une primitive, alors toutes les primitives sont les fonctions $F+k$ ($k$ constante réelle).
\end{itemize}

\textbf{Formules à connaître par c\oe ur}
\begin{center}
\begin{tabularx}{\linewidth}{|l|X|X|}\hline
\rowcolor{popblueL} \textbf{Fonction $f$} & \textbf{Dérivée $f'$} & \textbf{Une primitive $F$} \\ \hline
$k$ (constante) & $0$ & $kx$ \\ \hline
$x$ & $1$ & $\dfrac{x^{2}}{2}$ \\ \hline
$x^{n}$ ($n\in\mathbb{N}^{*}$) & $nx^{n-1}$ & $\dfrac{x^{n+1}}{n+1}$ \\ \hline
$\dfrac{1}{x}$ ($x\neq0$) & $-\dfrac{1}{x^{2}}$ & (hors polynômes) \\ \hline
$\sqrt{x}$ ($x>0$) & $\dfrac{1}{2\sqrt{x}}$ & $\dfrac{2}{3}x\sqrt{x}$ \\ \hline
$\cos x$ & $-\sin x$ & $\sin x$ \\ \hline
$\sin x$ & $\cos x$ & $-\cos x$ \\ \hline
\end{tabularx}
\end{center}

\textbf{Opérations sur les dérivées} : $(u+v)'=u'+v'$ ; $(ku)'=ku'$ ($k$ constante) ; $(uv)'=u'v+uv'$ ; $\left(\dfrac{u}{v}\right)'=\dfrac{u'v-uv'}{v^{2}}$ avec $v(x)\neq0$.

\textbf{Lien dérivée / variations} : si $f'\geq0$ sur un intervalle $I$, alors $f$ est croissante sur $I$ ; si $f'\leq0$ sur $I$, alors $f$ est décroissante sur $I$. Si $f$ admet un extremum local en $c$ et si $f'$ s'annule en changeant de signe en $c$, alors $f'(c)=0$.

\textbf{Méthodes express}
\begin{enumerate}
  \item \textbf{Dériver un polynôme} : dériver terme à terme avec $(ax^{n})'=anx^{n-1}$.
  \item \textbf{Primitiver un polynôme} : pour chaque terme $ax^{n}$, écrire $\dfrac{a\,x^{n+1}}{n+1}$, puis ajouter $+k$ ; \emph{vérifier en redérivant}.
  \item \textbf{Trouver la primitive vérifiant $F(x_{0})=y_{0}$} : écrire la primitive générale, puis résoudre l'équation $F(x_{0})=y_{0}$ pour déterminer $k$.
  \item \textbf{Équation de tangente} : calculer $f(a)$ et $f'(a)$, puis remplacer dans $y=f'(a)(x-a)+f(a)$.
  \item \textbf{Optimisation} : dériver, résoudre $f'(x)=0$, étudier le signe de $f'$, conclure avec un tableau de variations et une phrase en lien avec la situation.
\end{enumerate}
\end{aretenir}

\begin{attention}[Pièges fréquents au Probatoire]
\begin{itemize}
  \item Ne pas confondre dérivation et primitivation : $x^{n}$ a pour dérivée $nx^{n-1}$ (l'exposant \emph{diminue}) et pour primitive $\dfrac{x^{n+1}}{n+1}$ (l'exposant \emph{augmente}).
  \item La dérivée d'un produit n'est \textbf{pas} le produit des dérivées : $(uv)'=u'v+uv'$, et non $u'v'$.
  \item Au dénominateur du quotient, c'est $v^{2}$, et l'ordre du numérateur compte : $u'v-uv'$ (signe moins).
  \item Toute primitive générale s'écrit avec une constante $+k$ ; l'oublier fait perdre des points.
  \item Une primitive se \emph{vérifie} toujours en la redérivant : on doit retrouver $f$.
  \item Pour la tangente, ne pas oublier le terme $f(a)$ : $y=f'(a)(x-a)+f(a)$, et non $y=f'(a)(x-a)$.
  \item $f'(c)=0$ ne suffit pas à garantir un extremum : il faut vérifier le \emph{changement de signe} de $f'$ (tableau de variations).
\end{itemize}
\end{attention}

\begin{exoresolu}[Vérifier une primitive]
Énoncé. Soit $f$ la fonction définie sur $\mathbb{R}$ par $f(x)=3x^{2}-4x+5$. Déterminer la primitive $F$ de $f$ telle que $F(1)=2$.
\tcblower
Solution détaillée. On primitive terme à terme : $3x^{2}\to 3\times\dfrac{x^{3}}{3}=x^{3}$ ; $-4x\to -4\times\dfrac{x^{2}}{2}=-2x^{2}$ ; $5\to 5x$. Les primitives de $f$ sont donc
\[ F(x)=x^{3}-2x^{2}+5x+k,\quad k\in\mathbb{R}. \]
Vérification : $F'(x)=3x^{2}-4x+5=f(x)$. La condition $F(1)=2$ donne $1-2+5+k=2$, soit $4+k=2$, d'où $k=-2$. Conclusion :
\[ F(x)=x^{3}-2x^{2}+5x-2. \]
\end{exoresolu}

\begin{aretenir}[Checklist avant le Probatoire]
\begin{itemize}
  \item[$\square$] Je sais donner la définition du nombre dérivé $f'(a)$ comme limite du taux d'accroissement.
  \item[$\square$] Je connais par c\oe ur les dérivées de $x^{n}$, $\dfrac{1}{x}$, $\sqrt{x}$, $\cos x$, $\sin x$.
  \item[$\square$] Je sais dériver une somme, un produit par une constante, un produit et un quotient.
  \item[$\square$] Je sais écrire l'équation de la tangente en un point : $y=f'(a)(x-a)+f(a)$.
  \item[$\square$] Je sais utiliser le signe de $f'$ pour dresser le tableau de variations de $f$.
  \item[$\square$] Je sais qu'une primitive $F$ de $f$ vérifie $F'=f$ et que deux primitives diffèrent d'une constante.
  \item[$\square$] Je sais primitiver un polynôme terme à terme : $ax^{n}\to\dfrac{a\,x^{n+1}}{n+1}$.
  \item[$\square$] Je n'oublie jamais la constante $k$ dans une primitive générale, et je sais la déterminer avec une condition $F(x_{0})=y_{0}$.
  \item[$\square$] Je vérifie toujours une primitive en la redérivant.
  \item[$\square$] Je sais résoudre $f'(x)=0$ pour trouver les extremums et résoudre un problème d'optimisation (coût, bénéfice, aire).
  \item[$\square$] Je rédige proprement : justification par le signe de la dérivée, tableau de variations complet, conclusion en lien avec la situation concrète.
\end{itemize}
\end{aretenir}

\findecours

\end{document}
